% !TeX program = xelatex
% 数学 LaTeX 排版 Skill v4.0.0；版式保持 v3 金标准. XeLaTeX 编译至少两次. 不分发字体.
% WZ-LOCKED-CLASS-BEGIN
\begin{filecontents*}[overwrite]{elegantbook-original-adapter.cls}
\NeedsTeXFormat{LaTeX2e}
\ProvidesClass{elegantbook-original-adapter}[2026/09/23 v3.0 fixed mathematical book environments]
\LoadClass[10pt,letterpaper,oneside]{book}
\RequirePackage[UTF8,scheme=plain,fontset=fandol]{ctex}
\RequirePackage{fontspec}
\setmainfont{cmunrm.otf}[BoldFont=cmunbx.otf,ItalicFont=cmunti.otf,BoldItalicFont=cmunbi.otf]
\setCJKfamilyfont{sourceheader}[ItalicFont=FandolKai-Regular.otf]{FandolKai-Regular.otf}
\RequirePackage{amsmath,amssymb,mathrsfs}
\RequirePackage{amsthm,etoolbox}
\RequirePackage{xcolor,geometry,setspace,fancyhdr,enumitem,needspace,xurl}
\definecolor{structurecolor}{HTML}{880E4F}
\definecolor{main}{HTML}{9C27B0}
\geometry{paperwidth=612bp,paperheight=792bp,left=20mm,right=20mm,top=95.04bp,bottom=20mm,headheight=12pt,headsep=25pt,footskip=30pt}
\setstretch{1.6}
\setlength{\parindent}{2em}
\setlength{\parskip}{0pt}
\setlist{nosep,leftmargin=2em}
\AtBeginDocument{\setlength{\abovedisplayskip}{4pt plus 1pt minus 1pt}\setlength{\belowdisplayskip}{4pt plus 1pt minus 1pt}\setlength{\abovedisplayshortskip}{2pt}\setlength{\belowdisplayshortskip}{3pt}}
\fancyhf{}
\fancyhead[L]{{\itshape\CJKfamily{sourceheader}\rightmark}}
\fancyhead[R]{{\itshape\CJKfamily{sourceheader}\leftmark}}
\fancyfoot[C]{\thepage}
\renewcommand{\headrulewidth}{0.4pt}
\renewcommand{\headrule}{{\color{structurecolor}\hrule height\headrulewidth width\headwidth\vskip-\headrulewidth}}
\pagestyle{fancy}
\fancypagestyle{cover}{\fancyhf{}\fancyfoot[C]{\thepage}\renewcommand{\headrulewidth}{0pt}\renewcommand{\headrule}{}}
\RequirePackage{hyperref}
\hypersetup{unicode=true,colorlinks=true,linkcolor=structurecolor,urlcolor=structurecolor,citecolor=structurecolor,linktoc=section,bookmarksnumbered=true,bookmarksopen=true,pdfborder={0 0 0},pdfstartview=FitH}
\setcounter{tocdepth}{1}
\renewcommand*\l@section{\@dottedtocline{1}{1.5em}{2.4em}}
\renewcommand{\thesection}{\thechapter.\arabic{section}}
\newcommand{\originalchapter}[1]{\par\addvspace{14pt}\Needspace{5\baselineskip}\refstepcounter{chapter}\setcounter{section}{0}\addcontentsline{toc}{chapter}{\protect\numberline{\thechapter}#1}\markboth{\thechapter\quad #1}{}\begingroup\centering\color{structurecolor}\bfseries\fontsize{14.4}{20}\selectfont\thechapter\quad #1\par\endgroup\nobreak\vspace{8pt}}
\newcommand{\originalsection}[1]{\par\addvspace{9pt}\Needspace{4\baselineskip}\refstepcounter{section}\addcontentsline{toc}{section}{\protect\hspace*{1.5em}\protect\numberline{\protect\textcolor{structurecolor}{\thesection}}#1}\markright{\thesection\quad #1}\noindent{\fontsize{12}{17}\selectfont\color{structurecolor}\thesection\quad}{\bfseries\fontsize{12}{17}\selectfont #1}\par\nobreak\vspace{4pt}}

% Semantic environments are registered once here, not re-designed in manuscripts.
\newtheoremstyle{wzstatement}{6pt}{5pt}
  {\normalfont\kaishu}{0pt}{\normalfont\bfseries\color{structurecolor}}
  {}{.7em}{\thmname{#1}\thmnumber{ #2}\thmnote{\normalfont\ (#3)}}
\newtheoremstyle{wzdefinition}{6pt}{5pt}
  {\normalfont\songti}{0pt}{\normalfont\bfseries\color{structurecolor}}
  {}{.7em}{\thmname{#1}\thmnumber{ #2}\thmnote{\normalfont\ (#3)}}
\newtheoremstyle{wzproblem}{7pt}{5pt}
  {\normalfont\songti}{2em}{\normalfont\bfseries\color{main}}
  {}{.7em}{\thmname{#1}\thmnumber{ #2}}
\newtheoremstyle{wzremark}{4pt}{4pt}
  {\normalfont\songti}{0pt}{\normalfont\bfseries}
  {}{.7em}{\thmname{#1}}
\theoremstyle{wzstatement}
\newtheorem{theorem}{定理}[chapter]
\newtheorem{lemma}[theorem]{引理}
\newtheorem{proposition}[theorem]{命题}
\newtheorem{corollary}[theorem]{推论}
\theoremstyle{wzdefinition}
\newtheorem{definition}[theorem]{定义}
\theoremstyle{wzproblem}
\newtheorem{example}{例题}[chapter]
\newtheorem{exercise}{习题}[chapter]
\theoremstyle{wzremark}
\newtheorem*{remark}{注}
\renewcommand{\proofname}{证明}
% Retain the native amsthm QED stack; only adjust the fixed typography.
\expandafter\patchcmd\csname\string\proof\endcsname{\normalfont \topsep6\p@\@plus6\p@\relax}
  {\normalfont\songti \topsep5\p@\@plus1\p@\relax}
  {}{\ClassError{elegantbook-original-adapter}{Cannot patch proof spacing}{Check the amsthm version.}}
\expandafter\patchcmd\csname\string\proof\endcsname{\itshape}{\normalfont\bfseries}
  {}{\ClassError{elegantbook-original-adapter}{Cannot patch proof heading}{Check the amsthm version.}}
\expandafter\patchcmd\csname\string\proof\endcsname{\ignorespaces}
  {\setlength{\parindent}{2em}\ignorespaces}
  {}{\ClassError{elegantbook-original-adapter}{Cannot patch proof paragraphs}{Check the amsthm version.}}
\AtBeginEnvironment{proof}{\par\Needspace{3\baselineskip}}
\renewcommand{\qedsymbol}{\ensuremath{\square}}
\newenvironment{solution}{\begin{proof}[解答]}{\end{proof}}
\newcommand{\wzcheckchapter}{\ifnum\value{chapter}<1
  \ClassError{elegantbook-original-adapter}{Numbered mathematics before chapter 1}{Move it after the first originalchapter.}\fi}

\newcommand{\wzregisterguard}[1]{\AtBeginEnvironment{#1}{\wzcheckchapter\par\Needspace{3\baselineskip}}}
\forcsvlist{\wzregisterguard}{theorem,lemma,proposition,corollary,definition,example,exercise}
% Front matter and bibliography are not numbered mathematical chapters.
\newcommand{\originalpreface}{\par\addvspace{10pt}\Needspace{5\baselineskip}
  \noindent{\bfseries\fontsize{12}{17}\selectfont 前言}\par\nobreak\vspace{4pt}}
\newcommand{\originalcontents}{\par\addvspace{14pt}
  \begin{center}{\color{structurecolor}\bfseries\fontsize{14.4}{20}\selectfont 目录}\end{center}
  \vspace{-10pt}\@starttoc{toc}}
\renewcommand{\bibname}{参考文献}
\renewenvironment{thebibliography}[1]{
  \par\addvspace{12pt}\Needspace{3\baselineskip}\phantomsection
  \addcontentsline{toc}{chapter}{\bibname}
  \markboth{\bibname}{}
  \noindent{\color{structurecolor}\bfseries\fontsize{14.4}{20}\selectfont\bibname}\par\nobreak\vspace{6pt}
  \list{\@biblabel{\@arabic\c@enumiv}}
    {\settowidth\labelwidth{\@biblabel{#1}}\leftmargin\labelwidth
     \advance\leftmargin\labelsep\usecounter{enumiv}
     \let\p@enumiv\@empty\renewcommand\theenumiv{\@arabic\c@enumiv}
     \setlength{\itemsep}{3pt}\setlength{\parsep}{0pt}}
  \sloppy\clubpenalty4000\widowpenalty4000\sfcode`\.=1000\relax
}{\def\@noitemerr{\@latex@warning{Empty bibliography}}\endlist}

\numberwithin{equation}{chapter}
\renewcommand{\theequation}{\thechapter.\arabic{equation}}
\allowdisplaybreaks[2]
\emergencystretch=2em
\clubpenalty=6000
\widowpenalty=6000
\raggedbottom
\endinput
\end{filecontents*}
% WZ-LOCKED-CLASS-END
\documentclass{elegantbook-original-adapter}
% ===== 元数据内容槽 =====
\newcommand{\BookTitle}{泛函分析}
\hypersetup{pdftitle={泛函分析: MIT 18.102 中文重构本},pdfauthor={Casey Rodriguez; Andrew Lin; 中文译编}}
\usepackage{tikz,pgfplots,booktabs,longtable}
\pgfplotsset{compat=1.18}
\usetikzlibrary{arrows.meta,calc,positioning}
\newcommand{\R}{\mathbb R}
\newcommand{\C}{\mathbb C}
\newcommand{\K}{\mathbb K}
\newcommand{\N}{\mathbb N}
\newcommand{\Z}{\mathbb Z}
\newcommand{\Q}{\mathbb Q}
\newcommand{\B}{\mathcal B}
\newcommand{\dd}{\,\mathrm d}
\newcommand{\ii}{\mathrm i}
\newcommand{\ee}{\mathrm e}
\newcommand{\eps}{\varepsilon}
\newcommand{\ind}{\mathbf 1}
\newcommand{\norm}[1]{\lVert #1\rVert}
\newcommand{\ip}[2]{\langle #1,#2\rangle}
\DeclareMathOperator{\ran}{ran}
\DeclareMathOperator{\dist}{dist}
\DeclareMathOperator{\Span}{span}
\DeclareMathOperator{\esssup}{ess\,sup}
\begin{document}
\frontmatter
\begin{center}
{\color{structurecolor}\bfseries\fontsize{18}{25}\selectfont\BookTitle\par}
{\bfseries MIT 18.102 中文重构本\par}
Casey Rodriguez 授课; Andrew Lin 原笔记\par
Introduction to Functional Analysis, Spring 2021\par
\end{center}
\originalpreface
函数、数列和算子都可以作为向量. 赋范空间用一个数衡量向量的大小, 完备性保证无限过程仍留在空间内, 对偶空间提供可连续观测这些向量的标量函数. 本书从这些结构出发, 经 Banach 空间的基本定理、Hilbert 空间的正交分解和弱收敛, 进入紧算子、谱分解及边值问题. 读者需会线性代数、一元实分析、度量空间的收敛和连续性; 不要求事先学过无限维谱理论.

主要依据 MIT OpenCourseWare 的 \href{https://ocw.mit.edu/courses/18-102-introduction-to-functional-analysis-spring-2021/pages/lecture-notes-and-readings/}{18.102 Introduction to Functional Analysis, Spring 2021}, 授课者 Casey Rodriguez, 完整课堂笔记作者 Andrew Lin. 实际主源为125页合订本与23讲独立讲义; 二者是同一套内容的不同排版, 不重复计算内容. 官方页面一处写 ``Fall 2021'', 但课程元数据、合订本扉页和源码均标 Spring 2021, 本书据后者署名. Richard Melrose 的2020讲义仅作辅助核对, 不把其更大的内容范围冒充已经覆盖.

本书按依赖关系重排, 合并重复复习, 补全原稿的练习性证明和省略计算, 校正若干笔误. 书末列出逐讲对应、改编范围及勘误. 对偶的弱与弱星拓扑、可分情形的弱星序列紧性、完备化及若干例题为\textbf{编者补充}, 在相应位置标明. 这些补充支持本书的数学主线, 不表示原课程逐一讲授过.

测度是此处的先修接口. 第5章给出可测集、积分与几乎处处等价的必要定义, 证明所用收敛定理和 $L^p$ 完备性; 不重复逐步构造 Lebesgue 测度、Vitali 集及整本实变理论. 完整构造留待本归档的\href{https://github.com/xhc144/mit-ocw-notes/tree/main/subjects/real-analysis}{实变分析科目目录}, 当前链接只指定后续归档位置, 不宣称该科已完成. 本册独立讲清泛函分析部分所需的条件与用法. 本次另收入全部官方作业与考核材料, 测度构造中作业实际要求的平移、缩放、开集分解、可测集逼近及 Littlewood 原则在题解中补齐, 不据此宣称重编了整本实变课程.

原件依 MIT OCW 的 \href{https://ocw.mit.edu/pages/privacy-and-terms-of-use/}{使用条款}和\href{https://creativecommons.org/licenses/by-nc-sa/4.0/}{CC BY-NC-SA 4.0}非商业归档与改编. 逐文件核对实际字节、页数和单独权利标记后保存原件; 未发现主源中的第三方限制标记. 原 Fejér 示意图不复用, 本册改以公式生成的原生图展示. 中文翻译、重排、增补证明、勘误和配套题解采用同一许可, 不表示 MIT、授课者或原笔记者认可本整理本. 原署名、来源和校验和见 source-manifest.json 与原件目录. 版本日期: 2026-10-08.
\clearpage
\phantomsection\addcontentsline{toc}{chapter}{课程信息}
\noindent{\bfseries\fontsize{12}{17}\selectfont 课程信息}\par
MIT 18.102/18.1021, Introduction to Functional Analysis, 2021年春季, 本科课程, Casey Rodriguez 授课. 原课堂笔记由 Andrew Lin 整理. 官方\href{https://ocw.mit.edu/courses/18-102-introduction-to-functional-analysis-spring-2021/pages/syllabus/}{大纲}规定每周两次讲课, 每次1.5小时. 页面未公布具体星期、钟点或教室, 本书不补造这些信息.

先修要求为以下两组各一门: 18.06、18.700、18.701 中的一门线性代数课, 以及18.100A、18.100B、18.100P、18.100Q 中的一门实分析课; 也可经教师许可选课. 不指定教材, 官方提供 Richard Melrose 2020年讲义及2021年课堂笔记. 课程由赋范空间、完备性、线性泛函和 Hahn--Banach, 经对偶与有界算子, 进入 Lebesgue 测度、积分、$L^p$ 空间、Hilbert 空间、紧算子和自伴谱理论.

\href{https://ocw.mit.edu/courses/18-102-introduction-to-functional-analysis-spring-2021/pages/calendar/}{日历}列出23讲: 第1--5讲为 Banach 空间基本理论; 第6--12讲为双对偶、外测度、可测集、可测函数和积分收敛; 第13--18讲为 $L^p$、Hilbert、Fourier、投影、Riesz 表示和伴随; 第19--23讲为紧算子、谱和 Dirichlet 问题. 原日历将作业1--5依次列在第2、4、6、8、10讲, 期中列在第11与12讲之间, 作业6--10依次列在第14、16、18、20、22讲, 期末作业列在课程最后. 这是教学次序记录, 不能据此推算未公布的公历截止日期.

作业占总评50\%, 期中25\%, 期末作业25\%. 最低一次作业成绩不计; 评分界限不会高于通常的90分为 A、80分为 B. 作业允许讨论, 但每人须独立写出解答并在右上角注明合作者; 通常不收迟交, 延期须邮件向教师提出. 多数原作业规定经 Gradescope 在24:00前提交; 作业6仅注明 Gradescope, 作业8未写提交说明, 不能据其他份补造其截止时间. 大纲规定第7周有24小时带回完成的期中测验, 课程结束有48小时期末作业.

期中与期末原PDF允许查阅 Richard Melrose 讲义、教师手写笔记、Andrew Lin 笔记、Royden 的 \textit{Real Analysis}（如已购买）、既有作业解答、Piazza 讨论与录播, 禁止与其他学生或互联网协作; 违反者按原文所指 MIT Institute Policy 10.2 处理. 这些是2021年课程的历史考核规则, 本书题解供归档后的学习核对使用.

官方\href{https://ocw.mit.edu/courses/18-102-introduction-to-functional-analysis-spring-2021/pages/assignments-and-exams/}{考核目录}提供10份作业、1份期中和1份 Final Assignment, 共54道编号题、87个显式字母子问及18道无字母单题, 合计105个解答单元. 无字母单题即原编号题未划分(a)、(b)等的题目; 期末第2题虽含两个积分, 按原编号计一个单元, 两个积分均分别解答. 作业10第6题虽注明“不提交”, 仍完整收入. 期末PDF页眉写作 Final, 开头沿用 exam 措辞; 大纲与资源页明确其身份是期末作业, 本书据此命名. 未提供官方答案; 以下新增题解均标为\textbf{AI 独立解答}, 并复用正文已证明的结果. 原例题与习题解答保留, 与官方题目的对应见书末. 本册不以这些解答冒充教师原答案.

\clearpage
\originalcontents
\clearpage
\mainmatter
\originalchapter{赋范空间与完备性}\label{ch:norm}
\originalsection{范数与收敛}
在有限维问题中, 求解线性方程只涉及有限次运算. 对函数或无穷数列, 近似解往往形成一个序列; 能否把近似解的极限仍视为合法解, 取决于空间的范数与完备性. 全书的底域 $\K$ 为 $\R$ 或 $\C$, 讨论谱时会明确采用复 Hilbert 空间.

\begin{definition}
设 $X$ 为 $\K$ 上的向量空间. 范数是函数 $\norm{\cdot}:X\to[0,\infty)$, 满足 $\norm{x}=0$ 当且仅当 $x=0$, $\norm{\alpha x}=|\alpha|\norm{x}$, 以及 $\norm{x+y}\le\norm{x}+\norm{y}$. 半范数保留齐次性与三角不等式, 允许非零向量的半范数为零. 配备范数的空间称赋范空间.
\end{definition}
向量可以是数列或函数, 所以范数并不只表示欧氏长度. 范数导出距离 $d(x,y)=\norm{x-y}$: 对称性来自 $\norm{-z}=\norm{z}$, 三角不等式来自 $x-z=(x-y)+(y-z)$. 下文 $B_X(x,r)$ 表示 $\norm{y-x}<r$ 的开球, $\overline B_X(x,r)$ 表示 $\norm{y-x}\le r$ 的闭球.

\begin{proposition}\label{prop:normcontinuous}
赋范空间中有 $|\norm{x}-\norm{y}|\le\norm{x-y}$. 向量加法和数乘连续. 特别地, 若 $x_n\to x$, 则 $\norm{x_n}\to\norm{x}$, 且序列 $(x_n)$ 有界.
\end{proposition}
\begin{proof}
由 $\norm{x}\le\norm{x-y}+\norm{y}$ 得一个方向, 交换 $x,y$ 得另一个方向. 若 $x_n\to x$, $y_n\to y$, 则 $\norm{x_n+y_n-x-y}\le\norm{x_n-x}+\norm{y_n-y}\to0$. 若同时 $\alpha_n\to\alpha$, 则
\[
\norm{\alpha_n x_n-\alpha x}\le |\alpha_n|\norm{x_n-x}+|\alpha_n-\alpha|\norm{x}\to0,
\]
因为收敛标量序列有界. 对收敛向量序列, 充分大的 $n$ 满足 $\norm{x_n-x}<1$, 前面有限项再取最大值, 就得全序列的界.
\end{proof}

\begin{example}\label{legacy:ch01:example1}
比较 $C([0,1])$ 上的两个范数 $\norm{f}_\infty=\max|f|$ 与 $\norm{f}_1=\int_0^1|f|\dd x$.
\end{example}
\begin{solution}
二者均满足齐次性和三角不等式. 连续函数若不恒为零, 在某个小区间上有 $|f|\ge c>0$, 因而积分正, 所以第二个也是范数. 对 $f_n(x)=x^n$, 有 $\norm{f_n}_\infty=1$, 而 $\norm{f_n}_1=1/(n+1)\to0$. 因此同一向量空间的不同范数可给出不同的收敛概念.
\end{solution}

\begin{definition}
若对每个 $\eps>0$ 存在 $N$, 使 $n,m\ge N$ 时 $\norm{x_n-x_m}<\eps$, 则 $(x_n)$ 为 Cauchy 序列. 每个 Cauchy 序列都有空间内的极限时, 空间称为 Banach 空间.
\end{definition}
这一条件不是说所有序列收敛. 它只把已经互相接近的序列的极限补齐. 选择不同范数可能改变完备性, 不能只凭向量空间的名称判断.

\originalsection{函数空间与数列空间}
若 $S=\varnothing$, 将唯一空函数的范数规定为零, 得到零 Banach 空间.
\begin{theorem}\label{thm:cb}
设 $S$ 为非空度量空间. 有界连续函数空间 $C_b(S)$ 在上确界范数 $\norm{f}_\infty=\sup_{s\in S}|f(s)|$ 下是 Banach 空间.
\end{theorem}
\begin{proof}
设 $(f_n)$ 为 Cauchy 序列. 对每个固定 $s$, 有 $|f_n(s)-f_m(s)|\le\norm{f_n-f_m}_\infty$, 故标量完备性给出 $f(s)=\lim_n f_n(s)$. Cauchy 序列有界的证明同命题\ref{prop:normcontinuous}, 所以存在 $M$ 使每个 $\norm{f_n}_\infty\le M$, 取点态极限得 $|f(s)|\le M$.

给定 $\eps>0$, 取 $N$ 使 $n,m\ge N$ 时上确界差小于 $\eps/2$. 固定 $n\ge N$, 令 $m\to\infty$, 得对所有 $s$ 同时有 $|f_n(s)-f(s)|\le\eps/2$. 因而 $f_n\to f$ 一致. 为核对连续性, 固定 $s_0$, 先取 $n$ 使一致误差小于 $\eps/3$, 再用 $f_n$ 在 $s_0$ 的连续性控制 $|f_n(s)-f_n(s_0)|<\eps/3$. 三项相加得 $|f(s)-f(s_0)|<\eps$. 因此 $f\in C_b(S)$, 且收敛发生在所给范数中.
\end{proof}

\begin{definition}
对于 $1\le p<\infty$, 定义 $\ell^p=\{a=(a_j)_{j\ge1}:\sum_j|a_j|^p<\infty\}$, 范数为 $\norm{a}_p=(\sum_j|a_j|^p)^{1/p}$. 定义 $\ell^\infty=\{a:\sup_j|a_j|<\infty\}$, 范数为 $\norm{a}_\infty=\sup_j|a_j|$. 子空间 $c_0$ 由趋于零的数列组成, $c_{00}$ 由只有有限项非零的数列组成.
\end{definition}
这些集合对线性运算封闭. 当 $p>1$ 时, 三角不等式需要 Hölder 不等式支持, 不能只把求和号当有限求和使用.

\begin{lemma}[Hölder 与 Minkowski 不等式]\label{lem:holder-seq}
设 $1<p<\infty$, $q=p/(p-1)$. 则 $\sum_j|a_jb_j|\le\norm{a}_p\norm{b}_q$, 并且 $\norm{a+b}_p\le\norm{a}_p+\norm{b}_p$. 对 $p=1,q=\infty$ 也有相应 Hölder 不等式, $p=1,\infty$ 的三角不等式直接成立.
\end{lemma}
\begin{proof}
对 $s,t\ge0$, 有 Young 不等式 $st\le s^p/p+t^q/q$. 固定 $t$ 时, 函数 $s^p/p-st$ 的最小值在 $s=t^{q-1}$ 取得, 代入即得此式. 若两个范数都非零, 将 $a,b$ 分别除以各自范数, 对 Young 不等式求和, 得归一化后的和不超过 $1/p+1/q=1$. 范数为零的情形直接成立; 无穷和由有限部分和取极限得到.

先对有限项证明 Minkowski. 记 $s=\norm{a+b}_p$. 若 $s=0$ 无须证明. 否则
\[
s^p\le\sum_j(|a_j|+|b_j|)|a_j+b_j|^{p-1}
\le(\norm{a}_p+\norm{b}_p)s^{p-1},
\]
第二步分别用 Hölder, 并用 $(p-1)q=p$. 除以 $s^{p-1}$ 得结论. 对截断后的数列应用此式, 令截断长度趋于无穷, 同时说明 $a+b\in\ell^p$. 当 $p=1$ 逐项相加, 当 $p=\infty$ 逐项取上确界.
\end{proof}

\begin{theorem}\label{thm:lpseqbanach}
$\ell^p$ 对每个 $1\le p\le\infty$ 都完备. $c_0$ 是 $\ell^\infty$ 的闭子空间, 因而也完备.
\end{theorem}
\begin{proof}
设 $(a^{(n)})$ 为 $\ell^p$ 中的 Cauchy 序列. 因为 $|a_j^{(n)}-a_j^{(m)}|\le\norm{a^{(n)}-a^{(m)}}_p$, 每个坐标有极限 $a_j$. 当 $p<\infty$ 时, 若 $n,m\ge N$ 的范数差小于 $\eps$, 则固定 $n\ge N$ 和有限的 $J$, 令 $m\to\infty$ 得
\[
\sum_{j=1}^J|a_j^{(n)}-a_j|^p\le\eps^p.
\]
再令 $J\to\infty$, 得 $a^{(n)}-a\in\ell^p$ 且其范数不超过 $\eps$. 取一个这样的 $n$ 并用 Minkowski, 可知 $a\in\ell^p$, 于是得到范数收敛. 当 $p=\infty$, 先固定 $j$ 取 $m$ 的极限, 再对 $j$ 取上确界, 得同样结论.

若 $a^{(n)}\in c_0$ 且一致趋于 $a$, 则给定 $\eps>0$, 先取 $n$ 使 $\norm{a-a^{(n)}}_\infty<\eps/2$, 再取 $J$ 使 $j>J$ 时 $|a_j^{(n)}|<\eps/2$. 所以 $|a_j|<\eps$, 即 $a\in c_0$. 闭子空间的完备性将在下一章一般证明.
\end{proof}
有限截断在 $\ell^p$ 中趋于原数列, 当且仅当这里 $p<\infty$, 或所给数列属于 $c_0$ 并采用上确界范数. 对常数数列 $(1,1,\ldots)$, 有限截断的上确界误差始终为1.

\originalsection{绝对收敛与有限维}
\begin{theorem}\label{thm:series}
赋范空间 $X$ 完备, 当且仅当每个满足 $\sum_n\norm{x_n}<\infty$ 的向量级数 $\sum_n x_n$ 在 $X$ 中收敛.
\end{theorem}
\begin{proof}
若 $X$ 完备, 则部分和之差的范数不超过对应的标量尾和, 故部分和为 Cauchy 序列并收敛. 反之, 设绝对收敛的向量级数均收敛, 且 $(u_n)$ 为 Cauchy 序列. 递增选取 $n_k$, 使任意 $n,m\ge n_k$ 时 $\norm{u_n-u_m}<2^{-k}$. 于是 $\sum_k(u_{n_{k+1}}-u_{n_k})$ 绝对收敛. 其前 $m$ 项为 $u_{n_{m+1}}-u_{n_1}$, 所以 $(u_{n_k})$ 收敛于某个 $u\in X$. 对 $n,n_k$ 足够大,
$\norm{u_n-u}\le\norm{u_n-u_{n_k}}+\norm{u_{n_k}-u}$, 两项都可任意小, 因而原序列也收敛.
\end{proof}

\begin{theorem}\label{thm:finite}
有限维向量空间上的任意两个范数等价: 存在 $c,C>0$ 使 $c\norm{x}_a\le\norm{x}_b\le C\norm{x}_a$. 有限维赋范空间完备, 其闭有界集紧, 且任何赋范空间的有限维子空间都闭.
\end{theorem}
\begin{proof}
零维情形显然, 以下设维数 $d\ge1$. 取基 $v_1,\ldots,v_d$, 以 $|t|_2$ 表示坐标 $t\in\K^d$ 的欧氏范数. 对任意所给范数,
$\norm{\sum_jt_jv_j}\le\sum_j|t_j|\norm{v_j}\le C|t|_2$. 因而该范数是坐标的连续函数. 欧氏单位球面紧, 范数在其上达到最小值 $c$; 最小值不能为零, 因为基向量线性无关且单位坐标不是零. 齐次性给出 $c|t|_2\le\norm{\sum_jt_jv_j}$. 两个范数均与欧氏范数比较, 得等价性.

给定范数下的 Cauchy 序列坐标也是 Cauchy 序列, 故逐坐标收敛; 范数上界保证收敛回到原空间. 同样的比较把闭有界集化为 $\R^{d}$ 或 $\R^{2d}$ 的闭有界集, 可用 Heine--Borel 定理. 若有限维子空间中的序列在环境空间收敛, 则它在子空间范数中为 Cauchy, 完备性给出子空间内的极限; 环境空间的极限唯一, 所以子空间闭.
\end{proof}

\begin{exercise}\label{legacy:ch01:exercise1}
证明 $c_{00}$ 在 $\ell^2$ 范数下不完备, 但在 $\ell^2$ 中稠密.
\end{exercise}
\begin{solution}
令 $u^{(n)}=\sum_{j=1}^n j^{-1}e_j$. 平方尾和 $\sum_{j>n}j^{-2}\to0$, 所以此序列为 Cauchy. 其 $\ell^2$ 极限为 $(j^{-1})$, 不在 $c_{00}$, 由极限唯一性知在 $c_{00}$ 中没有极限. 任意 $a\in\ell^2$ 的截断 $P_na$ 满足 $\norm{a-P_na}_2^2=\sum_{j>n}|a_j|^2\to0$, 所以稠密.
\end{solution}

\begin{exercise}\label{legacy:ch01:exercise2}
若 $1\le p<q\le\infty$, 证明 $\ell^p\subset\ell^q$ 且 $\norm{a}_q\le\norm{a}_p$. 说明有限测度上的 $L^p$ 包含关系为什么不应照抄此方向.
\end{exercise}
\begin{solution}
先设 $\norm{a}_p=1$, 则每个 $|a_j|\le1$. 当 $q<\infty$, 有 $\sum|a_j|^q\le\sum|a_j|^p=1$; 当 $q=\infty$ 直接用逐项界. 一般情形由归一化得到. 计数测度下各单点测度为1, 而有限区间允许函数在很小集合上很大. 例如 $x^{-1/3}$ 在 $(0,1)$ 属于 $L^2$ 而不属于 $L^3$, 因为对应幂积分分别有限与发散. 有限测度上的正确包含方向在第\ref{ch:measure}章给出.
\end{solution}

\originalchapter{有界算子与空间构造}\label{ch:operators}
\originalsection{线性算子的连续性}
有限维矩阵的每一行都是连续线性函数. 在无限维空间, 线性只约束有限次加法和数乘, 不自动约束极限. 分析中需要的算子因此还必须满足范数估计.

\begin{theorem}\label{thm:bounded}
设 $X,Y$ 为赋范空间, $T:X\to Y$ 为处处定义的线性映射. 以下条件等价: $T$ 连续; $T$ 在零点连续; 存在 $C\ge0$ 使 $\norm{Tx}_Y\le C\norm{x}_X$ 对所有 $x$ 成立. 满足这些条件的映射称有界线性算子.
\end{theorem}
\begin{proof}
连续当然蕴含零点连续. 若在零点连续, 取 $r>0$ 使 $\norm{x}<r$ 时 $\norm{Tx}<1$. 对 $x\ne0$, 向量 $rx/(2\norm{x})$ 的范数为 $r/2$, 因而由线性性得 $\norm{Tx}<2\norm{x}/r$. $x=0$ 时也成立. 最后, 若有统一估计, 则 $\norm{Tx-Ty}\le C\norm{x-y}$, 所以处处连续.
\end{proof}
这里的有界是把有界集合映成有界集合. 一个非零线性算子在整个空间上的像不可能有界, 因为可以不断放大输入.

\begin{definition}
所有有界线性算子组成的空间记为 $\B(X,Y)$. 算子范数为 $\norm{T}=\sup_{\norm{x}\le1}\norm{Tx}$. 若 $X\ne\{0\}$, 也等于 $\sup_{\norm{x}=1}\norm{Tx}$. 标量值算子空间 $X^*=\B(X,\K)$ 称连续对偶空间, 其元素称连续线性泛函.
\end{definition}
算子范数是使 $\norm{Tx}\le C\norm{x}$ 成立的最小常数. 由逐向量的三角不等式和齐次性, 得 $\norm{S+T}\le\norm{S}+\norm{T}$, $\norm{\alpha T}=|\alpha|\norm{T}$. 若 $\norm{T}=0$, 所有单位向量映成零, 缩放后所有向量都映成零, 因而这确实是范数.

\begin{proposition}\label{prop:composition}
若 $T\in\B(X,Y)$, $S\in\B(Y,Z)$, 则 $ST\in\B(X,Z)$ 且 $\norm{ST}\le\norm{S}\norm{T}$.
\end{proposition}
\begin{proof}
复合保持线性, 且 $\norm{STx}\le\norm{S}\norm{Tx}\le\norm{S}\norm{T}\norm{x}$. 对单位球取上确界即得结论.
\end{proof}

\begin{theorem}\label{thm:operatorbanach}
若 $Y$ 完备, 则 $\B(X,Y)$ 完备, 不要求 $X$ 完备. 因而任意赋范空间的连续对偶都是 Banach 空间.
\end{theorem}
\begin{proof}
设 $(T_n)$ 在算子范数下为 Cauchy 序列. 对每个 $x\in X$, 有 $\norm{T_nx-T_mx}\le\norm{T_n-T_m}\norm{x}$, 所以 $Y$ 的完备性允许定义 $Tx=\lim_nT_nx$. 极限与有限次线性运算交换, 得 $T$ 线性. Cauchy 序列 $(T_n)$ 的范数有统一界 $M$, 故取极限得 $\norm{Tx}\le M\norm{x}$.

给定 $\eps>0$, 取 $N$ 使 $n,m\ge N$ 时 $\norm{T_n-T_m}<\eps$. 固定 $n\ge N$, 对每个 $x$ 令 $m\to\infty$, 得 $\norm{(T_n-T)x}\le\eps\norm{x}$. 再对单位球取上确界, 得 $\norm{T_n-T}\le\eps$. 这一步保证的是统一于单位球的收敛, 比逐个 $x$ 的收敛更强.
\end{proof}

\begin{example}\label{legacy:ch02:example1}
设 $k\in C([0,1]^2)$, 定义 $(Tf)(x)=\int_0^1 k(x,y)f(y)\dd y$. 证明 $T$ 在 $C([0,1])$ 上有界. 比较微分算子在不同范数下的行为.
\end{example}
\begin{solution}
积分线性. 核的一致连续性给出
$|Tf(x)-Tf(z)|\le\norm{f}_\infty\sup_y|k(x,y)-k(z,y)|\to0$, 故输出连续. 又有 $\norm{Tf}_\infty\le\norm{k}_\infty\norm{f}_\infty$, 从而 $\norm{T}\le\norm{k}_\infty$.

对 $D:C^1([0,1])\to C([0,1])$, 若定义域只给上确界范数, 则 $f_n(x)=\sin(nx)/n$ 一致趋于零, 但 $\norm{Df_n}_\infty=1$, 故 $D$ 不连续. 若定义域给 $\norm{f}_{C^1}=\norm{f}_\infty+\norm{f'}_\infty$, 则 $\norm{Df}_\infty\le\norm{f}_{C^1}$, 它变成有界算子. 算子是否有界必须连同定义域范数一起说.
\end{solution}

\originalsection{子空间与商空间}
\begin{proposition}\label{prop:closedsubspace}
Banach 空间 $X$ 的线性子空间 $M$, 在继承的范数下完备当且仅当 $M$ 闭.
\end{proposition}
\begin{proof}
若 $M$ 闭, 其中的 Cauchy 序列在 $X$ 中收敛, 极限由闭性仍在 $M$. 若 $M$ 完备, 而 $m_n\in M$ 在 $X$ 中趋于 $x$, 则 $(m_n)$ 在 $M$ 中也是 Cauchy 序列, 故在 $M$ 中趋于某个 $m$. 极限唯一给出 $x=m\in M$.
\end{proof}

\begin{definition}
设 $M$ 为 $X$ 的线性子空间. 商空间 $X/M$ 的元素为陪集 $x+M$, 加法与数乘定义为 $(x+M)+(y+M)=x+y+M$, $\alpha(x+M)=\alpha x+M$. 商半范数为 $\norm{x+M}_{X/M}=\inf_{m\in M}\norm{x-m}=\dist(x,M)$.
\end{definition}
改变代表元只平移遍历 $M$ 的变量, 所以运算与商半范数均良定义. 三角不等式来自对 $\norm{x+y-(m+n)}$ 作估计后分别逼近下确界; 非零标量的齐次性来自换元 $m=\alpha n$, 零标量单独成立.

\begin{theorem}\label{thm:quotient}
商半范数是范数当且仅当 $M$ 闭. 若 $X$ 为 Banach 空间且 $M$ 闭, 则 $X/M$ 为 Banach 空间. 商映射 $q:X\to X/M$, $q(x)=x+M$, 有界且满射, 并将开集映成开集.
\end{theorem}
\begin{proof}
商半范数为零等价于 $x\in\overline M$, 故正定性恰好要求 $M=\overline M$. 又有 $\norm{qx}\le\norm{x}$, 所以 $q$ 有界. 并且 $q(B_X(0,r))=B_{X/M}(0,r)$: 一边由估计得到, 另一边由下确界小于 $r$ 可选代表元范数小于 $r$. 平移球即得开映射性质.

证明完备性时, 使用定理\ref{thm:series}. 设 $\sum_n\norm{\xi_n}_{X/M}<\infty$. 对每个 $n$ 选代表元 $x_n$, 满足 $qx_n=\xi_n$ 且 $\norm{x_n}\le\norm{\xi_n}+2^{-n}$. 于是 $\sum_nx_n$ 在 $X$ 中绝对收敛, 得到和 $x$. $q$ 连续给出商空间部分和 $\sum_{n=1}^N\xi_n=q(\sum_{n=1}^Nx_n)\to qx$. 商空间的绝对收敛级数均收敛, 所以完备.
\end{proof}

\begin{proposition}
设 $p$ 是向量空间 $V$ 上的半范数, $N=\{x:p(x)=0\}$. 则 $N$ 为子空间, $V/N$ 上的 $\norm{x+N}=p(x)$ 是范数.
\end{proposition}
\begin{proof}
由 $p(\alpha x+\beta y)\le|\alpha|p(x)+|\beta|p(y)$ 可知零集对线性运算封闭. 若 $x-y\in N$, 则 $p(x)\le p(y)+p(x-y)=p(y)$, 反向也成立, 所以定义与代表元无关. 齐次性与三角不等式继承自 $p$, 正定性来自 $p(x)=0$ 当且仅当 $x+N=N$.
\end{proof}
后文的 $L^p$ 空间正是把积分半范数为零的函数识别成一个等价类. 这不是忽略定义的细节, 而是范数成立所必须的步骤.

\originalsection{完备化与稠密延拓}
本节为\textbf{编者补充}. 原课程主要在已经完备的空间上工作; 完备化解释了为何可以由较简单对象的近似构成新空间.

\begin{theorem}\label{thm:completion}
每个赋范空间 $X$ 都可等距嵌入一个 Banach 空间 $\widehat X$, 且其像稠密. 若 $Y$ 为 Banach 空间, 每个 $T\in\B(X,Y)$ 唯一延拓为 $\widehat T\in\B(\widehat X,Y)$, 并保持算子范数.
\end{theorem}
\begin{proof}
取 $X$ 中所有 Cauchy 序列, 将 $(x_n)$ 与 $(y_n)$ 在 $\norm{x_n-y_n}\to0$ 时视为等价. 三角不等式保证这是等价关系. 用逐项线性运算定义商空间的运算; 两个等价代表元之和仍等价, 因而良定义. 设 $[(x_n)]$ 的范数为 $\lim_n\norm{x_n}$. 反三角不等式保证标量极限存在且与代表元无关, 范数为零恰好表示该序列等价于零序列. 范数的其他性质由逐项取极限得到.

把 $x$ 映到常值序列得到等距嵌入 $j$. 若 $\xi=[(x_n)]$, 则 $\norm{j(x_n)-\xi}=\lim_m\norm{x_n-x_m}\to0$, 所以 $j(X)$ 稠密. 若 $(\xi_k)$ 为所构空间中的 Cauchy 序列, 用稠密性选 $x_k\in X$ 使 $\norm{j(x_k)-\xi_k}<2^{-k}$. 则 $(x_k)$ 为 $X$ 中的 Cauchy 序列. 令 $\xi=[(x_k)]$, 前面的稠密估计给出 $j(x_k)\to\xi$, 加上 $\norm{j(x_k)-\xi_k}\to0$ 得 $\xi_k\to\xi$. 因而所构空间完备.

定义 $\widehat T([(x_n)])=\lim_nTx_n$. 因 $Y$ 完备且 $T$ 有界, 极限存在. 等价代表元的像之差趋于零, 所以良定义. 取极限得线性性与 $\norm{\widehat T\xi}\le\norm{T}\norm{\xi}$. 限制回 $j(X)$ 又得反向范数比较. 唯一性来自稠密性和连续性. 以此延拓恒等嵌入也可证明任意两个完备化之间存在保持原空间的唯一等距同构.
\end{proof}

\begin{exercise}\label{legacy:ch02:exercise1}
设 $M$ 闭, $T\in\B(X,Y)$ 且 $M\subset\ker T$. 证明 $T$ 唯一分解为 $T=\widetilde Tq$, 其中 $\widetilde T:X/M\to Y$ 有界且 $\norm{\widetilde T}=\norm{T}$.
\end{exercise}
\begin{solution}
令 $\widetilde T(x+M)=Tx$. 若 $x-y\in M$, 则 $Tx=Ty$, 所以良定义; 满射的 $q$ 保证唯一性. 对每个 $m\in M$, 有 $\norm{Tx}=\norm{T(x-m)}\le\norm{T}\norm{x-m}$. 取下确界得 $\norm{\widetilde T}\le\norm{T}$. 反之 $T=\widetilde Tq$ 和 $\norm{q}\le1$ 给出 $\norm{T}\le\norm{\widetilde T}$.
\end{solution}

\begin{exercise}\label{legacy:ch02:exercise2}
证明 $C^1([0,1])$ 在 $\norm{f}_{C^1}$ 下完备.
\end{exercise}
\begin{solution}
若 $(f_n)$ 为 Cauchy 序列, 则 $f_n\to f$, $f_n'\to g$ 都一致收敛, 且 $f,g$ 连续. 由 $f_n(x)=f_n(0)+\int_0^xf_n'(t)\dd t$ 取极限, 得 $f(x)=f(0)+\int_0^xg(t)\dd t$. 这里一致误差使积分误差不超过 $\norm{f_n'-g}_\infty$, 所以极限交换有依据. 微积分基本定理给出 $f'=g$, 于是 $f\in C^1$, 且两部分范数误差均趋于零.
\end{solution}

\Needspace{9\baselineskip}
\originalchapter{Hahn--Banach 与对偶}\label{ch:dual}
\originalsection{线性泛函的延拓}
在 $\ell^p$ 中, 坐标投影已经给出非零连续泛函. 一般赋范空间没有指定坐标, 需要证明连续泛函仍能区分向量. 延拓定理从一个易于定义的小子空间出发, 在保持范数的同时把泛函扩展到全空间.

\begin{definition}
实向量空间上的函数 $p:X\to\R$ 称次线性函数, 若 $p(x+y)\le p(x)+p(y)$, 且 $t\ge0$ 时 $p(tx)=tp(x)$. 它不必为非负或偶函数. 范数及正数倍的范数都是次线性函数.
\end{definition}

\begin{lemma}\label{lem:hb-one}
设 $M$ 为实向量空间 $X$ 的子空间, $f:M\to\R$ 线性, 且 $f(m)\le p(m)$. 若 $x_0\notin M$, 则 $f$ 可延拓到 $M+\R x_0$, 并仍由 $p$ 控制.
\end{lemma}
\begin{proof}
每个新向量唯一写成 $m+tx_0$, 因为若两种表示相等且系数不同, 则 $x_0$ 将属于 $M$. 因而只需选择 $a=F(x_0)$, 再令 $F(m+tx_0)=f(m)+ta$. 当 $t>0$ 时, 除以 $t$ 后的条件为 $a\le p(m+x_0)-f(m)$ 对所有 $m\in M$ 成立. 当 $t<0$ 时, 令 $s=-t>0$, 得 $a\ge f(m)-p(m-x_0)$ 对所有 $m$ 成立. 所以选择范围为
\[
\sup_{m\in M}\{f(m)-p(m-x_0)\}
\le a\le
\inf_{n\in M}\{p(n+x_0)-f(n)\}.
\]
该区间非空: 对任意 $m,n\in M$,
$f(m)+f(n)=f(m+n)\le p(m+n)\le p(m-x_0)+p(n+x_0)$.
移项就得到 $f(m)-p(m-x_0)\le p(n+x_0)-f(n)$, 所以每一个下界不超过每一个上界. 用 $m=0$ 或 $n=0$ 还可知两端为有限实数. 任选 $a$ 在其中, 即同时满足正、负 $t$ 的控制, $t=0$ 则由原假设得到.
\end{proof}

\begin{theorem}[实 Hahn--Banach 定理]\label{thm:hb-real}
在引理\ref{lem:hb-one}的条件下, 存在线性 $F:X\to\R$ 使 $F|_M=f$, 且 $F(x)\le p(x)$ 对所有 $x$ 成立.
\end{theorem}
\begin{proof}
采用 Zorn 引理: 若非空偏序集中每条链都有上界, 则有极大元. 这是本书采用的选择公理形式, 不在此证明集合论公理的等价性.

考虑所有受 $p$ 控制的延拓 $(N,g)$, 其中 $M\subset N\subset X$, 按扩张定义域且保持原值排序. 集合非空, 因为包含 $(M,f)$. 对任一链, 定义域的并集 $N_\infty$ 是子空间: 两个向量分别属于链中的两个子空间, 其中一个包含另一个, 因而这两个向量及其线性组合都属于较大的子空间. 在并集上令 $g_\infty(x)$ 等于任一包含 $x$ 的链成员的值; 链中延拓相容保证值唯一. 两个输入可放在同一成员中, 所以线性性成立; 控制不等式也逐点继承. 这给出链的上界.

Zorn 引理给出极大延拓 $(N,F)$. 若 $N\ne X$, 可取 $x_0\notin N$, 由单维延拓引理再扩展一步, 与极大性矛盾. 所以 $N=X$.
\end{proof}

\begin{theorem}[保范数延拓]\label{thm:hb}
设 $X$ 为实或复赋范空间, $M$ 为任意线性子空间, $f\in M^*$. 则存在 $F\in X^*$ 满足 $F|_M=f$ 及 $\norm{F}=\norm{f}$. 不要求 $X$ 完备或 $M$ 闭.
\end{theorem}
\begin{proof}
实情形令 $p(x)=\norm{f}\norm{x}$, 由实定理有 $F(x)\le p(x)$. 对 $-x$ 应用得 $-F(x)\le p(x)$, 所以 $|F(x)|\le\norm{f}\norm{x}$. 限制回 $M$ 得范数相等.

复情形先把 $X$ 看作实向量空间, 在 $M$ 上令 $g=\operatorname{Re}f$. 实线性泛函 $g$ 满足 $|g(m)|\le\norm{f}\norm{m}$, 可保此界延拓为实线性 $G$. 定义
$F(x)=G(x)-\ii G(\ii x)$. 它对实数线性, 且
$F(\ii x)=G(\ii x)+\ii G(x)=\ii F(x)$, 所以复线性. 在 $M$ 上, $\operatorname{Re}f(\ii m)=-\operatorname{Im}f(m)$, 故 $F(m)=f(m)$.

还必须证明原来的界保持, 不能从实、虚部分别有界直接得到同一常数. 对 $F(x)\ne0$, 取 $|\alpha|=1$ 使 $\alpha F(x)=|F(x)|$. 则 $|F(x)|=\operatorname{Re}F(\alpha x)=G(\alpha x)\le\norm{f}\norm{\alpha x}=\norm{f}\norm{x}$. 零值时也成立. 限制给出反向比较, 完成保范数证明.
\end{proof}

\begin{remark}
原第5讲的复延拓证明对实、虚方向的相容性写得很略. 上述实化后重建复泛函的论证为编者补全, 保留了每一步的复线性和范数控制. Zorn 引理也给出 Hamel 基存在: 对线性无关子集按包含排序, 链的并仍线性无关, 因为任一有限关系落在一个链成员中; 极大线性无关集必张成全空间. Hamel 展开只含有限项, 与后文的正交级数展开不同. 标准坐标向量是 $c_{00}$ 的 Hamel 基, 却不是 $\ell^1$ 的 Hamel 基.
\end{remark}

\originalsection{分离与双对偶}
\begin{corollary}\label{cor:norming}
若 $x_0\ne0$, 则存在 $f\in X^*$ 满足 $\norm{f}=1$, $f(x_0)=\norm{x_0}$. 因而 $\norm{x}=\sup_{\norm{f}\le1}|f(x)|$, 且连续泛函区分不同向量.
\end{corollary}
\begin{proof}
在 $\K x_0$ 上定义 $f_0(\alpha x_0)=\alpha\norm{x_0}$. 它的范数为1, 由保范数延拓得到 $f$. 任意单位泛函都满足 $|f(x)|\le\norm{x}$, 而所构泛函使该界达到. 对不同 $x,y$, 将同一构造用于 $x-y$ 即得 $f(x)\ne f(y)$.
\end{proof}

\begin{theorem}\label{thm:distance-dual}
设 $M$ 为闭线性子空间, $x\notin M$, $d=\dist(x,M)>0$. 存在 $f\in X^*$ 使 $f|_M=0$, $\norm{f}=1$, $f(x)=d$. 更一般地,
\[
\dist(x,M)=\sup\{|f(x)|:f\in X^*,\ \norm{f}\le1,\ f|_M=0\}.
\]
\end{theorem}
\begin{proof}
在 $M+\K x$ 上令 $g(m+\alpha x)=\alpha d$. 因 $x\notin M$, 表示唯一. 对 $\alpha\ne0$,
$\norm{m+\alpha x}=|\alpha|\norm{x+m/\alpha}\ge|\alpha|d$, 所以 $\norm{g}\le1$. 取 $m_n\in M$ 使 $\norm{x-m_n}\to d$, 则 $|g(x-m_n)|/\norm{x-m_n}\to1$, 因而 $\norm{g}=1$. 用 Hahn--Banach 延拓. 任意在 $M$ 上为零的单位泛函满足 $|f(x)|=|f(x-m)|\le\norm{x-m}$, 取下确界得上界. 若 $x\in M$, 两边均为零.
\end{proof}

\begin{definition}
双对偶为 $X^{**}=(X^*)^*$. 典范映射 $J_X:X\to X^{**}$ 定义为 $(J_Xx)(f)=f(x)$. 若 $J_X$ 满射, 则称 $X$ 自反.
\end{definition}

\begin{proposition}\label{prop:doubledual}
$J_X$ 是线性等距嵌入. 因此 Banach 空间 $X$ 在双对偶中的像闭, 但未必等于双对偶.
\end{proposition}
\begin{proof}
对每个 $x$, 关于 $f$ 的函数 $f(x)$ 线性, 且 $|f(x)|\le\norm{x}\norm{f}$, 所以属于双对偶. $J_X$ 的线性性来自 $f$ 的线性性. 推论\ref{cor:norming}给出 $\norm{J_Xx}=\norm{x}$, 因而单射. 若 $X$ 完备, 等距像也完备, 由子空间完备性判别得闭.
\end{proof}

\originalsection{数列空间的对偶}
以下把抽象对偶具体化的计算为编者补全; 原第2讲主要给出结论. 所有泛函在输入上都取线性约定, 用配对 $\sum_j a_jb_j$ 表示, 不在 $b_j$ 上额外添加共轭.

\begin{theorem}\label{thm:dual-lp}
若 $1\le p<\infty$, $1/p+1/q=1$, 则每个 $f\in(\ell^p)^*$ 唯一写成 $f(a)=\sum_ja_jb_j$, 其中 $b\in\ell^q$, 且 $\norm{f}=\norm{b}_q$. 同样 $(c_0)^*$ 等距对应于 $\ell^1$.
\end{theorem}
\begin{proof}
给定 $b\in\ell^q$, Hölder 保证级数绝对收敛及 $\norm{f_b}\le\norm{b}_q$. 对一般 $f$, 令 $b_j=f(e_j)$. 在有限支撑向量上已有 $f(a)=\sum a_jb_j$.

若 $1<p<\infty$, 固定 $N$, 令 $S_N=\sum_{j=1}^N|b_j|^q$. 当 $S_N>0$, 取有限支撑向量
$a_j=\overline{b_j}|b_j|^{q-2}/S_N^{1/p}$, 零 $b_j$ 处定义为零. 因 $(q-1)p=q$, 有 $\norm{a}_p=1$ 且 $f(a)=S_N^{1/q}$. 所以 $S_N^{1/q}\le\norm{f}$, 令 $N\to\infty$ 得 $b\in\ell^q$ 及相应界. 若 $p=1$, 由 $|b_j|=|f(e_j)|\le\norm{f}$ 得 $b\in\ell^\infty$. $c_{00}$ 在 $\ell^p$ 中稠密, 用截断收敛和 $f$ 连续性可把有限求和恒等式延至所有 $a$. 同时前面的有限支撑测试使范数反向不等式成立. 每个坐标由 $f(e_j)$ 确定, 故表示唯一.

对 $c_0$, 固定 $N$, 取 $a_j=\overline{b_j}/|b_j|$ 于 $1\le j\le N$ 的非零 $b_j$ 处, 其余取零. 此向量上确界范数不超过1, 得 $\sum_{j=1}^N|b_j|\le\norm{f}$. 所以 $b\in\ell^1$. $c_{00}$ 在 $c_0$ 中稠密, 余下表示与范数证明同前.
\end{proof}

\begin{corollary}\label{cor:lp-reflexive}
当 $1<p<\infty$ 时, $\ell^p$ 自反.
\end{corollary}
\begin{proof}
令 $q=p/(p-1)$, 则 $1<q<\infty$. 定理\ref{thm:dual-lp}把 $(\ell^p)^*$ 线性等距识别为 $\ell^q$, 其中 $b$ 对应 $f_b(a)=\sum_ja_jb_j$. 再对 $\ell^q$ 应用同一定理, 任一 $\Phi\in(\ell^p)^{**}$ 都有唯一的 $a\in\ell^p$ 使 $\Phi(f_b)=\sum_jb_ja_j$. 而右边就是 $f_b(a)=(J_{\ell^p}a)(f_b)$. 全部连续泛函均为某个 $f_b$, 所以 $\Phi=J_{\ell^p}a$, 典范映射满射. 这里核对的是典范映射, 仅知道两个空间抽象同构还不足以证明自反.
\end{proof}

\begin{example}\label{legacy:ch03:example1}
证明 $(\ell^\infty)^*$ 不由 $\ell^1$ 的上述配对穷尽, 并证明 $\ell^1$ 与 $c_0$ 不自反.
\end{example}
\begin{solution}
令 $c\subset\ell^\infty$ 为所有收敛数列, 泛函 $L(a)=\lim_j a_j$ 在 $c$ 上线性且范数为1. Hahn--Banach 把它延拓为 $F\in(\ell^\infty)^*$. 它在 $c_0$ 上为零, 且 $F(1,1,\ldots)=1$. 若 $F(a)=\sum_ja_jb_j$ 对某个 $b\in\ell^1$ 成立, 则每个 $b_j=F(e_j)=0$, 将推出 $F=0$, 矛盾.

由于 $(\ell^1)^*=\ell^\infty$, 双对偶为 $(\ell^\infty)^*$. 对 $b\in\ell^1$, 典范像 $J_{\ell^1}b$ 在 $a\in\ell^\infty$ 上的值恰为 $\sum_jb_ja_j$. 刚构造的 $F$ 不可能是这样的典范像, 所以 $\ell^1$ 不自反. 这同时说明端点 $p=1$ 不能加入上一推论.

由 $(c_0)^*=\ell^1$ 和 $(\ell^1)^*=\ell^\infty$, 双对偶可识别为 $\ell^\infty$. 典范映射把 $a\in c_0$ 送到同一坐标数列 $a$, 所以常数数列不在像中, $c_0$ 不自反.
\end{solution}

\begin{exercise}\label{legacy:ch03:exercise1}
设 $M\subset X$ 为线性子空间. 证明 $\overline M=X$ 当且仅当每个在 $M$ 上为零的连续泛函都是零.
\end{exercise}
\begin{solution}
若 $M$ 稠密, 连续泛函在 $M$ 上为零, 则对任意 $x$ 取 $m_n\to x$, 得 $f(x)=\lim f(m_n)=0$. 若 $\overline M\ne X$, 对闭子空间 $\overline M$ 及其外的一个向量应用定理\ref{thm:distance-dual}, 得到在 $M$ 上为零却不恒零的泛函.
\end{solution}

\originalchapter{Banach 空间的基本定理}\label{ch:baire}
\originalsection{Baire 定理与统一有界}
逐个输入向量都受某个界控制, 还不一定存在控制整个单位球的共同常数. 完备性在此处的作用, 是把可数个局部限制转化成一个含有开球的限制.

\begin{theorem}[Baire 类定理]\label{thm:baire}
非空完备度量空间 $S$ 中, 可数个开稠密集合 $U_n$ 的交仍稠密. 等价地, 若 $S=\bigcup_nF_n$ 且各 $F_n$ 闭, 则至少一个 $F_n$ 有非空内部.
\end{theorem}
\begin{proof}
固定任意非空开集 $O$. 因 $U_1$ 开且稠密, 可在 $O\cap U_1$ 中选一个闭球 $\overline B(s_1,r_1)$, 半径 $0<r_1<1/2$. 这是因为选定中心后开性给出一个开球, 再把半径减半即可. 递归地在 $B(s_n,r_n)\cap U_{n+1}$ 中选闭球 $\overline B(s_{n+1},r_{n+1})$, 使 $0<r_{n+1}<2^{-n-1}$. 即使某点是孤立点, 小正半径的单点球也允许这个选择.

闭球逐级包含. 当 $m\ge n$, 中心 $s_m$ 在第 $n$ 个闭球中, 故 $d(s_m,s_n)\le r_n\to0$. 中心序列为 Cauchy, 完备性给出极限 $s$. 第 $n$ 个闭球闭且包含之后的中心, 因此 $s$ 属于每个闭球. 从而 $s\in O\cap\bigcap_nU_n$. 任意 $O$ 都与交相遇, 所以交稠密.

若所有闭 $F_n$ 的内部为空, 则补集开且稠密, 其交非空, 不可能被 $\bigcup F_n$ 覆盖. 反之, 若某开稠密集合的交不稠密, 在一个非空开集内用相同的闭球构造也会得到矛盾. 所以两个形式等价.
\end{proof}

\begin{theorem}[统一有界原理]\label{thm:ub}
设 $X$ 为 Banach 空间, $Y$ 为赋范空间, $\mathcal T\subset\B(X,Y)$ 为任意算子族. 若每个 $x\in X$ 都满足 $\sup_{T\in\mathcal T}\norm{Tx}<\infty$, 则 $\sup_{T\in\mathcal T}\norm{T}<\infty$.
\end{theorem}
\begin{proof}
令 $E_k=\{x:\sup_{T\in\mathcal T}\norm{Tx}\le k\}$. 每个集合是闭的, 因为它是闭集 $\{x:\norm{Tx}\le k\}$ 的任意交. 点态有界保证 $X=\bigcup_{k\ge1}E_k$. Baire 定理给出某 $E_k$ 含有 $B(x_0,r)$.

若 $\norm{h}<r$, 则 $x_0+h,x_0\in E_k$, 所以对所有 $T$ 同时有 $\norm{Th}\le\norm{T(x_0+h)}+\norm{Tx_0}\le2k$. 对任意单位向量 $u$ 取 $h=ru/2$, 得 $\norm{Tu}\le4k/r$. 取两个上确界, 即得 $\sup_T\norm{T}\le4k/r$. 值域只用到范数和各 $T$ 的连续性, 所以无需完备.
\end{proof}

\begin{corollary}\label{cor:stronglimit}
若 $X,Y$ 均为 Banach 空间, $T_n\in\B(X,Y)$ 且每个 $T_nx$ 都收敛, 则逐点极限 $Tx=\lim_nT_nx$ 为有界线性算子. 但未必有 $\norm{T_n-T}\to0$.
\end{corollary}
\begin{proof}
标量与向量的有限线性运算可取极限, 所以 $T$ 线性. 收敛序列逐点有界, 统一有界原理给出 $M=\sup_n\norm{T_n}<\infty$. 对 $\norm{T_nx}\le M\norm{x}$ 取极限, 得 $T$ 有界. 在 $\ell^2$ 上令 $P_n$ 截断前 $n$ 项, 则 $P_nx\to x$ 对每个 $x$ 成立, 但 $\norm{P_n-I}=1$, 因为它在 $e_{n+1}$ 上的差范数为1, 且投影的差范数至多1.
\end{proof}

\begin{example}\label{legacy:ch04:example1}
在 $c_{00}$ 的上确界范数下令 $T_nx=nx_n$. 检查统一有界原理的完备性条件.
\end{example}
\begin{solution}
每个 $T_n$ 连续, 且 $\norm{T_n}=n$, 在 $e_n$ 上达到. 每个固定 $x\in c_{00}$ 只有有限项非零, 所以 $\sup_n|nx_n|$ 有限. 但算子范数不统一有界. 此空间不完备: 数列 $(1/j)$ 的有限截断一致 Cauchy, 极限不属于 $c_{00}$. 因此该例准确显示去掉定义域的完备性后结论可能失效.
\end{solution}

\originalsection{开映射与有界逆}
\begin{theorem}[开映射定理]\label{thm:open}
若 $X,Y$ 为 Banach 空间, $T\in\B(X,Y)$ 满射, 则 $T$ 为开映射.
\end{theorem}
\begin{proof}
若 $Y=\{0\}$, 结论直接成立. 否则满射性给出 $Y=\bigcup_{n\ge1}\overline{T(B_X(0,n))}$. Baire 定理与缩放说明 $A=\overline{T(B_X(0,1))}$ 含有某个球 $B_Y(y_0,r)$. 因 $A-A\subset\overline{T(B_X(0,2))}$, 对 $\norm{y}<r$ 可把 $y$ 写成 $(y_0+y)-y_0$, 得 $B_Y(0,r)\subset\overline{T(B_X(0,2))}$. 缩放得到某 $\delta>0$ 满足
\[
B_Y(0,\delta)\subset\overline{T(B_X(0,1))}.
\]
目前只得到像的闭包中有球, 不能直接删掉闭包. 给定 $\norm{y}<\delta/2$, 利用缩放后的闭包关系选 $x_1$ 使 $\norm{x_1}<1/2$, $\norm{y-Tx_1}<\delta/4$. 再选 $x_2$ 使 $\norm{x_2}<1/4$, $\norm{y-Tx_1-Tx_2}<\delta/8$. 递归得到
\[
\norm{x_k}<2^{-k},\qquad
\norm{y-T\sum_{k=1}^n x_k}<\delta2^{-n-1}.
\]
$X$ 完备保证 $x=\sum_kx_k$ 存在; 且 $\norm{x}\le\sum_k\norm{x_k}<1$. 最后的严格不等式可由第一项严格小于 $1/2$、其余项和至多 $1/2$ 看出. 连续性给出 $Tx=y$. 因而 $B_Y(0,\delta/2)\subset T(B_X(0,1))$.

对开集 $U\subset X$ 中任意 $x$, 取 $s>0$ 使 $B_X(x,s)\subset U$. 线性性与上述球关系给出 $B_Y(Tx,s\delta/2)\subset T(U)$. 所以 $T(U)$ 开. 证明中 $Y$ 的完备性用于 Baire, $X$ 的完备性用于误差级数求和.
\end{proof}

\begin{corollary}[有界逆定理]\label{cor:inverse}
Banach 空间之间的有界线性双射具有有界线性逆. 同一向量空间上的两个完备范数若满足单向比较 $\norm{x}_a\le C\norm{x}_b$, 则两个范数等价.
\end{corollary}
\begin{proof}
双射 $T$ 的逆线性. 逆映射对开集 $U\subset X$ 的逆像为 $T(U)$, 由开映射定理为开集, 所以逆连续, 进而有界. 第二个结论应用于恒等映射 $I:(X,\norm{\cdot}_b)\to(X,\norm{\cdot}_a)$; 其逆的有界性正是反向比较.
\end{proof}
有界逆的作用不仅是存在解, 还给出稳定性 $\norm{x}\le\norm{T^{-1}}\norm{Tx}$. 满射而不单射时则允许多解; 开映射证明仍保证能选到大小受右端项控制的某个解, 不保证这个选解规则线性.

\originalsection{闭图定理}
\begin{theorem}[闭图定理]\label{thm:closedgraph}
若 $X,Y$ 为 Banach 空间, $T:X\to Y$ 为处处定义的线性映射, 则 $T$ 有界当且仅当图像 $\Gamma(T)=\{(x,Tx):x\in X\}$ 在 $X\times Y$ 中闭. 乘积使用范数 $\norm{(x,y)}=\norm{x}+\norm{y}$.
\end{theorem}
\begin{proof}
乘积 Cauchy 序列的两个坐标分别为 Cauchy 序列, 坐标完备性给出乘积完备性. 若 $T$ 有界, 且 $x_n\to x$, $Tx_n\to y$, 则连续性给出 $Tx_n\to Tx$, 极限唯一说明 $y=Tx$, 所以图像闭.

若图像闭, 它是乘积 Banach 空间的闭线性子空间, 因而完备. 坐标投影 $\pi_X:\Gamma(T)\to X$, $\pi_Y:\Gamma(T)\to Y$ 都有界, 范数至多1. 第一个投影双射, 其逆 $S:x\mapsto(x,Tx)$ 由有界逆定理有界. 所以 $T=\pi_YS$ 有界.
\end{proof}
闭图条件只要求: 当 $x_n$ 和 $Tx_n$ 都已经收敛时, 第二个极限必须是 $Tx$. 它不要求事先证明 $Tx_n$ 收敛. 此简化依赖两空间完备以及算子处处定义, 不能直接用来宣布偏定义的微分算子有界.

\begin{exercise}\label{legacy:ch04:exercise1}
设 $X$ 是向量空间, 两个范数都完备. 若对任意序列, $x_n\to x$ 于第一范数且 $x_n\to y$ 于第二范数时总有 $x=y$, 证明两个范数等价.
\end{exercise}
\begin{solution}
恒等映射的图像闭: 乘积中的收敛对 $(x_n,x_n)\to(x,y)$ 恰由假设给出 $x=y$. 闭图定理说明 $I:(X,\norm{\cdot}_a)\to(X,\norm{\cdot}_b)$ 有界. 交换两个范数同样得到反向有界性, 即等价.
\end{solution}

\begin{exercise}\label{legacy:ch04:exercise2}
设 $T:X\to Y$ 为 Banach 空间之间的有界单射. 证明 $\ran T$ 闭当且仅当存在 $c>0$ 使 $\norm{Tx}\ge c\norm{x}$.
\end{exercise}
\begin{solution}
若值域闭, 则它本身完备, $T:X\to\ran T$ 为有界双射. 有界逆定理给出下界, 可取 $c=1/\norm{T^{-1}}$, 非零空间情形下此常数正; 零空间时任意正 $c$ 可用. 反之, 若 $Tx_n$ 收敛于 $y$, 下界保证 $x_n$ 为 Cauchy 序列. $X$ 完备给出 $x_n\to x$, 连续性给出 $Tx=y$, 故值域闭.
\end{solution}

\originalchapter{积分与函数空间}\label{ch:measure}
\originalsection{测度与几乎处处}
积分范数能容纳不连续函数, 并把函数列的极限保存在空间内. 本章采用已经构造好的 Lebesgue 测度, 只建立后续泛函分析需要的积分工具. 原第6--12讲从外测度构造测度的详细过程, 是本册明确交由实变科处理的部分.

\begin{definition}
集合 $S$ 上的 $\sigma$ 代数 $\mathcal A$ 包含 $S$, 对补集及可数并封闭. 测度 $\mu:\mathcal A\to[0,\infty]$ 满足 $\mu(\varnothing)=0$, 且对两两不交的可测集 $E_n$ 有 $\mu(\bigcup_nE_n)=\sum_n\mu(E_n)$. $(S,\mathcal A,\mu)$ 称测度空间.
\end{definition}
若每个零测集的任意子集也可测, 则称测度完备. 本章的积分与 $L^p$ 定理适用于一般测度空间, 但其中所有被积分的函数均须可测; 不能在非完备测度空间里把任意零测集上的任意改值自动视作可测函数. 后续使用的 Lebesgue 测度是完备的.
Lebesgue 外测度在 $\R$ 上由
$m^*(E)=\inf\{\sum_n|I_n|:E\subset\bigcup_n I_n,\ I_n\text{为开区间}\}$ 定义. 满足 $m^*(A)=m^*(A\cap E)+m^*(A\setminus E)$ 对所有 $A\subset\R$ 成立的 $E$ 称 Lebesgue 可测集. 这些集合构成 $\sigma$ 代数, $m=m^*|_{\mathcal A}$ 为完备、平移不变的测度, 区间的测度等于长度. 这些构造事实在此作为测度先修使用, 不把定义当作它们的证明.

\begin{definition}
若某性质只在一个零测集上可能不成立, 则称该性质几乎处处成立. 实值函数 $f$ 可测, 指每个 $\{x:f(x)>t\}$ 可测; 复函数可测, 指实部和虚部都可测. 允许非负函数取值 $+\infty$.
\end{definition}
有限实值或复值可测函数的和、积、绝对值仍可测, 逐点极限在存在处也可测. 例如 $\{\sup_nf_n>t\}=\bigcup_n\{f_n>t\}$, 可数上下确界由补集与并得到. 函数加法可由 $\{f+g>t\}=\bigcup_{r\in\Q}\{f>r\}\cap\{g>t-r\}$ 验证; 连续实函数作用于可测函数仍可测, 从而得到积及绝对值. 扩展实值的非负函数可作非负加法, 但不作 $\infty-\infty$ 这样的未定义运算. 几乎处处相等是等价关系, 因为可数个零测集的并仍零测.

\begin{definition}
非负简单函数是只取有限个非负值的可测函数, 可写成 $s=\sum_{j=1}^Na_j\ind_{E_j}$, 各 $E_j$ 不交. 定义 $\int s\dd\mu=\sum_ja_j\mu(E_j)$, 采用 $0\cdot\infty=0$. 对非负可测 $f$, 定义 $\int f\dd\mu=\sup_{0\le s\le f}\int s\dd\mu$, 上确界遍历非负简单函数.
\end{definition}
把两个简单函数的分割共同细化, 可核对其积分的良定义、有限加法和单调性. 对一般非负函数, 上述定义立即给出积分的单调性. 每个非负可测函数可由递增简单函数逼近: 将 $\min(f,n)$ 的值按长度 $2^{-n}$ 向下取整, 得到 $s_n=2^{-n}\lfloor2^n\min(f,n)\rfloor$, 满足 $s_n\uparrow f$.

\originalsection{积分的收敛定理}
\begin{theorem}[单调收敛]\label{thm:mct}
若 $f_n,f$ 都为非负可测函数, 且 $f_n\uparrow f$ 几乎处处, 则 $\int f_n\dd\mu\uparrow\int f\dd\mu$.
\end{theorem}
\begin{proof}
在一个包含全部例外的可测零测集 $N$ 上把所有函数改为零, 使单调关系处处成立. 这样改值保持可测性, 且不改变积分: 简单函数在 $N$ 上的积分为零, 再用上确界定义. 令 $L=\lim_n\int f_n$, 单调性给出 $L\le\int f$. 固定非负简单函数 $s\le f$ 和 $0<c<1$, 设 $A_n=\{f_n\ge cs\}$. 在 $s>0$ 的每一点, 因 $f_n\to f\ge s>cs$, 该点最终属于 $A_n$; $s=0$ 的点也属于. 所以 $A_n\uparrow S$.

测度从下连续: 若 $E_n\uparrow E$, 把 $E$ 分解为 $E_1$ 及 $E_n\setminus E_{n-1}$ 的不交并, 可数可加性得 $\mu(E_n)\uparrow\mu(E)$. 对 $s$ 的有限个值域集应用此性质, 得 $\int s\ind_{A_n}\uparrow\int s$. 因 $f_n\ge cs\ind_{A_n}$, 有 $L\ge c\int s$. 若 $\int s=\infty$ 则 $L=\infty$; 否则令 $c\uparrow1$. 再对 $s$ 取上确界, 得 $L\ge\int f$.
\end{proof}
由简单函数递增逼近和单调收敛, 对非负 $f,g$ 有 $\int(f+g)=\int f+\int g$, 因为可对两组简单逼近之和取极限. 同样有非负齐次性. 对级数部分和应用单调收敛, 得 $\int\sum_nf_n=\sum_n\int f_n$; 此处非负性允许两边为无穷大.

\begin{theorem}[Fatou 引理]\label{thm:fatou}
若 $f_n\ge0$ 可测, 则 $\int\liminf_nf_n\le\liminf_n\int f_n$.
\end{theorem}
\begin{proof}
令 $g_k=\inf_{n\ge k}f_n$, 则 $g_k\uparrow\liminf_nf_n$. 对每个 $n\ge k$, 单调性给出 $\int g_k\le\int f_n$, 故 $\int g_k\le\inf_{n\ge k}\int f_n$. 令 $k\to\infty$, 左边用单调收敛, 右边按下极限定义收敛, 即得结论.
\end{proof}

\begin{definition}
若 $\int|f|\dd\mu<\infty$, 则实或复可测函数 $f$ 可积. 实函数的正、负部分为 $f^+=\max(f,0)$、$f^-=\max(-f,0)$, 以 $f=f^+-f^-$ 定义积分; 复函数分别积分实部与虚部. 两部分都有限, 所以不出现 $\infty-\infty$. 积分线性, 且 $|\int f|\le\int|f|$.
\end{definition}
实线性可由非负加法及正负部分的分解核对. 对复三角不等式, 若积分非零, 选 $|\alpha|=1$ 使 $\alpha\int f=|\int f|$, 则 $|\int f|=\int\operatorname{Re}(\alpha f)\le\int|f|$.

\begin{theorem}[控制收敛]\label{thm:dct}
若 $f_n,f$ 都可测, $f_n\to f$ 几乎处处, 且对同一个非负可积 $g$ 有 $|f_n|\le g$ 几乎处处, 则 $f$ 可积, $\int|f_n-f|\to0$, 因而 $\int f_n\to\int f$.
\end{theorem}
\begin{proof}
把所有控制与收敛关系的例外取可数并, 在这个可测零测集上把 $f_n,f,g$ 同时置零, 不改变积分. 此后取极限得 $|f|\le g$, 所以 $f$ 可积. 非负函数 $2g-|f_n-f|$ 由 Fatou 引理满足
\[
2\int g\le\liminf_n\int(2g-|f_n-f|)
=2\int g-\limsup_n\int|f_n-f|.
\]
因为 $\int g$ 有限, 可相减, 得距离积分的上极限至多零. 每项非负, 所以该积分趋于零. 积分三角不等式给出最后结论.
\end{proof}

\begin{proposition}\label{prop:riemann-lebesgue}
有限闭区间上的连续函数的 Lebesgue 积分等于其 Riemann 积分.
\end{proposition}
\begin{proof}
先设 $f$ 为实值连续函数. 对分割 $a=x_0<\cdots<x_N=b$, 在每段取 $m_j=\min f$ 和 $M_j=\max f$, 得下、上简单函数 $s,t$, 满足 $s\le f\le t$. 端点的取值不影响积分, 因为有限集零测. 两个积分分别为 Darboux 下和 $\sum_jm_j(x_j-x_{j-1})$ 与上和 $\sum_jM_j(x_j-x_{j-1})$. 一致连续性保证分割网格足够小时每个 $M_j-m_j\le\varepsilon$, 所以上、下和之差至多 $\varepsilon(b-a)$. 这同时证明 Riemann 可积, 并把 Lebesgue 积分和 Riemann 积分夹在任意接近的同一上下界之间, 因而两者相等. 实函数有界, 正、负部分均可积; 复函数分别处理实部、虚部.

特别地, 对 $\beta>0$, 令 $g_n(x)=x^{-\beta}\ind_{[1/n,1]}(x)$ 于 $(0,1)$. 它们递增趋于 $x^{-\beta}$. 单调收敛和上述连续函数结论给出 $\int_0^1x^{-\beta}\dd x=\lim_n\int_{1/n}^1x^{-\beta}\dd x$. 当 $\beta\ne1$, 截断积分为 $(1-n^{\beta-1})/(1-\beta)$; 当 $\beta=1$, 它为 $\log n$. 故极限有限当且仅当 $\beta<1$. 因此后文的幂函数分类使用的是已建立的积分衔接, 并未把不连续端点处的广义积分直接当作定义.
\end{proof}

\begin{remark}
本书在核积分和 Fourier 平均中还使用 Tonelli--Fubini 定理作为测度先修: 对乘积 Lebesgue 测度, 非负可测函数可交换积分次序, 两边允许无穷大; 若 $\iint|h|<\infty$, 则复可测 $h$ 可交换次序, 各截面只需几乎处处可积. 不能只因每个截面看起来可积就交换. 后文每次交换均核对非负性或绝对可积性. 乘积测度的构造与此定理的完整证明留待实变科.
\end{remark}

\originalsection{范数、等价类与完备性}
\begin{definition}
对于 $1\le p<\infty$, $L^p(\mu)$ 由 $\int|f|^p<\infty$ 的可测函数按几乎处处相等取等价类组成, 范数为 $\norm{f}_p=(\int|f|^p)^{1/p}$. $L^\infty(\mu)$ 由本质有界函数的等价类组成, 范数为 $\norm{f}_\infty=\inf\{M\ge0:|f|\le M\text{ 几乎处处}\}$. 使用 $L^p(E)$ 时默认 Lebesgue 测度.
\end{definition}
该定义的运算与代表元无关, 因为两组代表元之差只在两个零测集的并上可能不为零. 若 $\int|f|^p=0$, 则每个集合 $\{|f|>1/n\}$ 测度为零, 因为积分至少是 $n^{-p}$ 倍该集合的测度; 它们的并覆盖 $f\ne0$ 的点. 所以范数为零恰好是零等价类. $L^\infty$ 的最小界也几乎处处成立, 可将界 $\norm{f}_\infty+1/n$ 的零测例外取可数并.

\begin{theorem}\label{thm:holder-lp}
若 $1/p+1/q=1$, $1\le p\le\infty$, 则 $\int|fg|\le\norm{f}_p\norm{g}_q$. 对 $1\le p\le\infty$, 有 $\norm{f+g}_p\le\norm{f}_p+\norm{g}_p$. 因而 $L^p$ 是赋范空间.
\end{theorem}
\begin{proof}
当 $1<p<\infty$ 且范数非零, 将 $f,g$ 归一化后, 对 Young 不等式逐点积分, 得 Hölder. 当 $p=1$ 或 $q=1$, 用本质上界直接估计. 零范数意味着函数几乎处处为零, 单独成立.

为证明 Minkowski, 先用 $|f+g|^p\le2^{p-1}(|f|^p+|g|^p)$ 知 $f+g\in L^p$. 这个逐点界来自实函数 $t^p$ 的凸性. 若 $s=\norm{f+g}_p>0$, 分别对 $|f|\,|f+g|^{p-1}$ 和 $|g|\,|f+g|^{p-1}$ 用 Hölder, 得 $s^p\le(\norm{f}_p+\norm{g}_p)s^{p-1}$, 除法即得结论. $s=0$ 时直接成立. 端点 $p=1,\infty$ 分别积分或取本质上界.
\end{proof}

\begin{theorem}[Riesz--Fischer]\label{thm:lpbanach}
$L^p(\mu)$ 对每个 $1\le p\le\infty$ 都是 Banach 空间.
\end{theorem}
\begin{proof}
先设 $p<\infty$. 根据绝对级数判别, 只需证明 $\sum_n\norm{f_n}_p<\infty$ 时函数级数收敛于 $L^p$ 中某元素. 选各项代表元, 令 $g_N=\sum_{n=1}^N|f_n|$, $g=\sum_{n\ge1}|f_n|$. 由 Minkowski 有 $\norm{g_N}_p\le C=\sum_n\norm{f_n}_p$. 因 $g_N^p\uparrow g^p$, 单调收敛给出 $\int g^p\le C^p$. 所以 $g$ 几乎处处有限, $\sum_nf_n$ 几乎处处绝对收敛, 定义其和为 $f$, 并在零测例外上取零. 它可测, 且 $|f|\le g$, 故属于 $L^p$.

同一论证用于尾部, 得
\[
\norm{f-\sum_{n=1}^N f_n}_p
\le\norm{\sum_{n>N}|f_n|}_p
\le\sum_{n>N}\norm{f_n}_p\to0.
\]
因此绝对级数收敛, 空间完备.

对 $p=\infty$, 若 $(f_n)$ 为 Cauchy 序列, 给定任意两个指标, 其函数差的本质界在一个零测集外成立. 把所有指标对的例外取可数并, 再并上每个 $f_n$ 的本质界例外, 得共同零测集 $N$. 在 $S\setminus N$ 上, $(f_n)$ 为一致 Cauchy 的有界函数列, 因而一致收敛于有界函数 $f$. 在 $N$ 上令 $f=0$. 点态极限的可测性和一致差估计给出 $f\in L^\infty$, $\norm{f_n-f}_\infty\to0$.
\end{proof}

有可数稠密子集的度量空间称可分.
\begin{proposition}\label{prop:lp-density}
若 $[a,b]$ 为有限区间且 $1\le p<\infty$, 则端点为零的连续函数在 $L^p([a,b])$ 中稠密. 因而 $L^p([a,b])$ 可分.
\end{proposition}
\begin{proof}
先截断可测 $f$ 的大小, 再将实、虚部的有限值域分割成小格, 得简单函数逼近. 截断误差由控制收敛作用于 $|f|^p\ind_{\{|f|>N\}}$ 而趋于零; 值域格化的误差至多是常数倍格长乘 $(b-a)^{1/p}$.

这里采用 Lebesgue 测度的正则性作为构造先修: 有限区间内每个可测集都可由有限个区间的并在对称差测度上任意逼近. 于是可将简单函数中每个 $\ind_E$ 换成有限区间并的示性函数, $L^p$ 误差为对称差测度的 $1/p$ 次方. 对简单函数的有限和, Minkowski 给出 $\norm{\sum_jc_j(\ind_{E_j}-\ind_{U_j})}_p\le\sum_j|c_j|\mu(E_j\mathbin{\triangle}U_j)^{1/p}$, 因而可给每一项分别分配误差. 再把每个区间端点附近的跳跃改成宽度 $\delta$ 的线性过渡, 并在 $a,b$ 附近使函数降到零. 改动只发生在总长度趋于零的有限个小区间, 函数大小受固定常数 $C$ 控制, 误差不超过 $C$ 乘改动区域总长度的 $1/p$ 次方, 所以 $L^p$ 误差趋于零. 这给出稠密性.

固定 $a,b$ 和两个零端点值, 在内部支撑端点、折线节点和值上再取有理数 (复值取 $\Q+\ii\Q$), 可以任意逼近这些有限折线函数. 全部这样编码的函数构成可数集, 故得到可数稠密集. $p=\infty$ 不在本命题范围内: 跳跃函数不能被连续函数按本质上确界任意逼近.
\end{proof}

\begin{exercise}\label{legacy:ch05:exercise1}
若 $\mu(S)<\infty$ 且 $1\le p<q\le\infty$, 证明 $\norm{f}_p\le\mu(S)^{1/p-1/q}\norm{f}_q$. 同时核对 $f(x)=x^{-\alpha}$ 在 $(0,1)$ 的 $L^p$ 条件.
\end{exercise}
\begin{solution}
当 $q<\infty$, 对 $|f|^p$ 与1用共轭指数 $q/p$ 和 $q/(q-p)$ 的 Hölder, 得 $\int|f|^p\le(\int|f|^q)^{p/q}\mu(S)^{1-p/q}$. 取 $p$ 次根即得结论. 当 $q=\infty$, 逐点用本质界积分即可. 零测度时所有 $L^p$ 等价类均为零, 另行理解该结论. 对 $\alpha>0$, $\int_0^1x^{-\alpha p}\dd x$ 有限当且仅当 $\alpha p<1$, 所以它精确展示包含关系的方向.
\end{solution}

\begin{exercise}\label{legacy:ch05:exercise2}
设 $f_n\to f$ 于 $L^p$, $1\le p<\infty$. 证明存在一个子列几乎处处趋于 $f$. 说明全列不必几乎处处收敛.
\end{exercise}
\begin{solution}
选 $n_k$ 使 $\norm{f_{n_k}-f}_p^p\le2^{-k}$. 单调收敛给出 $\int\sum_k|f_{n_k}-f|^p\le1$, 故几乎处处该标量级数有限, 其项趋于零, 得所需子列. 对全列, 依次列出 $[0,1)$ 的每层二进区间示性函数: 第 $m$ 层的每项范数为 $2^{-m/p}\to0$. 每个 $x\in[0,1)$ 在每层有一项值为1, 而其他大量项为0, 因而全列在每点都不趋于零.
\end{solution}

\originalchapter{Hilbert 空间与正交展开}\label{ch:hilbert}
\originalsection{内积与几何}
范数可度量长度, 却未必给出角度和正交分解. Hilbert 空间增加内积, 并保留完备性; 这两个条件共同支持无限维的投影与展开. 本书采用内积对第一变量线性、对第二变量共轭线性的约定.

\begin{definition}
内积是函数 $\ip{x}{y}:H\times H\to\K$, 对第一变量线性, 满足 $\ip{x}{y}=\overline{\ip{y}{x}}$, 且 $\ip{x}{x}\ge0$, 等号当且仅当 $x=0$. 带内积的空间称预 Hilbert 空间. 由内积诱导的范数为 $\norm{x}=\sqrt{\ip{x}{x}}$. 此范数下完备时称 Hilbert 空间.
\end{definition}

\begin{theorem}[Cauchy--Schwarz]\label{thm:cs}
预 Hilbert 空间中 $|\ip{x}{y}|\le\norm{x}\norm{y}$. 等号成立当且仅当 $x,y$ 线性相关.
\end{theorem}
\begin{proof}
若 $y=0$, 结论成立. 否则令 $a=\ip{x}{y}/\norm{y}^2$, 则 $\ip{x-ay}{y}=0$. 展开得到
\[
0\le\norm{x-ay}^2=\norm{x}^2-\frac{|\ip{x}{y}|^2}{\norm{y}^2}.
\]
移项后取平方根即得不等式. 等号恰好要求 $\norm{x-ay}=0$, 即 $x=ay$; 包含零向量的情况也符合线性相关条件.
\end{proof}
正定性和齐次性直接给出诱导范数的前两个条件, 而 $\norm{x+y}^2=\norm{x}^2+\norm{y}^2+2\operatorname{Re}\ip{x}{y}\le(\norm{x}+\norm{y})^2$ 给出三角不等式. 内积连续, 因为
$|\ip{x_n}{y_n}-\ip{x}{y}|\le\norm{x_n-x}\norm{y_n}+\norm{x}\norm{y_n-y}$, 当两向量列收敛时右端趋于零.

\begin{proposition}[平行四边形恒等式]\label{prop:parallelogram}
内积范数满足 $\norm{x+y}^2+\norm{x-y}^2=2\norm{x}^2+2\norm{y}^2$. 反之, 满足此恒等式的实或复范数来自唯一内积.
\end{proposition}
\begin{proof}
正向展开两个平方, 混合项抵消. 反向的细节为编者补全. 先在底层实空间令 $B(x,y)=(\norm{x+y}^2-\norm{x-y}^2)/4$. 它连续、对称, 且 $B(-x,y)=-B(x,y)$. 对平行四边形恒等式分别取向量 $x+y,z$ 与 $x-y,z$, 两式相减, 得
\[
B(x+z,y)+B(x-z,y)=2B(x,y).
\]
取 $z=x$ 得 $B(2x,y)=2B(x,y)$. 再令 $x=(u+v)/2$, $z=(u-v)/2$, 得 $B(u,y)+B(v,y)=B(u+v,y)$. 因而有整数与有理数齐次性; 连续性再将齐次性延至所有实数. $B(x,x)=\norm{x}^2$, 所以实情形已得到内积.

复情形由 $\norm{\ii x}=\norm{x}$ 得 $B(\ii x,\ii y)=B(x,y)$, 从而 $B(\ii x,y)=-B(x,\ii y)$. 定义 $\ip{x}{y}=B(x,y)+\ii B(x,\ii y)$. 实双线性及上述关系给出第一变量的复线性; 对称性和 $B(y,\ii x)=-B(x,\ii y)$ 给出共轭对称. 又有 $B(x,\ii x)=0$, 故 $\ip{x}{x}=\norm{x}^2$. 由这些恒等式倒推可见内积被范数唯一确定.
\end{proof}

\begin{example}
核对 $\ell^2$ 和 $L^2(E)$ 的内积, 并说明 $\ell^p$, $p\ne2$, 的标准范数不来自内积.
\end{example}
\begin{solution}
分别取 $\ip{a}{b}=\sum_ja_j\overline{b_j}$ 与 $\ip{f}{g}=\int_Ef\overline g$. Hölder 保证绝对收敛或绝对可积, 所以定义合法. 对角线值恰为对应范数的平方, 其余内积公理由求和与积分得到. 前文已证明完备性, 因而二者为 Hilbert 空间.

在 $\ell^p$ 中取 $e_1,e_2$, 当 $p<\infty$ 时平行四边形左端为 $2\cdot2^{2/p}$, 右端为4, 仅 $p=2$ 相等. $p=\infty$ 时左端为2, 也不等于4. 这否定的是标准范数的内积表示, 不否定空间可改用其他拓扑或其他结构.
\end{solution}

\originalsection{正交系与 Bessel 不等式}
\begin{definition}
$x\perp y$ 表示 $\ip{x}{y}=0$. 正交归一系 $(e_j)$ 满足 $\ip{e_j}{e_k}=\delta_{jk}$. 一个正交归一系的有限线性组合在 $H$ 中稠密时, 称其为正交归一基. 基可以不可数; 本书的展开主要在可分 Hilbert 空间中使用可数基.
\end{definition}
正交归一系中的向量两两距离为 $\sqrt2$, 因此可分空间中的这种系至多可数: 以各向量为中心、半径小于 $\sqrt2/2$ 的球两两不交, 每球需包含可数稠密集的一个不同元素.

\begin{theorem}[Bessel 不等式]\label{thm:bessel}
对任一有限或可数正交归一系 $(e_j)$, 有 $\sum_j|\ip{x}{e_j}|^2\le\norm{x}^2$. 若 $P_Nx=\sum_{j=1}^N\ip{x}{e_j}e_j$, 则
\[
\norm{x-P_Nx}^2=\norm{x}^2-\sum_{j=1}^N|\ip{x}{e_j}|^2.
\]
\end{theorem}
\begin{proof}
由正交归一性, $x-P_Nx$ 与每个 $e_j$, $j\le N$, 正交, 且 $\norm{P_Nx}^2=\sum_{j=1}^N|\ip{x}{e_j}|^2$. 对 $x=P_Nx+(x-P_Nx)$ 展开范数平方, 正交使混合项为零. 得有限恒等式, 非负余项给出有限不等式, 最后令 $N\to\infty$.
\end{proof}

\begin{theorem}\label{thm:onb}
每个可分 Hilbert 空间都有有限或可数正交归一基. 若 $(e_j)$ 为此基, 则对每个 $x\in H$ 有范数收敛及 Parseval 恒等式
\[
x=\sum_j\ip{x}{e_j}e_j,\qquad
\norm{x}^2=\sum_j|\ip{x}{e_j}|^2.
\]
对有限或可数正交归一系, 它是基当且仅当只与整个系都正交的向量是零. 此处级数证明限于有限或可数系; 不可数情形的判别可在下一章用闭子空间的正交补证明.
\end{theorem}
\begin{proof}
从可数稠密集 $(v_n)$ 出发做 Gram--Schmidt. 在已有 $e_1,\ldots,e_k$ 时, 令 $w=v_n-\sum_{j=1}^k\ip{v_n}{e_j}e_j$. 若 $w=0$, 跳过; 否则加入 $w/\norm{w}$. 正交性逐项核对. 每个 $v_n$ 都在所构系的有限张成中, 所以最终张成稠密. 零空间取空基. 反过来, 有有限或可数正交归一基时, 系数取有理数 (复值取 $\Q+\ii\Q$) 的有限组合构成可数稠密集, 所以空间可分.

对任一可数正交归一系, Bessel 不等式和正交性给出
$\norm{P_Mx-P_Nx}^2=\sum_{j=N+1}^M|\ip{x}{e_j}|^2\to0$.
完备性给出极限 $y$. 对每个固定 $e_j$ 取内积极限, 得 $\ip{x-y}{e_j}=0$. 若系的张成稠密, 连续性说明 $x-y$ 与所有向量正交, 因而 $x=y$. 有限恒等式取极限给出 Parseval.

反之, 若只有零向量与全系正交, 则上述极限仍满足 $x-y$ 与全系正交, 从而 $y=x$, 有限张成稠密. 这一论证使用 Hilbert 完备性, 不把未完备预 Hilbert 空间中的最大正交系自动当作收敛展开依据.
\end{proof}

\originalsection{坐标表示与分类}
\begin{theorem}\label{thm:hilbert-l2}
每个无限维可分 Hilbert 空间等距同构于 $\ell^2$. 对选定的基 $(e_j)$, 同构为 $Ux=(\ip{x}{e_j})_j$.
\end{theorem}
\begin{proof}
第一变量线性使 $U$ 线性, Parseval 给出 $\norm{Ux}_2=\norm{x}$, 因而单射. 给定 $a\in\ell^2$, 级数 $\sum_ja_je_j$ 的部分和由平方尾和估计为 Cauchy, 完备性给出和 $x$. 对每个 $e_j$ 取内积得 $\ip{x}{e_j}=a_j$, 所以 $U$ 满射. 对两向量的部分和取内积并取极限, 得 $\ip{x}{y}=\sum_j\ip{x}{e_j}\overline{\ip{y}{e_j}}=\ip{Ux}{Uy}$, 因此也保持内积. 有限维情形同理对应于 $\K^d$.
\end{proof}
抽象 Hilbert 空间的可分分类不表示所有问题都成为有限矩阵运算. 算子在所选基中可以耦合无穷多个坐标, 求逆与谱仍有不同于有限维的现象.

\begin{exercise}
对任一正交归一基, 证明 $P_Nx$ 是 $x$ 在 $\Span(e_1,\ldots,e_N)$ 中的唯一最佳近似.
\end{exercise}
\begin{solution}
任意 $v$ 在该有限张成中, $x-v=(x-P_Nx)+(P_Nx-v)$ 的两项正交. 故 $\norm{x-v}^2=\norm{x-P_Nx}^2+\norm{P_Nx-v}^2\ge\norm{x-P_Nx}^2$, 等号恰好为 $v=P_Nx$.
\end{solution}

\begin{exercise}
证明 $\ell^2$ 中的向量级数 $\sum_{n\ge1}e_n/n$ 收敛但不绝对收敛. 说明 Hilbert 展开中的平方可和不能换成绝对可和.
\end{exercise}
\begin{solution}
部分和之差的范数平方为对应的 $\sum n^{-2}$ 尾和, 因而 Cauchy 并收敛. 但各项范数之和是 $\sum1/n$, 发散. 正交性消除了混合项, 所以这里只需平方可和. 绝对级数判别是足够条件, 并非每个收敛级数都必须满足的条件.
\end{solution}

\originalchapter{投影、表示与伴随}\label{ch:adjoint}
\originalsection{闭凸集的最短距离}
有限维闭有界集紧, 因而连续函数常能取到最值. 无限维的闭球一般不紧, 最短距离问题需要其他结构. Hilbert 空间的平行四边形恒等式与凸性, 能直接把极小化序列变成 Cauchy 序列.

\begin{theorem}[投影定理]\label{thm:projection}
设 $C$ 为 Hilbert 空间 $H$ 的非空闭凸集. 对每个 $x\in H$, 存在唯一 $p\in C$ 使 $\norm{x-p}=\dist(x,C)$. 此点等价地由
\[
\operatorname{Re}\ip{x-p}{z-p}\le0\quad(z\in C)
\]
刻画.
\end{theorem}
\begin{proof}
平移后可设 $x=0$. 令 $d=\inf_{z\in C}\norm{z}$, 选 $z_n\in C$ 使 $\norm{z_n}\to d$. 因中点仍在 $C$, 平行四边形恒等式给出
\[
\norm{z_n-z_m}^2
=2\norm{z_n}^2+2\norm{z_m}^2-4\norm{(z_n+z_m)/2}^2
\le2\norm{z_n}^2+2\norm{z_m}^2-4d^2\to0.
\]
所以 $(z_n)$ 为 Cauchy, 完备性给出极限 $p$, 闭性使 $p\in C$, 范数连续性使 $\norm{p}=d$. 若 $p,q$ 都取到最小范数, 同一不等式给出 $\norm{p-q}^2\le0$, 所以唯一.

回到一般 $x$, 若 $p$ 最优, 对 $z\in C$ 与 $0<t\le1$, 点 $p+t(z-p)$ 在 $C$, 所以
$0\le\norm{x-p-t(z-p)}^2-\norm{x-p}^2=-2t\operatorname{Re}\ip{x-p}{z-p}+t^2\norm{z-p}^2$.
除以 $t$ 后令 $t\downarrow0$, 得刻画式. 反之, 刻画式代入 $\norm{x-z}^2=\norm{x-p}^2+\norm{z-p}^2-2\operatorname{Re}\ip{x-p}{z-p}$, 得 $p$ 最优, 且 $z\ne p$ 时严格增大.
\end{proof}

\begin{corollary}\label{cor:orthodecompose}
若 $M$ 为 $H$ 的闭线性子空间, 则 $H=M\oplus M^\perp$, 其中 $M^\perp=\{y:\ip{y}{m}=0\text{ 对所有 }m\in M\}$. 分解 $x=P_Mx+(I-P_M)x$ 唯一; $P_M$ 线性, $P_M^2=P_M$, 且 $\norm{P_M}\le1$.
\end{corollary}
\begin{proof}
$M^\perp$ 是线性子空间, 并是连续泛函 $y\mapsto\ip{y}{m}$ 的零集之交, 因而闭. 取 $p$ 为 $x$ 在 $M$ 上的最优点. 投影刻画中可代入 $z=p\pm m$, 得 $\operatorname{Re}\ip{x-p}{m}=0$. 复情形再代入 $p\pm\ii m$, 得虚部也为零. 所以 $x-p\in M^\perp$. $M\cap M^\perp=\{0\}$, 因为同一个向量与自身正交则为零, 因而分解唯一.

两向量分解的线性组合仍为合法分解, 唯一性给出 $P_M$ 线性. 对 $m\in M$ 有 $P_Mm=m$, 故幂等. 正交性给出 $\norm{x}^2=\norm{P_Mx}^2+\norm{(I-P_M)x}^2$, 得范数界. 若 $M\ne\{0\}$, 在单位 $m\in M$ 上达到1.
\end{proof}
对不闭的子空间仍有 $M^\perp=(\overline M)^\perp$, 因内积连续. 用闭子空间分解可得 $(M^\perp)^\perp=\overline M$. 具体地, 将 $x$ 分解为 $m+n$ 于 $\overline M\oplus M^\perp$, 若 $x$ 与 $M^\perp$ 正交, 则 $0=\ip{x}{n}=\norm{n}^2$, 故 $x\in\overline M$.

\begin{center}
\begin{tikzpicture}[scale=1,>=Latex]
\draw[->](-.4,0)--(5,0) node[right]{$M$};
\draw[->](0,-.3)--(0,2.8) node[above]{$M^\perp$};
\draw[->,thick,structurecolor](0,0)--(3.6,2.1) node[above right]{$x$};
\draw[->,thick](0,0)--(3.6,0) node[below]{$P_Mx$};
\draw[dashed](3.6,0)--(3.6,2.1);
\draw[->,thick](0,0)--(0,2.1) node[left]{$x-P_Mx$};
\draw[dashed](0,2.1)--(3.6,2.1);
\draw(3.35,0)--(3.35,.25)--(3.6,.25);
\end{tikzpicture}
\end{center}
图中二维分解只表示正交关系. 一般子空间可以无限维, 投影定理不依赖把整个空间画成坐标平面.

\originalsection{Riesz 表示}
\begin{theorem}[Riesz 表示定理]\label{thm:riesz}
对 Hilbert 空间 $H$ 的每个 $f\in H^*$, 存在唯一 $y\in H$ 使 $f(x)=\ip{x}{y}$ 对所有 $x$ 成立, 且 $\norm{f}=\norm{y}$. 在复情形, $y\mapsto f_y$ 是共轭线性等距双射.
\end{theorem}
\begin{proof}
若 $f=0$, 取 $y=0$. 否则 $M=\ker f$ 为闭子空间, 取 $x_0$ 使 $f(x_0)=1$. 分解 $x_0=m+z$, 其中 $m\in M$, $z\in M^\perp$. 则 $f(z)=1$, 所以 $z\ne0$. 任意 $x$ 有 $x-f(x)z\in M$, 因而
$\ip{x}{z}=\ip{x-f(x)z}{z}+f(x)\norm{z}^2=f(x)\norm{z}^2$.
取 $y=z/\norm{z}^2$, 得表示. 若两个向量表示同一泛函, 它们之差与所有 $x$ 的内积为零; 取 $x$ 为此差, 得二者相等. Cauchy--Schwarz 给出范数上界, 在 $x=y/\norm{y}$ 上达到下界, 零 $y$ 单独成立. 第二变量共轭线性使 $f_{\alpha y}=\overline\alpha f_y$, 所以复情形不可把此对应写成普通线性映射.
\end{proof}

\begin{proposition}\label{prop:hilbert-reflexive}
每个 Hilbert 空间自反.
\end{proposition}
\begin{proof}
取 $\Phi\in H^{**}$. 复情形令 $g(y)=\overline{\Phi(f_y)}$, 实情形无须加共轭. 前述共轭线性对应使 $g$ 对 $y$ 线性, 且 $|g(y)|\le\norm{\Phi}\norm{y}$. 用 Riesz 表示得 $g(y)=\ip{y}{z}$ 对某个 $z$ 成立. 因而 $\Phi(f_y)=\ip{z}{y}=f_y(z)$. 所有 $f\in H^*$ 都有形式 $f_y$, 所以 $\Phi=J_Hz$, 典范映射满射.
\end{proof}

\originalsection{伴随与方程的可解性}
\begin{theorem}\label{thm:adjoint}
设 $H_1,H_2$ 为 Hilbert 空间, $T\in\B(H_1,H_2)$. 存在唯一 $T^*\in\B(H_2,H_1)$ 满足 $\ip{Tx}{y}=\ip{x}{T^*y}$, 且 $\norm{T^*}=\norm{T}$.
\end{theorem}
\begin{proof}
固定 $y\in H_2$, 函数 $x\mapsto\ip{Tx}{y}$ 线性且绝对值不超过 $\norm{T}\norm{x}\norm{y}$. Riesz 表示给出唯一向量 $T^*y$, 且 $\norm{T^*y}\le\norm{T}\norm{y}$. 比较 $\ip{x}{T^*(\alpha y+\beta z)}$ 与 $\ip{x}{\alpha T^*y+\beta T^*z}$, 用第二变量的共轭线性及表示的唯一性, 得 $T^*$ 线性. 所以已证明 $\norm{T^*}\le\norm{T}$.

共轭内积等式给出 $(T^*)^*=T$, 再用同一范数上界得反向比较. 唯一性来自若所有 $\ip{x}{(S-T^*)y}=0$, 则取 $x=(S-T^*)y$ 得 $Sy=T^*y$.
\end{proof}
从定义逐项核对可得 $(S+T)^*=S^*+T^*$, $(\alpha T)^*=\overline\alpha T^*$, $(ST)^*=T^*S^*$, $\norm{T^*T}=\norm{T}^2$. 最后一式的上界来自复合估计, 下界来自 $\norm{Tx}^2=\ip{T^*Tx}{x}\le\norm{T^*T}\norm{x}^2$. 算子称自伴, 指 $T=T^*$; 称酉, 指 $T^*T=TT^*=I$. 酉算子保持内积且满射, 反之满射的线性内积等距映射也是酉算子.

\begin{theorem}\label{thm:rangeperp}
有界算子满足 $(\ran T)^\perp=\ker T^*$, 从而 $\overline{\ran T}=(\ker T^*)^\perp$. 若值域闭, 则 $Tx=y$ 可解当且仅当 $y\perp\ker T^*$; 所有解之差属于 $\ker T$.
\end{theorem}
\begin{proof}
$y\in(\ran T)^\perp$ 当且仅当对每个 $x$ 有 $\ip{Tx}{y}=0$, 由伴随等式等价于对每个 $x$ 有 $\ip{x}{T^*y}=0$, 即 $T^*y=0$. 对等式再取正交补, 并使用双正交补等于闭包, 得第二式. 值域闭时闭包可删去, 正好给出方程的相容条件. 若 $Tx_1=Tx_2$, 则 $x_1-x_2\in\ker T$, 反之给任意解加核向量仍是解.
\end{proof}

\begin{example}\label{legacy:ch07:example1}
在 $\ell^2$ 上令右移 $S(a_1,a_2,\ldots)=(0,a_1,a_2,\ldots)$. 求伴随, 并判断它是否酉.
\end{example}
\begin{solution}
直接配对得 $\ip{Sa}{b}=\sum_{j\ge1}a_j\overline{b_{j+1}}$, 所以 $S^*b=(b_2,b_3,\ldots)$. 有 $S^*S=I$, 但 $SS^*b=b-b_1e_1$, 因而 $S$ 是等距嵌入却非满射, 不是酉算子. $\ker S^*=\K e_1$, 其正交补为首坐标为零的数列, 正好是 $S$ 的闭值域.
\end{solution}

\begin{exercise}\label{legacy:ch07:exercise1}
设 $P$ 为有界幂等算子. 证明它是某闭子空间上的正交投影, 当且仅当 $P=P^*$.
\end{exercise}
\begin{solution}
若 $P$ 为正交投影, 用两向量的正交分解得 $\ip{Px}{y}=\ip{Px}{Py}=\ip{x}{Py}$, 所以自伴. 反之, 幂等给出 $\ran P=\ker(I-P)$ 闭, 及 $x=Px+(I-P)x$. 若自伴, $\ip{Px}{(I-P)y}=\ip{x}{P(I-P)y}=0$, 所以这就是正交分解, $P$ 为对应正交投影.
\end{solution}

\begin{exercise}\label{legacy:ch07:exercise2}
连续核 $k$ 在 $L^2([0,1])$ 上给出积分算子 $Tf(x)=\int_0^1k(x,y)f(y)\dd y$. 求 $T^*$.
\end{exercise}
\begin{solution}
Cauchy--Schwarz 给出 $\norm{Tf}_2\le\norm{k}_{L^2([0,1]^2)}\norm{f}_2$. 对 $f,g\in L^2$, 双积分绝对值受 $\norm{k}_\infty\norm{f}_1\norm{g}_1$ 控制, 且有限区间的 Hölder 保证这两个 $L^1$ 范数有限. 因此 Fubini 合法, 得
\[
\ip{Tf}{g}=\int_0^1 f(y)\overline{\left(\int_0^1\overline{k(x,y)}g(x)\dd x\right)}\dd y.
\]
所以 $T^*g(y)=\int_0^1\overline{k(x,y)}g(x)\dd x$. 若 $k(x,y)=\overline{k(y,x)}$, 算子自伴.
\end{solution}

\originalchapter{弱拓扑与弱收敛}\label{ch:weak}
\originalsection{弱与弱星拓扑}
本章将主源的双对偶内容与 Melrose 讲义第3章第15节的弱收敛串联. 一般赋范空间的弱、弱星拓扑定义及可分对偶球的论证为\textbf{编者补充}. 弱收敛减少了需要一致控制的观测数量, 为有界序列提供较容易取得的极限.

\begin{definition}
$X$ 上的弱拓扑 $\sigma(X,X^*)$ 是使每个 $f\in X^*$ 连续的最弱拓扑. 在 $x$ 处的基本邻域形如
\[
\{y\in X:|f_j(y-x)|<\eps,\ 1\le j\le m\},
\]
其中只取有限个泛函. $X^*$ 上的弱星拓扑 $\sigma(X^*,X)$ 则用有限个 $x_j\in X$ 的取值 $|f(x_j)-g(x_j)|<\eps$ 定义邻域. 序列 $x_n\rightharpoonup x$ 指对每个 $f\in X^*$ 有 $f(x_n)\to f(x)$; 序列 $f_n\overset{*}{\rightharpoonup}f$ 指对每个 $x\in X$ 有 $f_n(x)\to f(x)$.
\end{definition}
有限交形式确实组成邻域基, 因为交集可合并有限个约束, 缩小阈值仍给邻域. 任一使所有泛函连续的拓扑必须包含这些邻域, 因而它是最弱的. 单个序列满足所有取值收敛, 就最终满足任意有限组约束, 所以与拓扑的序列收敛相符. 这里并不声称一般弱拓扑的闭性都能只靠序列刻画.

Hahn--Banach 说明弱拓扑是 Hausdorff: 若 $x\ne y$, 可选 $f(x)\ne f(y)$, 用分开的标量小球取逆像就能分离二者. 对弱星拓扑, 不同泛函本来就在某个向量上取值不同. 因而弱极限与弱星极限唯一. 范数收敛蕴含弱收敛, 因 $|f(x_n-x)|\le\norm{f}\norm{x_n-x}$. 对偶的范数收敛同样蕴含弱星收敛.

\begin{proposition}\label{prop:weakbounded}
任何赋范空间中的弱收敛序列在范数下有界. 若 $X$ 为 Banach 空间, $X^*$ 中弱星收敛的序列也在范数下有界.
\end{proposition}
\begin{proof}
若 $x_n\rightharpoonup x$, 对每个 $f\in X^*$, 序列 $(J_Xx_n)(f)=f(x_n)$ 有界. 对 Banach 空间 $X^*$ 上的算子 $J_Xx_n:X^*\to\K$ 用统一有界原理, 得 $\sup_n\norm{J_Xx_n}<\infty$. 等距性给出 $\sup_n\norm{x_n}<\infty$, 不要求 $X$ 完备. 对弱星收敛的 $f_n$, 直接在定义域 Banach 空间 $X$ 上用统一有界原理, 得第二结论.
\end{proof}

\begin{proposition}\label{prop:weaklsc}
若 $x_n\rightharpoonup x$, 则 $\norm{x}\le\liminf_n\norm{x_n}$. 有界线性算子保持弱收敛.
\end{proposition}
\begin{proof}
若 $x\ne0$, 选单位泛函 $f$ 使 $f(x)=\norm{x}$. 因 $f(x_n)\to\norm{x}$ 且 $|f(x_n)|\le\norm{x_n}$, 取下极限即得结论. $x=0$ 时直接成立. 对 $T\in\B(X,Y)$ 和任意 $g\in Y^*$, 复合 $gT\in X^*$, 所以 $g(Tx_n)\to g(Tx)$.
\end{proof}

\begin{example}\label{legacy:ch08:example1}
比较 $\ell^2$ 中的 $e_n$ 及 $\ell^\infty$ 中的 $u^{(n)}=(\underbrace{0,\ldots,0}_{n\text{项}},1,1,\ldots)$ 的收敛.
\end{example}
\begin{solution}
Riesz 表示说明 $\ell^2$ 中弱收敛等价于对每个 $a\in\ell^2$ 的内积收敛. 因 $a_n\to0$, 有 $\ip{e_n}{a}=\overline{a_n}\to0$, 所以 $e_n\rightharpoonup0$, 但范数恒为1, 不强收敛.

把 $\ell^\infty=(\ell^1)^*$ 视为对偶, 对任意 $a\in\ell^1$ 有 $\sum_{j>n}a_j\to0$, 所以 $u^{(n)}$ 弱星趋于零. 但第\ref{ch:dual}章构造的泛函 $F\in(\ell^\infty)^*$ 满足 $F(u^{(n)})=1$, 因各 $u^{(n)}$ 是极限为1的收敛数列. 所以它不弱趋于零. 弱星测试向量来自 $\ell^1$, 弱测试泛函来自更大的 $(\ell^\infty)^*$, 两者不可混淆.
\end{solution}

\originalsection{Hilbert 空间的弱极限}
\begin{theorem}\label{thm:weak-coordinates}
设 $(e_j)$ 为可数正交归一基. 范数有界的序列 $(x_n)$ 弱收敛, 当且仅当每个坐标 $\ip{x_n}{e_j}$ 收敛. 若坐标极限为 $a_j$, 则弱极限为 $x=\sum_ja_je_j$.
\end{theorem}
\begin{proof}
正向由测试泛函 $y\mapsto\ip{y}{e_j}$ 得到. 反向设 $\norm{x_n}\le M$. 对每个有限 $J$, 坐标收敛给出 $\sum_{j=1}^J|a_j|^2=\lim_n\sum_{j=1}^J|\ip{x_n}{e_j}|^2\le M^2$. 令 $J\to\infty$, 得 $a\in\ell^2$, 从而向量 $x$ 存在且 $\norm{x}\le M$.

给定 $y\in H$, 令 $y_J=P_Jy$. 有
\[
|\ip{x_n-x}{y}|\le |\ip{x_n-x}{y_J}|+2M\norm{y-y_J}.
\]
先选 $J$ 使第二项小, 这一步不依赖 $n$; 再由有限个坐标收敛使第一项小. 所以对所有 $y$ 内积收敛. Riesz 表示把这等价为弱收敛. $M=0$ 时全序列零, 另行成立.
\end{proof}

\begin{theorem}\label{thm:weaksubsequence}
Hilbert 空间中的每个范数有界序列都有弱收敛子列. 这里说的是序列结论.
\end{theorem}
\begin{proof}
先设空间可分. 每个坐标序列有界, 由标量 Bolzano--Weierstrass 定理, 对第一个坐标取收敛子列, 再对第二个取前一子列的收敛子列, 依次进行. 取第 $k$ 个嵌套子列中的第 $k$ 项, 并保持原指标严格递增, 得对角子列; 对任意固定坐标, 该子列从某项起属于相应的收敛子列, 因而坐标收敛. 用定理\ref{thm:weak-coordinates}得到弱极限.

若 $H$ 不可分, 令 $M$ 为原序列的闭线性张成. 有理实或有理复系数的有限线性组合给出可数稠密集, 所以 $M$ 可分. 在 $M$ 中用上述论证得到子列和极限. 对任意 $y\in H$, 分解 $y=P_My+y^\perp$, 序列及其极限均与 $y^\perp$ 正交, 所以内积测试化为 $M$ 中的测试. 因而该子列也在 $H$ 中弱收敛.
\end{proof}

\begin{proposition}\label{prop:weakstrong}
在 Hilbert 空间中, 若 $x_n\rightharpoonup x$ 且 $\norm{x_n}\to\norm{x}$, 则 $x_n\to x$ 于范数. 非空闭凸集对弱收敛序列的极限封闭.
\end{proposition}
\begin{proof}
展开得 $\norm{x_n-x}^2=\norm{x_n}^2+\norm{x}^2-2\operatorname{Re}\ip{x_n}{x}\to0$. 对闭凸集 $C$, 若 $x\notin C$, 取 $p=P_Cx$. 投影刻画给出对所有 $z\in C$ 有 $\operatorname{Re}\ip{z-p}{x-p}\le0$. 若 $x_n\in C$ 弱趋于 $x$, 则对此固定向量取内积极限, 得 $\norm{x-p}^2\le0$, 矛盾.
\end{proof}

\originalsection{可分情形的弱星紧性}
\begin{theorem}\label{thm:weakstar-separable}
若赋范空间 $X$ 可分, 则对偶闭单位球的每个序列都有弱星收敛子列, 极限仍在闭单位球. 该球上的弱星拓扑可由一个度量给出, 因而该球在弱星拓扑中紧.
\end{theorem}
\begin{proof}
取 $X$ 的可数稠密集 $(x_j)$. 对 $\norm{f_n}\le1$, 每个标量序列 $f_n(x_j)$ 有界, 对角选择得子列 $(g_n)$, 使所有这些取值收敛. 对任意 $x$, 选 $x_j$ 接近 $x$, 则
$|g_n(x)-g_m(x)|\le2\norm{x-x_j}+|g_n(x_j)-g_m(x_j)|$.
先选 $j$, 再取大 $n,m$, 可知 $(g_n(x))$ 为 Cauchy. 定义 $f(x)=\lim_ng_n(x)$, 有限线性运算取极限给出线性性, 且 $|f(x)|\le\norm{x}$, 所以 $f\in X^*$, $\norm{f}\le1$, 子列弱星收敛.

在此单位球上定义
\[
d(f,g)=\sum_{j=1}^\infty2^{-j}\min\{1,|f(x_j)-g(x_j)|\}.
\]
截断函数满足三角不等式, 总和有限. 若距离为零, 两泛函在稠密集上一致, 连续性使它们处处一致, 所以 $d$ 为度量. 对有限个 $x_j$ 的取值约束能使有限部分小, 剩余尾和统一小, 说明每个度量球包含弱星邻域. 反方向, 对给定的有限测试向量 $y_1,\ldots,y_m$, 各选稠密点 $x_{j_k}$ 使 $2\norm{y_k-x_{j_k}}$ 很小; 距离足够小就使每个 $|f(x_{j_k})-g(x_{j_k})|$ 小, 再用上述估计控制 $|f(y_k)-g(y_k)|$. 所以两拓扑相同.

已经证明每个序列有收敛子列. 在度量空间, 序列紧等价于紧, 因而得到最后结论. 此等价是本书采用的度量空间先修, 可参见已归档的基础拓扑册. 一般不可分空间的 Banach--Alaoglu 定理需要乘积紧性, 不在本册证明范围内.
\end{proof}

\begin{exercise}\label{legacy:ch08:exercise1}
设 $X$ 为 Banach 空间, $f_n\in X^*$, 且每个 $f_n(x)$ 都有极限. 证明这些极限构成一个 $f\in X^*$, 且 $f_n$ 弱星趋于 $f$.
\end{exercise}
\begin{solution}
逐点收敛给出逐点有界, 统一有界原理给出 $M=\sup_n\norm{f_n}<\infty$. 令 $f(x)=\lim_nf_n(x)$, 有限线性运算取极限说明线性, 且 $|f(x)|\le M\norm{x}$, 因而连续. 弱星收敛正是其定义. 不需声称范数收敛.
\end{solution}

\originalchapter{Fourier 级数}\label{ch:fourier}
\originalsection{Fourier 系数与 Dirichlet 核}
正交归一系给出投影, 但只有完整的正交归一基才保证投影趋于原向量. Fourier 级数的核心问题, 正是指数函数是否张成整个 $L^2$ 空间. 本章保留主源第15--16讲用 Fejér 平均证明完备性的路线.

在 $H=L^2([-\pi,\pi])$ 中取 $e_n(x)=(2\pi)^{-1/2}\ee^{\ii nx}$, $n\in\Z$. 因 $\int_{-\pi}^\pi\ee^{\ii(n-m)x}\dd x$ 在 $n=m$ 时为 $2\pi$, 否则为零, 这些函数组成正交归一系. 对 $f\in L^2$, 有 $f\in L^1$, 所以下述系数都有意义.

\begin{definition}
Fourier 系数为 $\widehat f(n)=(2\pi)^{-1}\int_{-\pi}^\pi f(t)\ee^{-\ii nt}\dd t$. 对称部分和为 $S_Nf=\sum_{|n|\le N}\widehat f(n)\ee^{\ii nx}$. 其形式级数称 Fourier 级数. Cesàro 平均为 $\sigma_Nf=(N+1)^{-1}\sum_{k=0}^NS_kf$.
\end{definition}
有限部分和与有限平均是定义明确的三角多项式. 在证明基的完备性之前, 不把形式级数当作已经等于 $f$ 的表达式.

\begin{proposition}
令 $D_N(t)=(2\pi)^{-1}\sum_{|n|\le N}\ee^{\ii nt}$. 则 $S_Nf(x)=\int_{-\pi}^\pi D_N(x-t)f(t)\dd t$, 且
\[
D_N(t)=\frac{\sin((N+1/2)t)}{2\pi\sin(t/2)}\quad(t\notin2\pi\Z),
\qquad D_N(0)=\frac{2N+1}{2\pi}.
\]
\end{proposition}
\begin{proof}
把 Fourier 系数代入有限和并交换有限求和与积分即可得到核表示. 对 $\ee^{\ii t}\ne1$, 几何级数给出
\[
\sum_{n=-N}^N\ee^{\ii nt}
=\ee^{-\ii Nt}\frac{1-\ee^{\ii(2N+1)t}}{1-\ee^{\ii t}}
=\frac{\sin((N+1/2)t)}{\sin(t/2)}.
\]
在 $t\in2\pi\Z$ 时直接代入有限和得 $2N+1$, 也是连续延拓值. 核对所有实数周期延拓, 不把商式的零分母当真正奇点.
\end{proof}
Dirichlet 核会改变符号, 所以其积分为1并不足以控制卷积误差的绝对值. 平均后的核非负, 才能用其总质量估计误差.

\originalsection{Fejér 核与一致收敛}
\begin{proposition}\label{prop:fejerkernel}
Cesàro 平均可写成 $\sigma_Nf(x)=\int_{-\pi}^\pi F_N(x-t)f(t)\dd t$, 其中
\[
F_N(t)=\frac1{2\pi(N+1)}\left|\sum_{k=0}^N\ee^{\ii kt}\right|^2
=\frac1{2\pi(N+1)}\frac{\sin^2((N+1)t/2)}{\sin^2(t/2)}.
\]
第二式在 $2\pi\Z$ 上按连续延拓解释, $F_N(0)=(N+1)/(2\pi)$. 核非负、偶、$2\pi$ 周期, $\int_{-\pi}^\pi F_N=1$, 且对于 $0<\delta<\pi$,
\[
\sup_{\delta\le|t|\le\pi}F_N(t)
\le\frac1{2\pi(N+1)\sin^2(\delta/2)}.
\]
\end{proposition}
\begin{proof}
平均各 $D_k$ 时, 某个频率 $n$, $|n|\le N$, 出现 $N+1-|n|$ 次. 因而 $F_N(t)=(2\pi)^{-1}\sum_{|n|\le N}(1-|n|/(N+1))\ee^{\ii nt}$. 展开 $|\sum_{k=0}^N\ee^{\ii kt}|^2$, 频率 $n=k-l$ 也出现相同次数, 得第一式. 几何级数的模给出正弦商式. 非负性由平方模得到, 偶性由取共轭得到, 周期性由有限指数和得到. 积分只保留零频率项, 所以一个周期的积分为1. 最后的界使用分子不超过1及 $\sin^2(t/2)\ge\sin^2(\delta/2)$.
\end{proof}

\begin{center}
\begin{tikzpicture}
\begin{axis}[width=11cm,height=4.8cm,axis lines=left,clip=false,
xmin=-3.1416,xmax=3.1416,ymin=0,ymax=1.6,
xtick={-3.14159,0,3.14159},xticklabels={$-\pi$,$0$,$\pi$},
ytick={0,.5,1,1.5},xlabel={$t$},
xlabel style={at={(axis description cs:1.03,0)},anchor=west},
tick align=outside]
\addplot[thick,structurecolor,domain=-pi:pi,samples=401]
{abs(x)<0.00001 ? 9/(2*pi) : (sin(deg(9*x/2))/sin(deg(x/2)))^2/(18*pi)};
\node[anchor=south west] at (axis description cs:0,1.05) {$F_8(t)$};
\end{axis}
\end{tikzpicture}
\end{center}
图采用弧度变量 $t\in[-\pi,\pi]$, 显式将 PGFPlots 的三角函数参数转成度, 中心附近用连续极限值作数值保护. 它仅展示 $N=8$ 的核; 质量和收敛结论由公式证明, 不由图形采样推断.

\begin{theorem}[Fejér 定理]\label{thm:fejer}
若 $f$ 连续且 $2\pi$ 周期, 则 $\sigma_Nf\to f$ 一致收敛.
\end{theorem}
\begin{proof}
周期连续函数在整个实轴上一致连续且有界. 由周期换元, 可把核表示写成 $\sigma_Nf(x)=\int_{-\pi}^\pi F_N(t)f(x-t)\dd t$. 利用核总质量为1,
\[
|\sigma_Nf(x)-f(x)|
\le\int_{-\pi}^\pi F_N(t)|f(x-t)-f(x)|\dd t.
\]
给定 $\eps>0$, 先选 $0<\delta<\pi$ 使 $|t|<\delta$ 时差值小于 $\eps/2$, 对所有 $x$ 同时成立. 在此小区间上的积分至多 $\eps/2$. 其余区间上差值不超过 $2\norm{f}_\infty$, 由核的远端界和区间长度至多 $2\pi$, 所以积分至多
$2\norm{f}_\infty/((N+1)\sin^2(\delta/2))$. 对充分大 $N$, 此数小于 $\eps/2$. 所以误差对所有 $x$ 同时小于 $\eps$, 得一致收敛.
\end{proof}

\originalsection{平方积分收敛与完整性}
\begin{lemma}\label{lem:fejercontraction}
对每个 $f\in L^2([-\pi,\pi])$, 有 $\norm{\sigma_Nf}_2\le\norm{f}_2$.
\end{lemma}
\begin{proof}
将 $f$ 几乎处处按周期延拓. 用非负密度 $F_N(t)\dd t$ 的 Cauchy--Schwarz 不等式及总质量1, 得
\[
|\sigma_Nf(x)|^2\le\int_{-\pi}^\pi F_N(t)|f(x-t)|^2\dd t.
\]
对 $x$ 积分, 被积函数非负, 可用 Tonelli 交换次序. 周期性使任意平移的一个周期平方积分相同, 所以右端双积分为 $\norm{f}_2^2\int F_N=\norm{f}_2^2$. 这也保证截面表达式几乎处处有限. 取平方根即得界.
\end{proof}

\begin{theorem}\label{thm:fouriercomplete}
对每个 $f\in L^2([-\pi,\pi])$, 有 $\sigma_Nf\to f$ 于 $L^2$. 指数系 $(e_n)_{n\in\Z}$ 为正交归一基, 所以也有 $S_Nf\to f$ 于 $L^2$, 及
\[
\norm{f}_2^2=2\pi\sum_{n\in\Z}|\widehat f(n)|^2.
\]
\end{theorem}
\begin{proof}
由连续端点为零函数的稠密性, 给定 $\eps>0$ 可选连续周期 $g$, 使 $\norm{f-g}_2<\eps/3$. Fejér 定理保证 $\norm{\sigma_Ng-g}_2\le\sqrt{2\pi}\norm{\sigma_Ng-g}_\infty<\eps/3$ 对大 $N$ 成立. 利用收缩界,
$\norm{\sigma_Nf-f}_2\le\norm{\sigma_N(f-g)}_2+\norm{\sigma_Ng-g}_2+\norm{g-f}_2<\eps$.

若 $f$ 与所有 $e_n$ 正交, 则所有 Fourier 系数为零, 每个平均也为零, 极限结论给出 $f=0$ 于 $L^2$. 定理\ref{thm:onb}于是保证指数系完整, 普通正交投影 $S_N$ 也收敛. Parseval 乘上系数归一化常数即得最后恒等式.
\end{proof}
这里的等式和收敛均针对几乎处处等价类及 $L^2$ 范数. 它不自动给每个点的普通 Fourier 部分和收敛. Carleson 定理等更深的调和分析结论不在本册证明范围内.

\begin{exercise}
求 $f(x)=x$, $-\pi<x<\pi$, 的 Fourier 系数, 并用 Parseval 求 $\sum_{n\ge1}n^{-2}$.
\end{exercise}
\begin{solution}
零频率为零. 对 $n\ne0$ 分部积分,
$\int_{-\pi}^\pi x\ee^{-\ii nx}\dd x=-2\pi(-1)^n/(\ii n)$, 剩余指数积分为零. 所以 $\widehat f(n)=\ii(-1)^n/n$. 又有 $\norm{f}_2^2=2\pi^3/3$. Parseval 给出 $2\pi^3/3=2\pi\sum_{n\ne0}n^{-2}=4\pi\sum_{n\ge1}n^{-2}$, 因而和为 $\pi^2/6$. 本推导只使用 $L^2$ 展开, 不需要在端点代入一个未证明逐点收敛的级数.
\end{solution}

\begin{exercise}
证明 $\sqrt2\sin(k\pi x)$, $k\ge1$, 在 $L^2([0,1])$ 中构成正交归一基.
\end{exercise}
\begin{solution}
正交归一性由积化和差积分得到. 若 $f$ 与所有正弦正交, 将 $f$ 奇延拓到 $[-1,1]$. 它与常数及所有余弦自动正交, 因为奇偶性使积分为零, 与正弦的配对则为原积分的两倍. 变量变换 $t=\pi x$ 把这些指数配对化为 $L^2([-\pi,\pi])$ 的 Fourier 系数, 它们全为零, 所以奇延拓为零, 从而 $f=0$. 完整性判别给出所需基.
\end{solution}

\originalchapter{紧集与紧算子}\label{ch:compact}
\originalsection{紧性与有限维逼近}
在 Hilbert 空间的单位球内, 正交归一序列没有范数收敛子列, 因为两项距离恒为 $\sqrt2$. 因而闭、有界不再保证紧. 缺少的条件是能用同一个有限维子空间同时逼近集合中的所有向量.

本章在度量空间中采用等价的序列紧定义. 全有界指对每个 $\eps>0$ 都能用有限个半径 $\eps$ 的球覆盖. 完备且全有界的度量空间紧: 对任意序列, 逐级取有限覆盖中含无穷项的球, 可抽出 Cauchy 子列; 完备性给出极限. 反之紧集若没有某个有限球覆盖, 可逐步选取两两距离至少 $\eps$ 的点, 与序列紧矛盾. 紧集闭且有界的结论也由序列紧得到. 这些是此处所需的度量紧性先修.

\begin{theorem}\label{thm:compact-approx}
Hilbert 空间中的集合 $C$ 紧, 当且仅当它闭、有界, 并对每个 $\eps>0$ 存在有限维子空间 $M$ 使 $\sup_{x\in C}\dist(x,M)<\eps$.
\end{theorem}
\begin{proof}
若 $C$ 紧, 取有限 $\eps/2$ 网 $x_1,\ldots,x_N$, 令 $M$ 为这些点的张成. 每个 $x$ 距离其中某个点小于 $\eps/2$, 因而距离 $M$ 小于 $\eps$, 给出必要性.

反之, 给定 $\eps>0$, 选 $M$ 使所有向量距 $M$ 小于 $\eps/3$. 正交投影 $P_M$ 的范数至多1, 故 $P_M(C)$ 有界. 它位于有限维空间, 其闭包紧, 可用有限个半径 $\eps/3$ 的球覆盖. 每个 $x\in C$ 距离其投影小于 $\eps/3$, 因而这些球的中心给出 $C$ 的有限 $\eps$ 网. 所以 $C$ 全有界, 闭性和 Hilbert 完备性给出 $C$ 完备, 从而紧.
\end{proof}

\begin{corollary}\label{cor:smalltails}
在有可数正交归一基 $(e_j)$ 的 Hilbert 空间中, 闭有界集 $C$ 紧, 当且仅当有一致小尾项:
\[
\lim_{N\to\infty}\sup_{x\in C}\sum_{j>N}|\ip{x}{e_j}|^2=0.
\]
\end{corollary}
\begin{proof}
若一致小尾项成立, Parseval 给出 $\norm{x-P_Nx}^2$ 等于尾和, 故有限维逼近判别给出紧性. 若 $C$ 紧, 对任意 $\eps>0$ 选有限 $\eps/3$ 网 $x_1,\ldots,x_m$. 对每个网点, $P_Nx_j\to x_j$; 取共同 $N$, 使每个误差小于 $\eps/3$. 因 $\norm{I-P_N}\le1$, 对接近网点的任意 $x$ 有
$\norm{(I-P_N)x}\le\norm{x-x_j}+\norm{(I-P_N)x_j}<2\eps/3$.
这个共同 $N$ 适用于所有 $x$, 给出尾和一致趋于零.
\end{proof}

\begin{example}
核对 Hilbert 立方体 $C=\{a\in\ell^2:|a_j|\le2^{-j}\}$ 的紧性.
\end{example}
\begin{solution}
坐标泛函连续, 所以这些闭不等式的交闭. 范数平方不超过 $\sum_j4^{-j}$, 故有界. 尾和统一受 $\sum_{j>N}4^{-j}\to0$ 控制, 用一致小尾项判别得紧. 这个集合包含无穷多个方向, 说明无限维空间中仍可以有紧集.
\end{solution}

\originalsection{有限秩与范数闭包}
\begin{definition}
有界算子 $T:H_1\to H_2$ 称有限秩, 若 $\ran T$ 有限维; 称紧算子, 若 $T$ 把定义域闭单位球映为相对紧集, 即像的闭包紧. 等价地, 每个有界序列的像都有范数收敛子列.
\end{definition}
正向把有界输入序列除以一个共同常数, 放入单位球, 在像的紧闭包中抽取子列, 再缩放回来. 反向若单位球像不全有界, 存在 $\eps>0$, 可递归选择像中两两距离至少 $\eps$ 的点及相应单位输入, 其像无 Cauchy 子列, 与假设矛盾. 故像全有界; $H_2$ 完备使其闭包紧. 有限秩算子紧, 因为其有界像位于有限维空间中. 若 $v_1,\ldots,v_r$ 是值域的正交归一基, 则
$Tx=\sum_{j=1}^r\ip{Tx}{v_j}v_j=\sum_{j=1}^r\ip{x}{T^*v_j}v_j$.
这给出有限秩算子的有限泛函表示, 不要求定义域本身有限维. 当定义域和值域均为同一 $H$, 对有限组 $v_j,T^*v_j$ 的联合张成做 Gram--Schmidt, 再展开这两组向量, 即得到原第19讲在同一有限正交系下的矩阵表示.

\begin{theorem}\label{thm:compactnorm}
紧算子的算子范数极限仍紧. Hilbert 空间之间的紧算子恰为有限秩算子的算子范数闭包.
\end{theorem}
\begin{proof}
若 $T_n\to T$ 于算子范数, 对 $\eps>0$ 选 $n$ 使 $\norm{T-T_n}<\eps/3$. $T_n$ 的单位球像有有限 $\eps/3$ 网, 因而同一网对 $T$ 的单位球像给出误差小于 $2\eps/3$. 所以 $T$ 的单位球像全有界, 它在完备值域中的闭包也完备且全有界, 从而紧. 这里值域完备性不可略去; 本章的 Hilbert 空间已满足.

若 $T$ 紧, 取其单位球像闭包 $C$. 有限维逼近判别给出 $M$ 使 $\norm{Tx-P_MTx}<\eps$ 对所有单位球向量成立. $P_MT$ 有限秩, 且 $\norm{T-P_MT}\le\eps$. 对逐渐减小的 $\eps$ 选择一列有限秩算子, 得范数逼近. 反方向由有限秩紧和前半结论得到.
\end{proof}

\begin{corollary}\label{cor:compactideal}
紧算子对加法、数乘、伴随以及左右复合有界算子封闭.
\end{corollary}
\begin{proof}
有限秩算子的和、数乘仍有限秩, 两侧复合也仍有限秩. 伴随由前述有限泛函表示仍有限秩: 若 $Tx=\sum_j\ip{x}{w_j}v_j$, 则 $T^*y=\sum_j\ip{y}{v_j}w_j$. 对紧算子的有限秩逼近取这些运算, 复合范数估计与 $\norm{T^*}=\norm{T}$ 保证逼近仍在算子范数中收敛, 再用定理\ref{thm:compactnorm}.
\end{proof}
Hilbert 空间的正交投影使有限秩逼近成立. 不能未经说明把此结论套到所有 Banach 空间之间的紧算子.

\originalsection{弱收敛与积分算子}
\begin{theorem}\label{thm:compactweak}
Hilbert 空间之间的有界算子 $T$ 紧, 当且仅当 $x_n\rightharpoonup x$ 总能推出 $Tx_n\to Tx$ 于范数.
\end{theorem}
\begin{proof}
若 $T$ 紧, 弱收敛序列范数有界, 所以任意子列的像都有范数收敛的再子列. 同时 $T$ 保持弱收敛, 范数极限也为弱极限, 唯一性迫使所有这样的极限都为 $Tx$. 若整个 $Tx_n$ 不趋于 $Tx$, 会有一条子列的距离始终至少为某 $\eps>0$, 与其再子列趋于 $Tx$ 矛盾.

反之, 给定有界 $(x_n)$, Hilbert 弱子列定理给出 $x_{n_k}\rightharpoonup x$. 假设使 $Tx_{n_k}\to Tx$ 于范数. 因而每个有界序列的像都有收敛子列, 算子紧.
\end{proof}

\begin{theorem}\label{thm:continuouskernelcompact}
若 $k\in C([0,1]^2)$, 则 $Tf(x)=\int_0^1k(x,y)f(y)\dd y$ 在 $L^2([0,1])$ 上紧. 更一般地, 对平方可积核有 $\norm{T}\le\norm{k}_{L^2([0,1]^2)}$; 本册只需连续核的紧性.
\end{theorem}
\begin{proof}
逐个 $x$ 使用 Cauchy--Schwarz,
$|Tf(x)|^2\le(\int|k(x,y)|^2\dd y)\norm{f}_2^2$.
对 $x$ 积分即得算子范数界, 非负积分用 Tonelli 解释. 连续核一致连续, 将正方形分成小矩形, 在每格取一个核值, 得有限和核 $k_N(x,y)=\sum_{i,j}c_{ij}\ind_{I_i}(x)\ind_{I_j}(y)$, 满足 $\norm{k-k_N}_\infty\to0$. 对应算子 $T_N$ 的像位于有限个 $\ind_{I_i}$ 的张成中, 所以有限秩. 核范数界给出 $\norm{T-T_N}\le\norm{k-k_N}_2\le\norm{k-k_N}_\infty\to0$, 故 $T$ 紧. 边界格线取值不影响 $L^2$ 算子.
\end{proof}
对平方可和矩阵 $a_{ij}$ 也有相同机制: $Tx=(\sum_ja_{ij}x_j)_i$ 满足 $\norm{T}\le(\sum_{i,j}|a_{ij}|^2)^{1/2}$. 截断到有限正方形, 剩余平方和控制算子范数误差, 所以算子紧. 这是 Hilbert--Schmidt 类型例子的必要部分; 更一般的迹类和 Schatten 理论不在主源覆盖范围.

\begin{exercise}
对 $a\in\ell^\infty$, 定义对角算子 $D_ax=(a_jx_j)$. 证明 $\norm{D_a}=\sup_j|a_j|$, 且 $D_a$ 紧当且仅当 $a_j\to0$.
\end{exercise}
\begin{solution}
平方求和给范数上界, 单位向量 $e_j$ 给所有坐标的下界, 所以范数等式成立. 若 $a_j\to0$, 截断到前 $N$ 个坐标的有限秩算子与 $D_a$ 的范数差为 $\sup_{j>N}|a_j|\to0$, 所以紧. 若不趋于零, 有无限子列 $|a_{j_k}|\ge\eps>0$. 对不同 $k,l$, 像 $a_{j_k}e_{j_k},a_{j_l}e_{j_l}$ 的距离平方至少为 $2\eps^2$, 无 Cauchy 子列, 故不紧.
\end{solution}

\begin{exercise}
证明无限维 Hilbert 空间的恒等算子不紧, 并说明任何紧算子都不可能在此空间上具有有界逆.
\end{exercise}
\begin{solution}
无限维空间可逐步取不在已有有限张成中的向量并做 Gram--Schmidt, 得无限正交归一序列. 它在单位球中, 却没有范数收敛子列, 所以恒等算子不紧. 若紧 $T$ 有有界逆, 则 $I=T^{-1}T$ 由复合封闭性紧, 与前述矛盾.
\end{solution}

\originalchapter{谱与紧自伴算子}\label{ch:spectrum}
\originalsection{预解式与谱}
本章在非零复 Hilbert 空间 $H$ 上讨论谱. 实空间上的自伴结论可通过复化理解; 特征值与特征向量的主要论证同样适用于实空间. 无限维的不可逆性不只来自核, 还可能来自非满射或不闭的值域.

\begin{theorem}[Neumann 级数]\label{thm:neumann}
若 $A\in\B(H)$ 且 $\norm{A}<1$, 则 $I-A$ 可逆, 其有界逆为 $\sum_{n=0}^\infty A^n$, 范数不超过 $1/(1-\norm{A})$. 可逆算子集合在算子范数下开.
\end{theorem}
\begin{proof}
因为 $\norm{A^n}\le\norm{A}^n$, 该级数在 Banach 空间 $\B(H)$ 中绝对收敛. 部分和 $S_N$ 满足 $(I-A)S_N=S_N(I-A)=I-A^{N+1}$. 乘法连续且余项范数趋于零, 取极限得到双侧逆. 几何级数给出范数估计.

若 $T$ 可逆, 且 $\norm{T^{-1}E}<1$, 则 $T+E=T(I+T^{-1}E)$ 可逆. 特别地 $\norm{E}<1/\norm{T^{-1}}$ 足够, 所以可逆集合开.
\end{proof}

\begin{definition}
预解集为 $\rho(T)=\{\lambda\in\C:T-\lambda I\text{ 有有界双侧逆}\}$, 谱为 $\sigma(T)=\C\setminus\rho(T)$. 预解式取 $R(\lambda)=(T-\lambda I)^{-1}$. 若 $(T-\lambda I)x=0$ 有非零解, 则 $\lambda$ 称特征值, 此时必属于谱.
\end{definition}
有界逆定理保证这里的有界算子双射已经具有有界逆. 但单射不足以推出可逆.

\begin{proposition}\label{prop:spectrumcompact}
$\sigma(T)$ 闭且包含于 $\{\lambda:|\lambda|\le\norm{T}\}$. 在预解集上 $R$ 范数连续, 且有预解式恒等式
\[
R(\lambda)-R(\mu)=(\lambda-\mu)R(\lambda)R(\mu).
\]
\end{proposition}
\begin{proof}
可逆集合开, 而 $\lambda\mapsto T-\lambda I$ 范数连续, 所以预解集开, 谱闭. 若 $|\lambda|>\norm{T}$, 则 $T-\lambda I=-\lambda(I-T/\lambda)$ 由 Neumann 级数可逆. 两个逆的差等于
$R(\lambda)[(T-\mu I)-(T-\lambda I)]R(\mu)$, 即恒等式. 对固定 $\lambda_0$, 分解
$T-\lambda I=(T-\lambda_0I)(I-(\lambda-\lambda_0)R(\lambda_0))$,
Neumann 级数说明逆在 $\lambda_0$ 附近连续, 并给出局部幂级数.
\end{proof}

\begin{remark}
下面的可选说明额外调用复分析的 Liouville 定理, 证明一般复有界算子的谱非空: 若谱空, 则对每个 $x,y$, $\ip{R(\lambda)x}{y}$ 由局部幂级数为整函数. 对 $|\lambda|>\norm{T}$, Neumann 界给出 $\norm{R(\lambda)}\le1/(|\lambda|-\norm{T})$, 所以这些整函数在无穷远趋于零, 在有限闭圆盘上由连续性有界. Liouville 说明它们恒为零, 将推出每个 $R(\lambda)=0$, 与其为逆矛盾. 此一般非空性说明调用了明确的复分析定理, 本册不重复证明 Cauchy 积分公式和 Liouville. 下文自伴谱的非空性另有完整的实估计证明, 不依赖这条调用.
\end{remark}

\begin{example}
在 $\ell^2$ 上令 $Dx=(x_j/j)$. 求谱并说明零为什么不是特征值. 在 $L^2([0,1])$ 上令 $Mf(x)=xf(x)$, 求谱和特征值.
\end{example}
\begin{solution}
$D$ 的每个 $1/j$ 是特征值. 零在谱中, 因 $De_j=e_j/j\to0$, 若存在有界逆则 $1=\norm{e_j}\le\norm{D^{-1}}/j$, 矛盾; 但 $Dx=0$ 只可能 $x=0$. 若 $\lambda\notin\{0,1,1/2,\ldots\}$, 它与这个闭集距离正, 对角逆 $(D-\lambda I)^{-1}x=(x_j/(1/j-\lambda))$ 有界, 故谱恰为上述集合.

若 $\lambda\notin[0,1]$, 则 $1/(x-\lambda)$ 有界, 给出 $M-\lambda I$ 的逆. 若 $\lambda\in[0,1]$, 取区间 $I_n=[0,1]\cap(\lambda-1/n,\lambda+1/n)$, 以 $f_n=\ind_{I_n}/\sqrt{|I_n|}$ 归一化, 得 $\norm{f_n}_2=1$, $\norm{(M-\lambda I)f_n}_2\le1/n$, 否定有界逆. 所以谱为 $[0,1]$. 但 $(x-\lambda)f(x)=0$ 只允许 $f$ 在单点上可能非零, 在 $L^2$ 中即为零, 因而没有特征值.
\end{solution}

\originalsection{自伴谱的端点}
\begin{lemma}\label{lem:selfnorm}
若 $A=A^*$, 则 $\ip{Ax}{x}$ 实, 且 $\norm{A}=\sup_{\norm{x}=1}|\ip{Ax}{x}|$.
\end{lemma}
\begin{proof}
共轭对称与自伴性给出 $\overline{\ip{Ax}{x}}=\ip{x}{Ax}=\ip{Ax}{x}$. 记右端为 $a$, Cauchy--Schwarz 给出 $a\le\norm{A}$. 对单位 $x$ 若 $Ax\ne0$, 取单位 $y=Ax/\norm{Ax}$. 则 $\ip{Ax}{y}=\norm{Ax}$ 为实数. 自伴性和展开给出
\[
\norm{Ax}=\tfrac14\{\ip{A(x+y)}{x+y}-\ip{A(x-y)}{x-y}\}
\le\tfrac a4(\norm{x+y}^2+\norm{x-y}^2)=a.
\]
最后用平行四边形及两个单位向量. $Ax=0$ 时也有此界, 对所有单位 $x$ 取上确界得到反向不等式.
\end{proof}

\begin{theorem}\label{thm:selfspec}
设 $A=A^*$, $a=\inf_{\norm{x}=1}\ip{Ax}{x}$, $b=\sup_{\norm{x}=1}\ip{Ax}{x}$. 则 $\sigma(A)\subset[a,b]\subset[-\norm{A},\norm{A}]$, 且 $a,b\in\sigma(A)$. 特别地自伴谱非空, $\norm{A}=\max_{\lambda\in\sigma(A)}|\lambda|$.
\end{theorem}
\begin{proof}
若 $\lambda=s+\ii t$, $t\ne0$, 则
$|t|\norm{x}^2=|\operatorname{Im}\ip{(A-\lambda I)x}{x}|\le\norm{(A-\lambda I)x}\norm{x}$.
所以 $A-\lambda I$ 有下界 $|t|$, 单射且值域闭; 后一个结论由输入序列的 Cauchy 估计与完备性得到. 其伴随 $A-\overline\lambda I$ 也单射, 因而值域的正交补为零. 值域同时闭且稠密, 就是整个空间, 所以 $\lambda\in\rho(A)$.

若实 $\lambda>b$, 则 $|\ip{(A-\lambda I)x}{x}|\ge(\lambda-b)\norm{x}^2$, 同样给出下界. 自伴性使伴随也单射, 因而满射. $\lambda<a$ 同理. 这说明谱包含于 $[a,b]$.

为证明端点在谱中, 令 $c=(a+b)/2$, $B=A-cI$, $r=(b-a)/2$. 引理\ref{lem:selfnorm}给出 $\norm{B}=r$. 若 $r=0$, 则 $B=0$, $A=cI$, 所以 $a=b=c$ 是特征值. 若 $r>0$, 选单位 $x_n$ 使 $\ip{Bx_n}{x_n}\to r$. 则关键估计为
\[
\norm{(B-rI)x_n}^2
=\norm{Bx_n}^2-2r\ip{Bx_n}{x_n}+r^2
\le2r(r-\ip{Bx_n}{x_n})\to0.
\]
若 $A-bI=B-rI$ 有有界逆, 将得 $1\le\norm{(A-bI)^{-1}}\norm{(A-bI)x_n}\to0$, 矛盾. 所以 $b\in\sigma(A)$. 对下端选 $\ip{Bx_n}{x_n}\to-r$, 使用 $B+rI$ 的同样估计, 得 $a\in\sigma(A)$. 上述范数估计补齐了原稿从二次型趋于零直接跳到向量范数趋于零的缺口. 最后 $\norm{A}=\max(|a|,|b|)$, 由端点在谱中得到谱范数公式.
\end{proof}

\originalsection{Fredholm 相容条件}
\begin{theorem}\label{thm:compact-eigen}
设 $A$ 紧且自伴. 非零特征值为实数, 其特征空间 $E_\lambda=\ker(A-\lambda I)$ 有限维. 不同特征值的特征空间正交. 对每个 $\eps>0$, 绝对值至少为 $\eps$ 的特征值只有有限个, 包括各自的有限重数.
\end{theorem}
\begin{proof}
单位特征向量 $x$ 满足 $\lambda=\ip{Ax}{x}$, 所以 $\lambda$ 实. 若 $E_\lambda$ 无限维, 可取无限正交归一系 $(x_n)$ 于其中; 像 $Ax_n=\lambda x_n$ 的不同项距离为 $\sqrt2|\lambda|>0$, 没有收敛子列, 违反紧性. 若 $x\in E_\lambda,y\in E_\mu$, 则
$\lambda\ip{x}{y}=\ip{Ax}{y}=\ip{x}{Ay}=\mu\ip{x}{y}$, 不同实特征值迫使正交.

如果绝对值至少为 $\eps$ 的特征空间合计无限维, 可从这些两两正交空间中选无限正交归一特征向量 $x_n$, 对应 $|\lambda_n|\ge\eps$. 不同像的距离平方为 $|\lambda_n|^2+|\lambda_m|^2\ge2\eps^2$, 再次违反紧性. 所以每个这样的区域只有有限个重数. 取 $\eps=1/k$ 后并集至多可数, 非零特征值若按重数无限列出, 必趋于零.
\end{proof}

\begin{theorem}[自伴 Fredholm 选择定理]\label{thm:fredholm}
若 $A$ 紧且自伴, $\lambda\in\R\setminus\{0\}$, 则 $\ran(A-\lambda I)$ 闭, 且等于 $E_\lambda^\perp$. 因而方程 $(A-\lambda I)x=f$ 可解当且仅当 $f\perp E_\lambda$; 在 $E_\lambda^\perp$ 内解唯一且连续依赖 $f$. 若 $E_\lambda=\{0\}$, 则 $A-\lambda I$ 可逆. 所以每个非零谱点都是特征值.
\end{theorem}
\begin{proof}
记 $B=A-\lambda I$, $M=(\ker B)^\perp$. 先证明 $B$ 在 $M$ 上有正下界. 若无, 可取 $x_n\in M$, $\norm{x_n}=1$, $Bx_n\to0$. 紧性给出子列 $Ax_{n_k}\to y$. 由 $x_{n_k}=(Ax_{n_k}-Bx_{n_k})/\lambda$ 得 $x_{n_k}\to x=y/\lambda$. 闭性使 $x\in M$, 连续性使 $Bx=0$, 同时 $\norm{x}=1$. 这与 $M\cap\ker B=\{0\}$ 矛盾. 因而存在 $c>0$ 使 $\norm{Bx}\ge c\norm{x}$ 对 $x\in M$ 成立.

任意 $Bz_n\to f$, 将 $z_n$ 分解为核分量加 $x_n\in M$, 则 $Bz_n=Bx_n$. 下界说明 $(x_n)$ 为 Cauchy, $M$ 完备给出 $x_n\to x\in M$, 从而 $Bx=f$. 值域闭. 自伴性和伴随值域公式给出 $\ran B=(\ker B)^\perp=M$. 下界给出逆在此子空间上的范数界 $1/c$, 核与 $M$ 无交给出唯一性. 若核零, 值域就是全空间, 所以可逆. 非实谱点已经由自伴谱定理排除, 得最后结论.
\end{proof}

\originalsection{谱分解与函数演算}
\begin{theorem}[紧自伴谱定理]\label{thm:spectral}
设 $A$ 为可分 Hilbert 空间上的紧自伴算子. 存在有限或可数正交归一特征向量 $(u_j)$, 对应所有非零特征值按重数列出的 $(\lambda_j)$, 可使 $|\lambda_1|\ge|\lambda_2|\ge\cdots$. 则
\[
H=\ker A\oplus\overline{\Span\{u_j\}},\qquad
Ax=\sum_j\lambda_j\ip{x}{u_j}u_j.
\]
非零特征值无限时 $\lambda_j\to0$, 且有限部分和算子按算子范数趋于 $A$. 在核中补上一组正交归一基, 得到全空间由特征向量组成的基.
\end{theorem}
\begin{proof}
若 $A=0$, 只需核中的基. 若 $A\ne0$, 自伴谱的范数端点至少有一个非零, 由 Fredholm 选择定理是特征值; 可取单位特征向量 $u_1$, 满足 $|\lambda_1|=\norm{A}$. $u_1^\perp$ 对 $A$ 不变, 因 $\ip{Ax}{u_1}=\ip{x}{Au_1}=\lambda_1\ip{x}{u_1}$. 在此闭子空间上限制算子仍紧、自伴. 若限制非零, 重复选择范数大小的特征值与单位特征向量; 若限制为零, 过程终止.

第 $N$ 步之后的余算子
$A_N=A-\sum_{j=1}^N\lambda_j\ip{\cdot}{u_j}u_j$
在已选有限张成上为零, 在其正交补上等于 $A$. 下一步选择给出 $\norm{A_N}=|\lambda_{N+1}|$, 所以绝对值递减. 如果无限进行, 定理\ref{thm:compact-eigen}说明这些按重数选出的值趋于零. 因而 $\norm{A_N}\to0$, 得算子范数级数展开. 若有限终止, 展开就是有限恒等式.

令 $M=\overline{\Span\{u_j\}}$. 若 $x\perp M$, 所有展开系数为零, 因而 $Ax=0$. 若 $Ax=0$, 则 $0=\ip{Ax}{u_j}=\lambda_j\ip{x}{u_j}$, 非零 $\lambda_j$ 给出 $x\perp M$. 所以 $\ker A=M^\perp$, 得正交分解. 核是可分 Hilbert 空间的闭子空间, 也可分: 对环境中的可数稠密集投影到核, 得可数稠密集. 因而核有正交归一基, 与 $(u_j)$ 合并即得全空间基.
\end{proof}
同时有 $\overline{\ran A}=(\ker A)^\perp=M$. 一般不能去掉值域闭包: 对角算子 $x_j\mapsto x_j/j$ 的值域稠密却不闭. 谱分解的无限和在 $H$ 的范数中解释, 不把非闭值域本身当作完整的 Hilbert 空间.

\begin{proposition}\label{prop:calculus}
在上述谱基下, 给定谱上的有界函数 $\phi$, 可定义
\[
\phi(A)x=\phi(0)P_{\ker A}x+\sum_j\phi(\lambda_j)\ip{x}{u_j}u_j,
\]
当核为零时第一项省略, 只需给实际谱点定义 $\phi$. 该算子有界, 范数为 $\sup_{t\in\sigma(A)}|\phi(t)|$ 若 $\phi$ 连续. 它满足和、积及共轭函数对应和、积及伴随. 若 $A$ 非负, 则有非负自伴平方根 $A^{1/2}$, 它紧且 $(A^{1/2})^2=A$.
\end{proposition}
\begin{proof}
在每个正交坐标上乘一个有界标量, Parseval 保证级数收敛并给出范数上界. 单位特征向量给出每个非零坐标的下界; 核非零时核中的单位向量给出零点下界. 若零只为谱的积聚点而非特征值, 连续性使 $\phi(\lambda_j)\to\phi(0)$, 所以这些下界仍逼近零点值. 这说明连续函数情形的范数等式. 和、积、伴随逐坐标核对, 有界性允许把有限部分和的等式取极限.

若 $\ip{Ax}{x}\ge0$, 单位特征向量给出 $\lambda_j>0$. 取 $\phi(t)=\sqrt t$ 于非负谱, 便得自伴非负算子, 逐坐标平方为 $A$. 特征值 $\sqrt{\lambda_j}\to0$ 时, 有限谱截断按算子范数逼近它; 有限情况直接有限秩. 因而平方根紧. 此为紧自伴情形所需的函数演算, 不在这里扩展到一般非紧自伴算子的谱测度理论.
\end{proof}

\begin{exercise}
给定紧自伴 $A$, 说明方程 $(I-A)x=f$ 在1为特征值及不为特征值时的可解性, 并给出相容时的解公式.
\end{exercise}
\begin{solution}
若1不是特征值, Fredholm 定理保证唯一解. 若1是特征值, 必须且只需 $f\perp E_1$. 谱分解中取 $x_0=P_{\ker A}f+\sum_{\lambda_j\ne1}\ip{f}{u_j}u_j/(1-\lambda_j)$. 非零特征值只能积聚于零, 所以排除等于1的有限重数后, 剩余 $|1-\lambda_j|$ 有正下界, 级数在 $H$ 中收敛. 每个坐标代入给出 $(I-A)x_0=f$. 全部解为 $x_0+z$, $z\in E_1$; 没有相容条件时不得把零分母项直接删去.
\end{solution}

\originalchapter{区间上的 Dirichlet 问题}\label{ch:dirichlet}
\originalsection{Green 算子}
最后的应用把微分方程转化为有界算子方程, 再由紧自伴谱理论判断可解性. 主源考虑实连续势 $V$ 与连续右端 $f$, 寻找 $u\in C^2([0,1])$ 满足 $-u''+Vu=f$, $u(0)=u(1)=0$. 为使用 Hilbert 理论, 中间的算子先定义在 $L^2([0,1])$ 上, 最后再证明解的经典正则性.

\begin{theorem}\label{thm:green}
定义 $G(x,y)=\min(x,y)(1-\max(x,y))$, 以及 $Af(x)=\int_0^1G(x,y)f(y)\dd y$. 则 $A$ 是 $L^2([0,1])$ 上的紧自伴算子. 它将 $L^2$ 有界地映到连续且端点为零的函数. 若 $f$ 连续, 则 $u=Af\in C^2$, 并为 $-u''=f$, $u(0)=u(1)=0$ 的唯一解.
\end{theorem}
\begin{proof}
核连续且实对称, 连续核紧性定理和伴随核公式分别给出紧性与自伴性. 因 $\norm{f}_1\le\norm{f}_2$, 有 $\norm{Af}_\infty\le\norm{G}_\infty\norm{f}_2$. 核的一致连续性还给出
$|Af(x)-Af(z)|\le\sup_y|G(x,y)-G(z,y)|\norm{f}_2\to0$,
所以输出连续. 核在 $x=0,1$ 时均为零, 因而端点值为零.

对连续 $f$ 写成
\[
u(x)=(1-x)\int_0^xyf(y)\dd y+x\int_x^1(1-y)f(y)\dd y.
\]
求导时两个移动端点项抵消, 得 $u'(x)=-\int_0^xyf(y)\dd y+\int_x^1(1-y)f(y)\dd y$. 再求导得 $u''(x)=-xf(x)-(1-x)f(x)=-f(x)$. 连续性说明 $u\in C^2$.

若两个解之差为 $w$, 则 $w''=0$, 微积分基本定理说明 $w$ 为仿射函数, 两端为零迫使 $w=0$. 这里核为正的 $G$, 原第23讲所写核的相反号会解成 $u''=f$, 因而本稿明确校正符号.
\end{proof}

\originalsection{正弦谱与平方根}
\begin{theorem}\label{thm:greenspec}
令 $e_k(x)=\sqrt2\sin(k\pi x)$. 则 $Ae_k=(k\pi)^{-2}e_k$, $\ker A=\{0\}$, 并且
\[
Af=\sum_{k\ge1}\frac{\ip{f}{e_k}}{k^2\pi^2}e_k,
\qquad \norm{A}=\pi^{-2}.
\]
\end{theorem}
\begin{proof}
函数 $e_k/(k\pi)^2$ 满足 $-u''=e_k$ 和零端点条件, 由 Green 解的唯一性即得特征等式. 第\ref{ch:fourier}章已证明正弦系为基, 所以在有限部分和上使用线性性, 再由 $A$ 连续取极限, 得级数展开. 所有对角系数非零, 所以 $Af=0$ 蕴含各 Fourier 系数为零, 即核零. Parseval 给出范数上界 $1/\pi^2$, 在 $e_1$ 上达到. 不需要假设无限级数可以逐点求两次导数.
\end{proof}

由 Parseval 和刚得到的展开, $\ip{Af}{f}=\sum_{k\ge1}|\ip{f}{e_k}|^2/(k^2\pi^2)\ge0$. 因而 $A$ 非负, 可以应用上一章的谱平方根构造.

\begin{proposition}\label{prop:greensqrt}
令 $B=A^{1/2}$, 即 $Bf=\sum_{k\ge1}\ip{f}{e_k}e_k/(k\pi)$. 则 $B$ 紧、自伴, $B^2=A$, $\norm{B}=1/\pi$. 此外 $B:L^2\to C([0,1])$ 有界, 输出端点为零, 且 $\norm{Bf}_\infty\le\norm{f}_2/\sqrt3$.
\end{proposition}
\begin{proof}
紧性、自伴性与平方等式由谱函数演算得到; 也可用对角系数 $1/(k\pi)\to0$ 的有限截断直接验证. 对每个 $x$,
\[
\sum_{k\ge1}\left|\frac{\ip{f}{e_k}}{k\pi}e_k(x)\right|
\le\frac{\sqrt2}{\pi}\left(\sum_{k\ge1}|\ip{f}{e_k}|^2\right)^{1/2}
\left(\sum_{k\ge1}k^{-2}\right)^{1/2}
=\frac{\norm{f}_2}{\sqrt3}.
\]
各项上界之和有限, 所以 Weierstrass 判别保证级数绝对一致收敛, 极限连续并有给出的上确界估计. 每一有限部分和在两端为零, 一致极限也为零. 这里使用正确的平方根范数 $\norm{f}_2$, 不把它误写成范数的平方.
\end{proof}

\originalsection{非负势的存在、唯一与稳定}
从积分方程 $u+AM_Vu=Af$ 出发, $A=B^2$ 提示先构造 $u=Bv$, 将辅助算子写成自伴的 $BM_VB$. 接下来的证明先解辅助方程、再验证输出满足原问题; 这一存在性构造无需预先断言每个候选解都属于 $B$ 的值域.
\begin{theorem}\label{thm:dirichlet}
若 $V\in C([0,1])$ 为实非负函数, 则每个 $f\in C([0,1])$ 都有唯一 $u\in C^2([0,1])$ 满足 $-u''+Vu=f$, $u(0)=u(1)=0$. 该解满足 $\norm{u}_2\le\pi^{-2}\norm{f}_2$.
\end{theorem}
\begin{proof}
令 $M_V$ 为乘法算子, $M_Vg=Vg$. 有 $\norm{M_V}\le\norm{V}_\infty$, 实值性使 $M_V$ 自伴, 非负性使 $\ip{M_Vg}{g}=\int V|g|^2\ge0$. 它本身不要求紧. 令 $C=BM_VB$, 因 $B$ 紧且 $M_V$ 有界, 得 $C$ 紧; 伴随反转次序给出 $C=C^*$.

对任意 $v$, $\ip{(I+C)v}{v}=\norm{v}^2+\int V|Bv|^2\ge\norm{v}^2$. 所以 $I+C$ 核为零, 且有下界1. 将 Fredholm 定理用于紧自伴 $C$ 和非零参数 $-1$, 得 $I+C$ 满射并可逆. 同一估计与 Cauchy--Schwarz 给出 $\norm{(I+C)^{-1}}\le1$.

定义 $v=(I+C)^{-1}Bf$, $u=Bv$. 这一定义没有使用未必存在的 $B^{-1}$. 对 $(I+C)v=Bf$ 左乘 $B$, 使用 $B^2=A$, 得 $u+AM_Vu=Af$, 即 $u=A(f-Vu)$. $B:L^2\to C$ 的有界性使 $u$ 连续且端点为零, 因而 $f-Vu$ 连续. Green 定理于是使 $u\in C^2$, 且 $-u''=f-Vu$. 这完成由 Hilbert 空间解到经典解的提升.

若两个经典解之差为 $w$, 分部积分及端点条件给出
\[
0=\int_0^1(-w''+Vw)\overline w\dd x
=\int_0^1|w'|^2\dd x+\int_0^1V|w|^2\dd x.
\]
两项非负, 所以 $w'=0$ 几乎处处; 连续性进一步给出处处为零. 端点条件迫使常数 $w=0$. 最后由构造有
$\norm{u}_2\le\norm{B}\norm{(I+C)^{-1}}\norm{B}\norm{f}_2\le\pi^{-2}\norm{f}_2$, 得稳定估计.
\end{proof}
用谱平方根使 $BM_VB$ 自伴, 避免了 $AM_V$ 一般不自伴的困难. 即使 $A$ 与 $M_V$ 分别自伴, 也不能因此断言它们的乘积自伴; 需要交换性, 或如这里的对称夹心结构.

\begin{example}
当 $V=0$, $f=1$ 时求解; 当 $V=-\pi^2$ 时说明唯一性为何可能失效.
\end{example}
\begin{solution}
第一种情形直接积分得 $u(x)=x(1-x)/2$, 它满足 $-u''=1$ 和端点条件, 与 Green 算子结果一致. 第二种情形齐次方程为 $-u''-\pi^2u=0$, 有非零解 $\sin(\pi x)$ 满足两个端点条件. 因而非负势保证的核零性质不再自动成立. 若此时有非齐次经典解, 与 $\sin(\pi x)$ 配对并分部积分, 得必要相容条件 $\int_0^1f(x)\sin(\pi x)\dd x=0$; 对 $f=1$ 该积分为 $2/\pi\ne0$, 因而无解.
\end{solution}

\begin{exercise}
取常数势 $V=v\ge0$ 与 $f\in C([0,1])$. 给出解的正弦系数, 并说明级数与经典解的关系.
\end{exercise}
\begin{solution}
经典解存在唯一. 对方程与 $e_k$ 配对, 两次分部积分的边界项由 $u(0)=u(1)=e_k(0)=e_k(1)=0$ 消失, 得 $((k\pi)^2+v)\ip{u}{e_k}=\ip{f}{e_k}$. 因而
$u=\sum_{k\ge1}\ip{f}{e_k}e_k/((k\pi)^2+v)$ 于 $L^2$ 成立. 系数绝对值乘 $\sqrt2$ 的和由 Cauchy--Schwarz 及 $\sum k^{-4}<\infty$ 控制, 所以级数也一致收敛到连续函数. 它与已经证明存在的经典解在 $L^2$ 中一致, 两个连续函数几乎处处一致必处处一致. 无须未经验证地对该级数逐项求两次导数.
\end{solution}

\originalchapter{课程作业}\label{ch:assignments}
\originalsection{作业1}\label{sec:ps01}
原件3页, 其中数学题面2页. 采用 $1/\infty=0$; 对偶中的配对不加共轭. 原件规定作业经 Gradescope 在24:00前提交, 未给本份的具体日期.
\begin{exercise}\label{ps01:1}
原题1. (a) 对 $n\in\N$, $a_k,b_k\in\R$ 及 $1<p<\infty$, $1/p+1/q=1$, 证明 $\sum_{k=1}^n|a_kb_k|\le(\sum|a_k|^p)^{1/p}(\sum|b_k|^q)^{1/q}$. 原提示: 对 $A,B>0$, $0<t<1$, 在 $x>0$ 上由 $t x+(1-t)B-x^tB^{1-t}$ 在 $x=B$ 处取最小值得 $A^tB^{1-t}\le tA+(1-t)B$.

(b) 对 $1\le p<\infty$, 证明 $(\sum|a_k+b_k|^p)^{1/p}\le(\sum|a_k|^p)^{1/p}+(\sum|b_k|^p)^{1/p}$. 原提示: 先把 $|a_k+b_k|^p$ 控制为 $(|a_k|+|b_k|)|a_k+b_k|^{p-1}$, 再用(a).
\end{exercise}
\begin{proof}[AI 独立解答]
(a) 这是引理\ref{lem:holder-seq}的有限版本. 为核对原提示, 上述函数的导数为 $t[1-(B/x)^{1-t}]$, 在 $x<B$ 时负、$x>B$ 时正. 取 $t=1/p$, $A=s^p,B=r^q$ 得 $sr\le s^p/p+r^q/q$, 零值由连续性得到. 两列范数非零时分别归一化再求和, 右端为 $1/p+1/q=1$; 任一范数为零时左端也为零.

(b) $p=1$ 逐项使用三角不等式. $p>1$ 时令 $S=(\sum|a_k+b_k|^p)^{1/p}$. 若 $S=0$ 结论成立; 否则原提示与(a)给出 $S^p\le(\norm a_p+\norm b_p)S^{p-1}$, 因为 $(p-1)q=p$. 除以 $S^{p-1}$ 即得结论.
\end{proof}
\begin{exercise}\label{ps01:2}
原题2. 证明对 $1\le p<\infty$, $\ell^p$ 是 Banach 空间; 须同时证明它是赋范空间及完备.
\end{exercise}
\begin{proof}[AI 独立解答]
正定性、齐次性由定义得到. 将上一题(b)用于前 $N$ 项再令 $N\to\infty$, 同时得到加法封闭和三角不等式; 数乘封闭显然. 完备性的逐坐标证明已在定理\ref{thm:lpseqbanach}写出: Cauchy 列 $a^{(n)}$ 的每一坐标有极限 $a_k$, 固定前 $J$ 项取极限, 得 $\sum_{k\le J}|a_k^{(n)}-a_k|^p\le\eps^p$; 令 $J\to\infty$ 得 $\norm{a^{(n)}-a}_p\le\eps$. 选定一个 $n$ 再用三角不等式保证 $a\in\ell^p$, 因而极限在本空间内.
\end{proof}
\begin{exercise}\label{ps01:3}
原题3. 已知 $\ell^\infty=C_b(\N,d_{\rm disc})$ 是 Banach 空间, 证明 $c_0=\{a\in\ell^\infty:a_k\to0\}$ 是闭子空间, 从而是 Banach 空间.
\end{exercise}
\begin{proof}[AI 独立解答]
线性组合的极限仍为零, 故为子空间. 若 $a^{(n)}\in c_0$ 且 $\norm{a^{(n)}-a}_\infty\to0$, 给定 $\eps>0$, 先令一致误差小于 $\eps/2$, 再令该固定 $a^{(n)}$ 的尾项小于 $\eps/2$, 得 $|a_k|<\eps$ 于充分大的 $k$. 所以 $a\in c_0$. 此为定理\ref{thm:lpseqbanach}末段的同一论证; 闭子空间完备性由命题\ref{prop:closedsubspace}保证.
\end{proof}
\begin{exercise}\label{ps01:4}
原题4. 对 $1\le p\le\infty$, 令 $S=\{a\in\ell^p:\norm a_p=1\}$. (a) 证明 $S$ 闭. (b) 证明 $S$ 不紧. 原提示: 用标准坐标向量 $e_n=(\delta_{kn})_k$ 构造无收敛子列的序列.
\end{exercise}
\begin{proof}[AI 独立解答]
(a) 范数连续, $S$ 是闭集 $\{1\}$ 的原像. (b) 每个 $e_n\in S$, 且 $m\ne n$ 时 $\norm{e_n-e_m}_p=2^{1/p}$ 于有限 $p$, 在 $p=\infty$ 时等于1. 任一子列都非 Cauchy, 从而没有收敛子列. 度量紧集中的每条序列应有收敛子列, 故 $S$ 不紧.
\end{proof}
\begin{exercise}\label{ps01:5}
原题5. 设 $1\le p<\infty$, $1/p+1/q=1$. (a) 对 $a\in\ell^p,b\in\ell^q$, 证明 $\sum_{k\ge1}|a_kb_k|\le\norm a_p\norm b_q$. (b) 对 $b\in\ell^q$, 证明 $F_b(a)=\sum a_kb_k$ 属于 $(\ell^p)^*$ 且 $\norm{F_b}=\norm b_q$. (c) 证明 $F:\ell^q\to(\ell^p)^*$, $b\mapsto F_b$, 是双射有界线性算子.
\end{exercise}
\begin{proof}[AI 独立解答]
(a) 当 $p>1$, 对有限部分和用题1(a), 右端不超过全列范数的积, 再取极限. 当 $p=1,q=\infty$, 逐项用 $|b_k|\le\norm b_\infty$ 后求和.

(b) 绝对收敛使 $F_b$ 良定义且线性, (a)给出范数上界. $p>1$ 时令 $S_N=\sum_{k\le N}|b_k|^q$; 若 $S_N>0$, 取 $a_k=\overline{b_k}|b_k|^{q-2}/S_N^{1/p}$ 于 $k\le N$ 且 $b_k\ne0$ 的坐标, 其余置零, 则 $\norm a_p=1$, $F_b(a)=S_N^{1/q}$. 取 $N\to\infty$ 得反向界. $p=1$ 时取相位调整后的 $e_k$, 得 $\norm{F_b}\ge|b_k|$, 再取上确界. 零 $b$ 情形直接成立.

(c) $F$ 线性且等距, 因而有界和单射. 对任一 $f\in(\ell^p)^*$, 令 $b_k=f(e_k)$. 上一有限支撑测试给出 $\sum_{k\le N}|b_k|^q\le\norm f^q$ 于 $p>1$, $p=1$ 时给出 $|b_k|\le\norm f$. 所以 $b\in\ell^q$. 对 $a$ 的有限截断用线性性, 再由截断在 $\ell^p$ 中收敛和 $f$ 连续性得到 $f(a)=\sum a_kb_k$. 这正是定理\ref{thm:dual-lp}的表示证明, 满射由此成立.
\end{proof}

\originalsection{作业2}\label{sec:ps02}
原件3页, 数学题面2页.
\begin{exercise}\label{ps02:1}
原题1. 设 $B$ 为 Banach 空间. (a) 若 $T\in\B(B,B)$ 且 $\norm{I-T}<1$, 证明 $T$ 可逆, 且 $\sum_{n=0}^\infty(I-T)^n$ 在算子范数中收敛于 $T^{-1}$. (b) 证明可逆算子的集合在 $\B(B,B)$ 中开.
\end{exercise}
\begin{proof}[AI 独立解答]
(a) 令 $A=I-T$, 几何级数界 $\sum\norm{A^n}\le\sum\norm A^n<\infty$ 与算子空间完备性保证和 $S$. 有 $T\sum_{n=0}^NA^n=(\sum_{n=0}^NA^n)T=I-A^{N+1}$; 取范数极限得 $TS=ST=I$. 这是定理\ref{thm:neumann}的证明, 此处只依赖算子空间完备性, 不依赖谱理论.

(b) 若 $T$ 可逆且 $\norm{T-S}<\norm{T^{-1}}^{-1}$, 则 $S=T[I-T^{-1}(T-S)]$, 方括号由(a)可逆, 从而 $S$ 可逆. 零空间的唯一算子可逆, 其算子空间本身也是开集, 另行满足结论.
\end{proof}
\begin{exercise}\label{ps02:2}
原题2. 设 $V$ 为赋范空间, $W\subsetneq V$ 为闭子空间. (a) 证明 $\norm{v+W}=\inf_{w\in W}\norm{v+w}$ 是 $V/W$ 上的范数. (b) 对每个 $\eps>0$, 证明存在 $v\in V$ 使 $\norm v=1$ 且 $\norm{v+W}\ge1-\eps$. 原提示: 对 $u\notin W$ 选接近距离下确界的代表元, 再归一化.
\end{exercise}
\begin{proof}[AI 独立解答]
(a) 改换代表元不改变被取下确界的陪集. 齐次性由 $W$ 为子空间得到; 对两个陪集各取误差小于 $\delta$ 的代表元, 用三角不等式再令 $\delta\downarrow0$, 得商范数的三角不等式. 若商范数为零, 则存在 $w_n\in W$ 使 $v+w_n\to0$, 闭性给出 $v\in W$; 反向显然. 此即定理\ref{thm:quotient}的范数部分.

(b) $\eps\ge1$ 时任取一个单位向量即可. $0<\eps<1$ 时选 $u\notin W$, 令 $d=\norm{u+W}>0$, 并选 $w\in W$ 使 $d\le\norm{u+w}<(1+\eps)d$. 则 $v=(u+w)/\norm{u+w}$ 为单位向量, 且 $\norm{v+W}=d/\norm{u+w}>1/(1+\eps)\ge1-\eps$.
\end{proof}
\begin{exercise}\label{ps02:3}
原题3. 若 $V$ 为 Banach 空间, $W\subsetneq V$ 为闭子空间, 证明上述商空间为 Banach 空间. 原提示: 对绝对可和的陪集级数选 $\norm{v_n+w_n}\le\norm{v_n+W}+2^{-n}$, 利用 $V$ 内的绝对级数收敛.
\end{exercise}
\begin{proof}[AI 独立解答]
按提示选代表元, 则 $\sum\norm{v_n+w_n}<\infty$. 定理\ref{thm:series}给出 $v=\sum(v_n+w_n)\in V$, 且
\[
\left\|v+W-\sum_{n=1}^N(v_n+W)\right\|\le\left\|v-\sum_{n=1}^N(v_n+w_n)\right\|\longrightarrow0.
\]
所以商空间中每个绝对级数收敛, 再用同一定理的逆向证明得完备. 这是定理\ref{thm:quotient}的完备性结论的另一种证明.
\end{proof}
\begin{exercise}\label{ps02:4}
原题4. 设 $V,W$ 为 Banach 空间, $T\in\B(V,W)$. (a) 证明 $\ker T$ 是闭子空间. (b) 称双射线性 $S:V_1\to V_2$ 为同构, 若 $S$ 及 $S^{-1}$ 都有界. 证明 $V/\ker T$ 与 $\ran T$ 同构当且仅当 $\ran T$ 闭. 原提示: 考察 $S(v+\ker T)=Tv$. 原件同构定义的箭头误写 $V_1\to V_1$, 此处按其定义域、值域及后续公式校正为 $V_1\to V_2$.
\end{exercise}
\begin{proof}[AI 独立解答]
(a) 线性性保证核为子空间, 连续性保证它是闭集 $\{0\}$ 的原像.

(b) $S$ 与代表元无关, 是线性双射, 且 $\norm{Tv}\le\norm T\norm{v+w}$ 对每个 $w\in\ker T$ 成立, 因而 $\norm S\le\norm T$. 若值域闭, 它与 $V/\ker T$ 都完备, 有界逆定理\ref{cor:inverse}给出 $S^{-1}$ 有界. 反之, 即使假设的是任何一个有界双向同构, 它也把完备的商空间映为完备空间: 值域的 Cauchy 列在逆像中 Cauchy, 再映回极限. 完备子空间在 $W$ 内闭, 故值域闭. 这不把“代数同构”误当成“拓扑同构”.
\end{proof}
\begin{exercise}\label{ps02:5}
原题5. 令 $W=\{a:\sum_{k\ge1}k|a_k|<\infty\}$, 配备 $\ell^1$ 范数. (a) 证明 $W$ 是 $\ell^1$ 的真稠密子空间, 从而不完备. (b) 定义 $(Ta)_k=ka_k$, 证明 $T:W\to\ell^1$ 的图闭, 但 $T$ 无界. (c) 证明 $S=T^{-1}:\ell^1\to W$ 有界且满射, 但不开. 原提示: 用有限截断证明稠密.
\end{exercise}
\begin{proof}[AI 独立解答]
(a) $c_{00}\subset W$, 截断在 $\ell^1$ 中逼近每个元素. 数列 $(k^{-2})$ 属于 $\ell^1$ 而不属于 $W$, 故 $W$ 真且不闭; 若完备就会在 $\ell^1$ 内闭, 矛盾.

(b) 若 $a^{(n)}\to a$ 于 $W$ 的所给范数, $Ta^{(n)}\to b$ 于 $\ell^1$, 坐标收敛给出 $b_k=ka_k$. 因 $b\in\ell^1$, $a\in W$ 且 $Ta=b$, 故图闭. 但 $\norm{Te_k}_1=k$, $\norm{e_k}_1=1$, 所以无界.

(c) $Sb=(b_k/k)_k$, 有 $\sum k|(Sb)_k|=\norm b_1$, $\norm{Sb}_1\le\norm b_1$, 故良定义、有界且与 $T$ 互逆. 若 $S$ 开, $S$ 的单位开球像应包含 $W$ 中某个半径 $\delta$ 的零点邻域. 取 $a=(\delta/2)e_k$ 且 $k>2/\delta$, 则 $\norm a_1<\delta$ 而 $\norm{Ta}_1>1$, 所以 $a$ 不在该像中, 矛盾. 这具体说明闭图定理不能去掉定义域完备性, 开映射定理不能去掉值域完备性.
\end{proof}

\originalsection{作业3}\label{sec:ps03}
原件3页, 数学题面2页.
\begin{exercise}\label{ps03:1}
原题1. 令 $M=\{a\in\ell^\infty:\lim a_k\text{存在}\}$. (a) 证明 $u(a)=\lim a_k$ 是 $M$ 上的有界线性泛函. (b) Hahn--Banach 给出延拓 $v\in(\ell^\infty)^*$, $v|_M=u$; 证明不存在 $b\in\ell^1$ 使 $v(a)=\sum a_kb_k$ 对所有 $a\in\ell^\infty$ 成立. 原提示: 测试 $e_n$.
\end{exercise}
\begin{proof}[AI 独立解答]
(a) 极限的线性性给出线性, $|u(a)|\le\norm a_\infty$ 给出有界. 常数数列1说明 $\norm u=1$. (b) 因 $e_n\in M$ 且极限为零, 若存在此 $b$, 则 $b_n=v(e_n)=0$ 对所有 $n$ 成立. 但 $v(1,1,\ldots)=1$, 与零配对矛盾. 这复用第\ref{ch:dual}章关于 $\ell^\infty$ 对偶的例题解答, 未对 $p=\infty$ 套用题\ref{ps01:5}的结论.
\end{proof}
\begin{exercise}\label{ps03:2}
原题2. 设 $V,W$ 为赋范空间, $T\in\B(V,W)$. (a) 定义 $T^\dagger f=f\circ T$ 于 $f\in W^*$; 证明 $T^\dagger\in\B(W^*,V^*)$ 且 $\norm{T^\dagger}=\norm T$. (b) 对 $1\le p\le\infty$, 证明右移 $Ra=(0,a_1,a_2,\ldots)$ 在 $\ell^p$ 上有界, 并求范数. (c) 对 $1\le p<\infty$ 及共轭指数 $q$, 用题\ref{ps01:5}的配对识别对偶为 $\ell^q$, 求 $R^\dagger b$. 原提示(a): Hahn--Banach 给出达到 $\norm{Tx}$ 的范数1泛函; (c)以 $R^\dagger e_1=0$ 为例.
\end{exercise}
\begin{proof}[AI 独立解答]
(a) 对每个 $f$, $f\circ T$ 线性且 $\norm{f\circ T}\le\norm f\norm T$; 关于 $f$ 也线性, 故得上界. 若 $Tx\ne0$, 推论\ref{cor:norming}给出 $\norm f=1$, $f(Tx)=\norm{Tx}$, 所以 $\norm{Tx}\le\norm{T^\dagger}\norm x$. $Tx=0$ 时同式直接成立. 对单位向量取上确界得反向界, 零空间亦满足等式.

(b) 右移不改变有限 $p$ 的幂和或上确界, 所以 $\norm{Ra}_p=\norm a_p$, $\norm R=1$.

(c) 绝对收敛允许重编号: $(R^\dagger b)(a)=\sum_{k\ge2}a_{k-1}b_k=\sum_{j\ge1}a_jb_{j+1}$. 因而 $R^\dagger b=(b_2,b_3,\ldots)$, 是 $\ell^q$ 上的左移; 此也包括 $p=1,q=\infty$.
\end{proof}
\begin{exercise}\label{ps03:3}
原题3. 对 $E\subset\R$, $x\in\R$, 令 $E+x=\{y+x:y\in E\}$, 证明 $m^*(E+x)=m^*(E)$.
\end{exercise}
\begin{proof}[AI 独立解答]
任意覆盖 $E$ 的开区间列 $(I_n)$ 平移为覆盖 $E+x$ 的开区间列 $(I_n+x)$, 长度总和不变. 对覆盖取下确界得 $m^*(E+x)\le m^*(E)$. 再以 $-x$ 平移 $E+x$ 得反向界. 两边为无穷时论证仍成立, 不涉及无穷相减.
\end{proof}
\begin{exercise}\label{ps03:4}
原题4. 设 $U\subset\R$ 为非空开集, 对 $x\in U$ 令 $a_x=\inf\{a\in\R:(a,x]\subset U\}$, $b_x=\sup\{b\in\R:[x,b)\subset U\}$. (a) 证明 $(a_x,b_x)\subset U$. (b) 若 $y\in(a_x,b_x)$, 证明 $(a_y,b_y)=(a_x,b_x)$. (c) 证明 $U=\bigcup_{q\in U\cap\Q}(a_q,b_q)$. 原提示: $\Q$ 在 $\R$ 中稠密. 端点允许取 $\pm\infty$; 这是原定义在无界开集上的必要解释.
\end{exercise}
\begin{proof}[AI 独立解答]
(a) 开性给出 $a_x<x<b_x$. 若 $a_x<t<x$, 下确界定义给出某个 $a<t$ 使 $(a,x]\subset U$, 所以 $t\in U$; 右侧同理, $x$ 自身也在 $U$.

(b) 区间 $I_x=(a_x,b_x)$ 是包含 $x$ 且包含于 $U$ 的最大开区间: 若另一个这样的区间左端更小, 其与 $I_x$ 的并仍为 $U$ 内包含 $x$ 的开区间, 与 $a_x$ 的定义矛盾; 右端同理. 对 $y\in I_x$, $I_y$ 与 $I_x$ 相交, 二者的并为 $U$ 内开区间, 含 $x$ 和 $y$. 分别用两个最大性得 $I_x\cup I_y=I_x=I_y$.

(c) 每个 $I_q\subset U$. 反之, 对 $x\in U$ 在 $I_x$ 内选有理数 $q$, 由(b)得 $I_q=I_x$, 所以 $x$ 属于右端. 指标集可数, 给出可数个开区间的并; 重复区间可去重, 不同区间由(b)必不交.
\end{proof}

\originalsection{作业4}\label{sec:ps04}
原件2页, 数学题面1页. 本组使用第\ref{ch:measure}章已明示的 Lebesgue 测度构造事实; 原题明确允许使用有限开区间并可测, 并说明相关第8讲移至第5周.
\begin{exercise}\label{ps04:1}
原题1. 对任意函数 $f:\R\to\R$, 证明 $\mathcal A=\{E\subset\R:f^{-1}(E)\text{为 Lebesgue 可测集}\}$ 是 $\sigma$ 代数.
\end{exercise}
\begin{proof}[AI 独立解答]
$f^{-1}(\R)=\R$ 可测. 原像满足 $f^{-1}(E^c)=(f^{-1}E)^c$, $f^{-1}(\bigcup E_n)=\bigcup f^{-1}(E_n)$. Lebesgue 可测集对补集和可数并封闭, 所以 $\mathcal A$ 具备全部 $\sigma$ 代数公理. 不要求 $f$ 本身可测.
\end{proof}
\begin{exercise}\label{ps04:2}
原题2. 设 $E\subset\R$, $m^*(E)<\infty$. 证明 $E$ 可测当且仅当每个 $\eps>0$ 都存在有限个开区间的并 $U$, 使 $m^*(U\triangle E)<\eps$. 这是 Littlewood 第一原则. 原提示: 逆向证明 $m^*(A\cap E)+m^*(A\cap E^c)\le m^*(A)+\eps$ 对任意 $A\subset\R$ 成立.
\end{exercise}
\begin{proof}[AI 独立解答]
正向: 由外测度定义取开区间覆盖 $E\subset O=\bigcup_{j\ge1}I_j$, 使 $\sum|I_j|<m(E)+\eps/2$. $O$ 可测且有限测度, $m(O\setminus E)=m(O)-m(E)<\eps/2$. 取 $N$ 使 $\sum_{j>N}|I_j|<\eps/2$, 令 $U=\bigcup_{j\le N}I_j$. 则 $U\setminus E\subset O\setminus E$, $E\setminus U\subset\bigcup_{j>N}I_j$, 故对称差的外测度小于 $\eps$.

逆向: 固定 $A$. 若 $m^*(A)=\infty$, Carathéodory 的所需方向自动成立. 否则取 $m^*(U\triangle E)<\delta$. 因 $U$ 可测,
\begin{align*}
m^*(A\cap E)&\le m^*(A\cap U)+\delta,\\
m^*(A\setminus E)&\le m^*(A\setminus U)+\delta,\\
m^*(A\cap E)+m^*(A\setminus E)&\le m^*(A)+2\delta.
\end{align*}
令 $\delta\downarrow0$ 得分割不等式; 另一方向由外测度的次可加性得到. 故 $E$ 满足可测性定义. 取 $\delta=\eps/2$ 即为原提示的不等式.
\end{proof}
\begin{exercise}\label{ps04:3}
原题3. 设 $E$ 可测. (a) 对每个 $x\in\R$, 证明 $E+x$ 可测. (b) 对每个 $r>0$, 证明 $rE=\{ry:y\in E\}$ 可测.
\end{exercise}
\begin{proof}[AI 独立解答]
(a) 对任意 $A\subset\R$, 将 $A-x$ 由 $E$ 分割, 再用题\ref{ps03:3}的平移不变性, 得
$m^*(A)=m^*(A\cap(E+x))+m^*(A\setminus(E+x))$, 故可测.

(b) 对区间覆盖缩放, 再反向用 $r^{-1}$, 得 $m^*(rB)=r\,m^*(B)$ 对任意 $B$ 成立. 将 $r^{-1}A$ 由 $E$ 分割后乘 $r$, 得同一可测性等式. $r>0$ 使扩展非负实数的乘法不出现未定义运算.
\end{proof}

\originalsection{作业5}\label{sec:ps05}
原件3页, 数学题面2页.
\begin{exercise}\label{ps05:1}
原题1. (a) 若 $E,F$ 可测, 证明 $m(E\cup F)+m(E\cap F)=m(E)+m(F)$. (b) 若 $E_1\supset E_2\supset\cdots$ 可测且 $m(E_1)<\infty$, 证明 $m(\bigcap_{k\ge1}E_k)=\lim_n m(E_n)$. 原件将(b)称为“从下连续”; 此处条件为递减列, 应称\textbf{从上连续}.
\end{exercise}
\begin{proof}[AI 独立解答]
(a) 将 $E\cup F$ 分为 $E\setminus F$, $E\cap F$, $F\setminus E$ 三个不交部分, 对非负扩展实数作加法即可得到等式. 不用可能未定义的 $\infty-\infty$.

(b) $A_n=E_1\setminus E_n$ 递增, 其并为 $E_1\setminus\bigcap E_n$. 从下连续由不交差集分解和可数可加性得到, 已在定理\ref{thm:mct}的证明中给出. 因 $m(E_1)$ 有限, 可用 $m(E_n)=m(E_1)-m(A_n)$ 并取极限. 有限性必要: $E_n=[n,\infty)$ 的每项测度为无穷而交集为空.
\end{proof}
\begin{exercise}\label{ps05:2}
原题2. 设 $E$ 可测, $f,g:E\to[-\infty,\infty]$ 可测. (a) 证明 $fg$ 可测, 其中异号无穷相乘取 $-\infty$, 同号取 $+\infty$, $(\pm\infty)0=0$. (b) 对 $a\in\R$, 在 $f=-g=\pm\infty$ 处令 $h=a$, 其余处令 $h=f+g$, 证明 $h$ 可测.
\end{exercise}
\begin{proof}[AI 独立解答]
(a) 令 $f_n=\max(-n,\min(f,n))$, $g_n$ 同理. 它们有限且可测, 所以 $f_ng_n$ 可测. 对有限值其极限为通常的积; 一个因子为无穷、另一个为非零有限值时极限有正确符号; 两个无穷时同样成立; 某因子为零时乘积恒为零. 因而 $f_ng_n\to fg$ 于扩展实数, 可测函数的逐点极限可测.

(b) 集合 $A=\{f=+\infty,g=-\infty\}\cup\{f=-\infty,g=+\infty\}$ 可测, 因无穷取值集由可数阈值交表示. 在 $E\setminus A$ 上 $f_n+g_n\to f+g$, 包括两个同号无穷的情形. 在 $A$ 上将每个 $f_n+g_n$ 的值换成常数 $a$, 得可测分段函数, 其逐点极限就是 $h$. 这避免在 $A$ 上使用未定义的无穷相减.
\end{proof}
\begin{exercise}\label{ps05:3}
原题3. (a) 若 $E,F$ 可测, $f:E\cup F\to[-\infty,\infty]$, 证明 $f$ 可测当且仅当两限制 $f|_E,f|_F$ 可测. (b) 若 $E$ 可测, 将 $f:E\to[-\infty,\infty]$ 在 $E^c$ 上延为零得 $g$, 证明 $f$ 可测当且仅当 $g$ 可测. (c) 若 $u,v:E\to\R$ 可测, 证明 $(u^2+v^2)^{1/2}$ 可测.
\end{exercise}
\begin{proof}[AI 独立解答]
(a) 限制的阈值集为原阈值集与可测定义域的交; 反向将两限制的阈值集取并, 恰为 $\{f>t\}$, 故可测. 这里定义域上“可测”指阈值集作为 $\R$ 的子集可测.

(b) $g$ 可测蕴含限制可测. 反向若 $t\ge0$, $\{g>t\}=\{x\in E:f(x)>t\}$; 若 $t<0$, 再并上 $E^c$, 仍可测.

(c) 可测有限函数乘积及和可测, 连续函数 $t\mapsto\sqrt t$ 在非负半轴上的复合保持可测. 后一事实可由连续函数的开集原像为可数开区间并来验证, 每个区间原像由两个阈值集相交表示.
\end{proof}
\begin{exercise}\label{ps05:4}
原题4. 设 $E$ 可测, $m(E)<\infty$, $f_n,f:E\to\R$ 可测且 $f_n\to f$ 几乎处处. (a) 对 $k,n\in\N$, 令 $F_n(k)=\bigcup_{m\ge n}\{|f_m-f|\ge1/k\}$, 证明 $m(F_n(k))\to0$. (b) 给定 $\eps>0$, 选 $n_k$ 使 $m(F_{n_k}(k))<2^{-k}\eps$, 令 $F=\bigcup_{k\ge1}F_{n_k}(k)$, 证明 $m(F)<\eps$ 且 $f_n\to f$ 在 $E\setminus F$ 上一致. 这是 Littlewood 第二原则, 亦称 Egorov 定理. 原件的 $F^c$ 在此解释为定义域内的相对补集.
\end{exercise}
\begin{proof}[AI 独立解答]
(a) 固定 $k$, $F_n(k)$ 可测且递减. 在收敛点, 充分大的 $m$ 都有 $|f_m-f|<1/k$, 所以该点不在交集中. 交集因此包含于收敛失败的零测集, 测度为零. 所有 $F_n(k)\subset E$ 有有限测度, 用题\ref{ps05:1}(b)得结论.

(b) 次可加性给出 $m(F)\le\sum m(F_{n_k}(k))<\sum2^{-k}\eps=\eps$; 严格性例如由第一项的严格差额保证. 对每个 $x\in E\setminus F$, 任意 $n\ge n_k$ 都有 $|f_n(x)-f(x)|<1/k$. 给定误差 $\delta>0$, 选 $1/k<\delta$, 同一个 $n_k$ 对所有这些 $x$ 有效, 即一致收敛.
\end{proof}

\originalsection{作业6}\label{sec:ps06}
原件3页, 数学题面2页. 区间均有 $a<b$. 阶梯函数指有限值简单函数且各取值集为有限个区间的并; 区间可开、闭或半开, 单点也允许作为退化区间. 有限端点的值不影响积分.
\begin{exercise}\label{ps06:1}
原题1. (a) 若实阶梯函数 $\psi:[a,b]\to\R$ 满足 $\sup|\psi|\le B$, 对每个 $\eps>0$, 构造连续 $g$ 与可测 $E\subset[a,b]$, 使 $g(a)=g(b)=0$, $\sup|g|\le B$, $m(E)<\eps$, 且在 $[a,b]\setminus E$ 上 $|\psi-g|<\eps$.

(b) 若实简单函数 $\phi:[a,b]\to\R$ 满足 $\sup|\phi|\le B$, 对每个 $\eps>0$, 构造阶梯 $\psi$ 与可测 $E$, 使 $\sup|\psi|\le B$, $m(E)<\eps$, 且在补集上 $|\phi-\psi|<\eps$. 原提示: 用 Littlewood 第一原则.
\end{exercise}
\begin{proof}[AI 独立解答]
(a) 收集所有取值区间端点及 $a,b$ 为有限集 $D$. 在 $[a,b]\setminus D$ 的每个分割区间上 $\psi$ 恒定. 在各端点选互不相交的小邻域, 相对总长度小于 $\eps$, 令它们的并为 $E$, 并包含全部端点. 在 $E$ 外令 $g=\psi$; 在内部端点的小邻域中用线段连接两侧常值, 在 $a,b$ 邻域连接零与邻近常值. 得连续 $g$, 所有线段值是 $[-B,B]$ 中两点的凸组合, 故不超 $B$. 补集上误差为零. $B=0$ 时直接取 $g=0,E=\varnothing$.

(b) 写 $\phi=\sum_{j=1}^m c_j\ind_{A_j}$, $A_j$ 为可测分割, $|c_j|\le B$. 按题\ref{ps04:2}选有限区间并 $U_j$, 使 $m(A_j\triangle U_j)<\eps/m$. 与 $[a,b]$ 相交, 定义 $D_1=U_1$, $D_j=U_j\setminus\bigcup_{i<j}U_i$, 并在所有 $U_j$ 的补集取零. 这些集合由有限端点分割而成, 故 $\psi=\sum c_j\ind_{D_j}$ 为阶梯函数, 且各点至多取一个 $c_j$, 保持 $B$ 界. 除去 $E=\bigcup(A_j\triangle U_j)$, 各点属于恰好正确的 $U_j$, 从而 $\psi=\phi$. 总例外测度小于 $\eps$.
\end{proof}
\begin{exercise}\label{ps06:2}
原题2. 若实可测 $f:[a,b]\to\R$ 满足 $\sup|f|\le B$, 对每个 $\eps>0$, 构造连续 $g$ 和可测 $E$, 使 $g(a)=g(b)=0$, $\sup|g|\le B$, $m(E)<\eps$, 且补集上 $|f-g|<\eps$. 这是 Littlewood 第三原则的有界版本.
\end{exercise}
\begin{proof}[AI 独立解答]
将 $[-B,B]$ 分成有限个长度小于 $\eps/3$ 的值域段, 以各段内一点代替 $f$ 得简单 $\phi$, 满足 $|\phi|\le B$, $\sup|f-\phi|<\eps/3$. 对题\ref{ps06:1}(b),(a)依次使用参数 $\eps/3$, 得阶梯 $\psi$、连续 $g$ 及两个例外集, 总测度小于 $2\eps/3<\eps$. 例外并集外三项误差各小于 $\eps/3$, 故总误差小于 $\eps$; 所有逼近均保持 $B$ 界与最终端点零. $B=0$ 直接取零函数.
\end{proof}
\begin{exercise}\label{ps06:3}
原题3. 设实函数 $f:[a,b]\to\R$ Lebesgue 可积. 对每个 $\eps>0$, (a) 求有界可测 $h$ 使 $\int_a^b|f-h|<\eps$; (b) 求阶梯 $\psi$ 使 $\int_a^b|f-\psi|<\eps$; (c) 求连续 $g$, $g(a)=g(b)=0$, 使 $\int_a^b|f-g|<\eps$. 原提示(a): 取 $h_n=f\ind_{\{|f|\le n\}}$.
\end{exercise}
\begin{proof}[AI 独立解答]
(a) $|f-h_n|=|f|\ind_{\{|f|>n\}}\to0$ 几乎处处, 受可积 $|f|$ 控制. 定理\ref{thm:dct}给出积分趋于零, 取充分大的 $n$.

(b) 先取(a)的 $h$, 误差小于 $\eps/3$, 令 $|h|\le B$. 对有界函数的值域作简单量化后用题\ref{ps06:1}(b), 得阶梯 $\psi$ 保持 $B$ 界, 在某个测度小于 $\delta$ 的例外集外误差小于 $2\delta$. 因而 $\int|h-\psi|\le2\delta(b-a)+2B\delta$. 选 $\delta$ 使此值小于 $\eps/3$, 与截断误差相加得到要求.

(c) 先取(b)中的 $\psi$, 误差小于 $\eps/2$, 设 $|\psi|\le B_1$. 用题\ref{ps06:1}(a)的构造, $g=\psi$ 于例外集外且 $|g|\le B_1$, 所以 $\int|\psi-g|\le2B_1m(E)$. 选例外测度使此项小于 $\eps/2$; $B_1=0$ 时取零. 相加即得结论. 这些构造细化了命题\ref{prop:lp-density}的 $p=1$ 密度证明.
\end{proof}
\begin{exercise}\label{ps06:4}
原题4. 若复函数 $f\in L^1([-\pi,\pi])$, 定义 $\widehat f(n)=(2\pi)^{-1}\int_{-\pi}^{\pi}f(x)\ee^{-\ii nx}\dd x$, $n\in\Z$. 证明 Riemann--Lebesgue 引理 $|\widehat f(n)|\to0$ 当 $|n|\to\infty$. 原提示: 先对阶梯函数证明, 再用题3逼近.
\end{exercise}
\begin{proof}[AI 独立解答]
对单个区间 $(c,d)$ 及 $n\ne0$, $\left|\int_c^d\ee^{-\ii nx}\dd x\right|\le2/|n|$, 因而有限线性组合的阶梯函数 Fourier 系数趋于零. 将 $f$ 的实、虚部分别按题\ref{ps06:3}(b)逼近, 得复阶梯 $\psi$ 使 $\norm{f-\psi}_1$ 任意小. 有 $|\widehat f(n)-\widehat\psi(n)|\le(2\pi)^{-1}\norm{f-\psi}_1$, 对所有 $n$ 一致. 先固定逼近, 再令 $|n|\to\infty$, 最后令逼近误差趋于零, 即得结论.
\end{proof}

\originalsection{作业7}\label{sec:ps07}
原件2页, 数学题面1页. 原题允许使用: 对 $a<b$, $1\le p<\infty$, $f\in L^p([a,b])$ 和每个 $\eps>0$, 存在阶梯 $\psi$ 及连续 $g$, $g(a)=g(b)=0$, 使 $\norm{f-\psi}_p+\norm{f-g}_p<\eps$. 此结论亦由命题\ref{prop:lp-density}的证明给出.
\begin{exercise}\label{ps07:1}
原题1. 设 $1\le p<\infty$. (a) 证明 $1\le q\le p$ 时 $L^p([a,b])\subset L^q([a,b])$. (b) 证明 $L^p([a,b])$ 可分. (c) 对 $f\in L^p(\R)$ 与 $\eps>0$, 求连续紧支集函数 $g$ 使 $\norm{f-g}_p<\eps$. (d) 证明 $L^p(\R)$ 可分. 原提示(b): 参看当周第二讲的证明提纲; (c): 先截断至 $[-n,n]$, 再用上述密度定理.
\end{exercise}
\begin{proof}[AI 独立解答]
(a) $q=p$ 直接成立. 若 $q<p$, 对 $|f|^q$ 与1用共轭指数 $p/q$, $p/(p-q)$ 的 Hölder, 得 $\norm f_q\le(b-a)^{1/q-1/p}\norm f_p$. 此复用第\ref{ch:measure}章有限测度包含关系的习题解答.

(b) 取端点为有理数（另允许固定 $a,b$）、系数为 $\Q+\ii\Q$ 的有限区间阶梯函数, 形成可数集. 任意阶梯函数可用这些函数逼近: 端点移动 $\delta$ 只改变测度至多 $2\delta$ 的集合, 对区间指示函数的 $L^p$ 误差为对称差测度的 $1/p$ 次方; 系数移动的误差由有限区间长度控制. 与原题允许的阶梯密度结合, 该可数集稠密.

(c) $\int_{|x|>n}|f|^p\to0$, 所以先选 $n$ 使截断误差小于 $\eps/2$. 在 $[-n,n]$ 内用零端点连续逼近, 误差小于 $\eps/2$, 再延为区间外零. 端点零保证全轴连续, 支集包含于 $[-n,n]$, 与截断误差相加小于 $\eps$.

(d) 取有限个有理端点的有界区间指示函数与有理复系数的线性组合. 此集可数, 对 $[-n,n]$ 上的截断用(b)逼近, 再令 $n$ 足够大控制全轴尾积分, 得稠密.
\end{proof}
\begin{exercise}\label{ps07:2}
原题2. (a) 对可测 $E\subset\R$, 证明 $L^\infty(E)$ 为 Banach 空间. (b) 对 $f\in C([a,b])$, 证明 $\norm f_{L^\infty}=\sup_{[a,b]}|f|$, 并推出 $C([a,b])$ 不在 $L^\infty([a,b])$ 中稠密.
\end{exercise}
\begin{proof}[AI 独立解答]
(a) 等价类上的范数正定性、齐次性和三角不等式见第\ref{ch:measure}章定义及其后论证. 设 $(f_n)$ 为 Cauchy 列, 选择可测代表元. 对每对 $n,m$, 本质界 $|f_n-f_m|\le\norm{f_n-f_m}_\infty$ 在一个零测集外成立. 再除去每个 $f_n$ 本质界的例外, 可数并仍为零测集 $N$. 在 $E\setminus N$ 上 Cauchy 条件给出一致 Cauchy, 因而逐点有有限极限 $f$, 在 $N$ 上置零. 此函数可测. 固定 $n_0$ 足够大, 在共同零测集外有 $|f|\le\norm{f_{n_0}}_\infty+1$, 所以 $f$ 本质有界. Cauchy 界取点态极限得 $\norm{f_n-f}_\infty\to0$. 空或零测定义域得到零 Banach 空间, 同样成立.

(b) 令 $M=\max|f|$, 本质上确界至多 $M$. 若 $C<M$, 连续性使 $|f|>C$ 在某个正长度相对区间内成立, 故 $C$ 不是本质界; 于是二者相等. 再取 $h=\ind_{[a,c]}$, $a<c<b$. 若连续 $g$ 满足 $\norm{g-h}_{L^\infty}<1/2$, 在左侧由连续性及几乎处处界得 $\operatorname{Re}g>1/2$, 在右侧得 $\operatorname{Re}g<1/2$, 且实际各有统一正差额. 在 $c$ 两侧取极限矛盾. 因而该 $h$ 到连续函数的距离至少 $1/2$, 不稠密.
\end{proof}
\begin{exercise}\label{ps07:3}
原题3. 对 $a<b$, $1\le p\le\infty$, $g\in L^\infty([a,b])$, 定义 $Tf=gf$. 证明 $T\in\B(L^p([a,b]))$ 且 $\norm T=\norm g_\infty$.
\end{exercise}
\begin{proof}[AI 独立解答]
乘法对等价类良定义且线性, 几乎处处界给出 $\norm{gf}_p\le\norm g_\infty\norm f_p$, 包括 $p=\infty$. 令 $M=\norm g_\infty$. $M=0$ 时 $T=0$. $M>0$ 时每个 $0<\delta<M$ 的集合 $A=\{|g|>M-\delta\}$ 有正测度, 否则 $M-\delta$ 是本质界. 取 $f=\ind_A/m(A)^{1/p}$ 于有限 $p$, 或取 $f=\ind_A$ 于 $p=\infty$. 它们范数为1, 且 $\norm{Tf}_p\ge M-\delta$. 令 $\delta\downarrow0$ 得反向界.
\end{proof}
\begin{exercise}\label{ps07:4}
原题4. (a) 在预 Hilbert 空间中证明极化恒等式
\[
\ip uv=\tfrac14\bigl(\norm{u+v}^2-\norm{u-v}^2+\ii\norm{u+\ii v}^2-\ii\norm{u-\ii v}^2\bigr).
\]
(b) 若复赋范空间的范数满足平行四边形恒等式, 证明存在 Hermitian 内积使 $\norm u=\ip uu^{1/2}$. 原提示: 用(a)右端定义内积.
\end{exercise}
\begin{proof}[AI 独立解答]
(a) 记 $z=\ip uv$. 前两个平方之差为 $4\operatorname{Re}z$, 后两个之差为 $4\operatorname{Im}z$, 因第一变量线性而 $\ip u{\ii v}=-\ii z$. 代回即得恒等式.

(b) 复用命题\ref{prop:parallelogram}的完整逆向证明. 具体地, 在底层实空间设 $B(u,v)=(\norm{u+v}^2-\norm{u-v}^2)/4$. 平行四边形恒等式相减给出 $B(x+z,y)+B(x-z,y)=2B(x,y)$. 先取 $z=x$ 得二倍齐次性, 再令 $x=(u+v)/2,z=(u-v)/2$ 得加法性; 整数、有理齐次性及连续性给出实双线性. 对称性与 $B(u,u)=\norm u^2$ 直接成立. 由 $\norm{\ii u}=\norm u$ 有 $B(\ii u,\ii v)=B(u,v)$, $B(\ii u,v)=-B(u,\ii v)$. 因而 $B(u,v)+\ii B(u,\ii v)$ 对第一变量复线性、共轭对称, 且对角值为 $\norm u^2$. 这建立了内积, 没有预先假设待证的双线性或使用其 Cauchy--Schwarz.
\end{proof}

\originalsection{作业8}\label{sec:ps08}
原件3页, 数学题面2页. 此组按原作业引入周期 Sobolev 空间, 只讨论 Fourier 加权定义及其指定性质.
\begin{exercise}\label{ps08:1}
原题1. 设 $H$ 为 Hilbert 空间, $(e_n)_{n\ge1}$ 为可数无限正交归一系, $W=\Span\{e_n:n\ge1\}$ 指有限线性组合. (a) 证明 $w\in\overline W$ 当且仅当 $w=\sum_{n\ge1}c_ne_n$ 对某个 $c\in\ell^2$ 成立. (b) 对 $u\in H$ 和 $w\in\overline W$, 证明
\[
\left\|u-\sum_{n\ge1}\ip u{e_n}e_n\right\|\le\norm{u-w},
\]
且等号当且仅当 $w=\sum\ip u{e_n}e_n$. 原提示: 展开平方并配方. 原PDF在(b)左端漏印范数符号; 向量不能直接与实数比较, 此处依题意补回.
\end{exercise}
\begin{proof}[AI 独立解答]
(a) 若 $c\in\ell^2$, 正交性给出级数部分和之差的范数平方为平方尾和, 完备性使它收敛, 极限在 $\overline W$ 内. 反之, Bessel 使 $c_n=\ip w{e_n}\in\ell^2$; 级数和 $v$ 满足 $w-v$ 与每个 $e_n$ 正交, 从而与 $\overline W$ 正交. 因 $w-v\in\overline W$, 得 $w-v=0$. 此复用定理\ref{thm:onb}中针对正交系的展开证明.

(b) 令 $v=\sum\ip u{e_n}e_n$. 同一证明给出 $u-v\perp\overline W$. 对任意 $w\in\overline W$, 两项 $u-v$ 与 $v-w$ 正交, 所以 $\norm{u-w}^2=\norm{u-v}^2+\norm{v-w}^2$. 不等式与等号条件同时得到; 也扩展了第\ref{ch:hilbert}章有限投影的习题.
\end{proof}
\begin{exercise}\label{ps08:2}
原题2. 设 $W$ 为 Hilbert 空间 $H$ 的线性子空间. (a) 证明 $W^\perp=\{u:\ip uw=0\text{对所有 }w\in W\}$ 是闭线性子空间. (b) 证明 $\overline W=(W^\perp)^\perp$. 原PDF有闭包横线; 未假设 $W$ 本身闭.
\end{exercise}
\begin{proof}[AI 独立解答]
(a) 每个固定 $w$ 的内积泛函连续线性, 所以 $W^\perp$ 是这些泛函的闭核的交. (b) 连续性给出 $W^\perp=(\overline W)^\perp$ 且 $\overline W\subset(W^\perp)^\perp$. 若 $x\in(W^\perp)^\perp$, 对闭子空间 $\overline W$ 用正交分解\ref{cor:orthodecompose}, 得 $x=v+z$, $v\in\overline W,z\in W^\perp$. 则 $0=\ip xz=\norm z^2$, 故 $x=v\in\overline W$.
\end{proof}
\begin{exercise}\label{ps08:3}
原题3. 对 $s\ge0$, 定义 $H^s(\mathbb T)=\{f\in L^2([-\pi,\pi]):\sum_{n\in\Z}(1+|n|^2)^s|\widehat f(n)|^2<\infty\}$. (a) 若 $f\in C^k([-\pi,\pi])$ 且 $f^{(j)}(\pi)=f^{(j)}(-\pi)$ 对 $0\le j<k$ 成立, 证明 $f\in H^s(\mathbb T)$ 对 $0\le s\le k$ 成立. (b) 证明 $H^s$ 是向量空间, 且 $\ip fg_{H^s}=\sum_{n\in\Z}\widehat f(n)\overline{\widehat g(n)}(1+|n|^2)^s$ 为 Hermitian 内积. (c) 证明它是 Hilbert 空间. 原提示(a): 分部积分.
\end{exercise}
\begin{proof}[AI 独立解答]
(a) $k=0$ 时由 Parseval 得 $H^0=L^2$. $k\ge1$ 时对 $n\ne0$ 分部积分 $k$ 次, 端点项因指定的周期条件消失, 得 $\widehat{f^{(k)}}(n)=(\ii n)^k\widehat f(n)$. $f^{(k)}$ 连续且在有限区间平方可积, Parseval 给出 $\sum_{n\ne0}|n|^{2k}|\widehat f(n)|^2<\infty$. 对 $s\le k$, $n\ne0$ 时 $(1+n^2)^s\le2^s|n|^{2k}$; 零频项有限, 得结论. 仅写 $O(|n|^{-k})$ 会不足以证明端点 $s=k$, 此处使用导数的平方可和.

(b) 加权系数列相加、数乘仍在 $\ell^2(\Z)$, 因而为向量空间. Cauchy--Schwarz 保证所给级数绝对收敛, 线性与共轭对称逐项成立. 若对角值零, 全部 Fourier 系数零, 由定理\ref{thm:fouriercomplete}的完备性得 $f=0$ 于 $L^2$ 等价类, 所以正定.

(c) 映射 $U_sf=((1+n^2)^{s/2}\widehat f(n))_{n\in\Z}$ 线性等距至 $\ell^2(\Z)$. 对任意 $c\in\ell^2(\Z)$, 令 $a_n=(1+n^2)^{-s/2}c_n$. 因权重至少1, $a\in\ell^2$. Fourier 完备性保证 $\sum a_n\ee^{\ii nx}$ 在 $L^2$ 中收敛于一个 $f$, 其 Fourier 系数恰为 $a_n$; 正交归一基为 $(2\pi)^{-1/2}\ee^{\ii nx}$, 故部分和的 $L^2$ 范数平方是 $2\pi\sum|a_n|^2$. 于是 $f\in H^s$, $U_sf=c$, 映射满射. $\ell^2$ 完备给出 $H^s$ 完备.
\end{proof}
\begin{exercise}\label{ps08:4}
原题4. 设 $s>1/2$. (a) 证明存在 $C(s)>0$ 使 $\sum_{n\in\Z}|\widehat f(n)|\le C(s)\norm f_{H^s}$ 对所有 $f\in H^s$ 成立. (b) 证明每个 $f\in H^s$ 有连续代表元 $g\in C([-\pi,\pi])$, $f=g$ 几乎处处, 且 $\norm g_\infty\le C(s)\norm f_{H^s}$. 原提示: Weierstrass 一致收敛判别.
\end{exercise}
\begin{proof}[AI 独立解答]
(a) 用加权 Cauchy--Schwarz 得所需界, 其中 $C(s)=(\sum_{n\in\Z}(1+n^2)^{-s})^{1/2}<\infty$, 因 $2s>1$ 且非零项可由 $|n|^{-2s}$ 控制.

(b) 级数 $\sum\widehat f(n)\ee^{\ii nx}$ 由(a)绝对一致收敛, 其和 $g$ 连续且周期, 上确界受系数绝对和控制. 同一部分和在 $L^2$ 中趋于 $f$, 一致收敛也蕴含 $L^2$ 收敛于 $g$, 极限唯一给出 $f=g$ 几乎处处. 所指是唯一连续代表元, 不是任意等价类代表元都连续.
\end{proof}

\originalsection{作业9}\label{sec:ps09}
原件4页, 数学题面及前置定理3页. 原题先给出如下子列定理及证明: 若 $h_n\in L^1(E)$ 且 $\norm{h_n}_1\to0$, 可选 $n_j$ 递增使 $m(\{|h_{n_j}|\ge2^{-j}\})\le2^{-j}$, 因为该测度不超过 $2^j\norm{h_{n_j}}_1$. 这些例外集的尾并测度至多 $\sum_{j\ge k}2^{-j}$, 所以其上极限集合零测; 在补集上 $h_{n_j}\to0$. 原证明中 $h/f$ 名称混写及阈值 $\ge/>$ 混用, 此处统一为 $h$ 及 $\ge$.

原题还固定光滑函数 $\phi:\R\to\R$, 满足 $\phi(-x)=\phi(x)$, $\phi\ge0$, $|x|>1$ 时为零, $\int_\R\phi=1$. 定义磨光核 $\phi_\eps(x)=\eps^{-1}\phi(x/\eps)$, $\eps>0$, 故支集在 $[-\eps,\eps]$, 积分仍为1. 后续交换积分使用第\ref{ch:measure}章明示的 Tonelli--Fubini 先修, 并核对非负或绝对可积条件.
\begin{exercise}\label{ps09:1}
原题1. 设 $g\in C(\R)$, $|x|>R$ 时 $g(x)=0$, 定义 $g_\eps(x)=\int_\R\phi_\eps(x-y)g(y)\dd y=\int_\R\phi_\eps(z)g(x-z)\dd z$. (a) 证明 $g_\eps$ 无限可微且 $|x|>R+\eps$ 时为零. (b) 证明 $\norm{g_\eps-g}_\infty\to0$ 当 $\eps\downarrow0$. (c) 对 $f\in L^2(\R)$ 和 $\delta>0$, 求光滑紧支集 $h$ 使 $\norm{f-h}_2<\delta$. 原提示: Fourier 平均收敛的思路可用于前两题.
\end{exercise}
\begin{proof}[AI 独立解答]
(a) 积分的 $y$ 限于 $[-R,R]$. 对任何阶导数, $\phi_\eps^{(m)}$ 连续有界, 其差商由再高一阶导数的上界控制. 控制收敛允许逐阶微分, 得 $g_\eps^{(m)}(x)=\int\phi_\eps^{(m)}(x-y)g(y)\dd y$. 若 $|x|>R+\eps$, $|y|\le R$ 时核为零, 所以函数为零.

(b) $g$ 在全轴上一致连续, 因紧区间外为零. 令其模连续性为 $\omega_g(\eps)=\sup_{|x-y|\le\eps}|g(x)-g(y)|\to0$. 核非负且质量1, 得 $|g_\eps(x)-g(x)|\le\omega_g(\eps)$ 对所有 $x$ 成立.

(c) 题\ref{ps07:1}(c)先给出连续紧支集 $g$, $\norm{f-g}_2<\delta/2$. 对 $0<\eps<1$, $g_\eps-g$ 支集在共同有限区间 $[-R-1,R+1]$, 故 $\norm{g_\eps-g}_2\le\sqrt{2R+2}\norm{g_\eps-g}_\infty\to0$. 取小 $\eps$, 令 $h=g_\eps$, 用(a)及三角不等式即得.
\end{proof}
\begin{exercise}\label{ps09:2}
原题2. 对 $f\in L^2(\R)$, $\eps>0$, 定义 $f_\eps(x)=\int_\R\phi_\eps(x-y)f(y)\dd y$. (a) 证明 $\norm{f_\eps}_2\le\norm f_2$. (b) 证明 $\norm{f_\eps-f}_2\to0$ 当 $\eps\downarrow0$. 原提示(a): 先对连续函数估计, 再作 $L^2$ 逼近, 用前置子列定理识别极限.
\end{exercise}
\begin{proof}[AI 独立解答]
(a) 每个 $x$ 的积分绝对收敛, 因 $\int|\phi_\eps(x-y)f(y)|\dd y\le\norm{\phi_\eps}_2\norm f_2$. 用质量1的非负权重作 Cauchy--Schwarz,
\[
|f_\eps(x)|^2\le\int_\R\phi_\eps(z)|f(x-z)|^2\dd z.
\]
该非负被积函数可测, Tonelli 与平移不变性得 $\int|f_\eps|^2\le\int\phi_\eps(z)\norm f_2^2\dd z=\norm f_2^2$. 定义的卷积可测也由积分的可测性给出, 故属于 $L^2$. 这一直接证明比原提示的极限识别少一步, 不额外假设 $f$ 连续.

(b) 对给定 $\delta>0$, 选连续紧支集 $g$ 使 $\norm{f-g}_2<\delta$. (a)作用于 $f-g$ 给出
$\norm{f_\eps-f}_2\le2\delta+\norm{g_\eps-g}_2$. 最后一项由题\ref{ps09:1}的共同支集与一致收敛趋于零. 先取上极限再令 $\delta\downarrow0$.
\end{proof}
\begin{exercise}\label{ps09:3}
原题3. 若 $H$ 为 Hilbert 空间, $P\in\B(H,H)$ 是投影（$P^2=P$）且 $\ip{Pu}v=\ip u{Pv}$ 对所有 $u,v$ 成立, 证明 $W=\ran P$ 是闭子空间, 且 $P=\Pi_W$ 为正交投影.
\end{exercise}
\begin{proof}[AI 独立解答]
此为第\ref{ch:adjoint}章关于自伴幂等算子的习题. $\ran P=\ker(I-P)$, 故闭. 对 $w=Pv\in W$, $\ip{u-Pu}w=\ip{P(u-Pu)}v=0$. 又 $Pu\in W$, 所以 $u=Pu+(u-Pu)$ 为正交分解, 其唯一性给出 $Pu=\Pi_Wu$.
\end{proof}
\begin{exercise}\label{ps09:4}
原题4. 设 $C\subset H$ 为闭凸集. (a) 证明每个 $u\in H$ 有唯一最近点 $P_Cu\in C$, 即 $\norm{u-P_Cu}=\inf_{w\in C}\norm{u-w}$. (b) 对 $u\in H$ 和 $v\in C$, 证明 $v=P_Cu$ 当且仅当 $\operatorname{Re}\ip{u-v}{w-v}\le0$ 对所有 $w\in C$ 成立. 原提示(b): 展开 $\norm{u-[(1-t)v+tw]}^2$ 再令 $t\downarrow0$.

\textbf{条件勘误}: 原题未声明 $C$ 非空. $C=\varnothing$ 是闭凸集但无最近点, 所以(a)按字面为假. 以下在明确补充 $C\ne\varnothing$ 后证明; 不把此条件隐藏进原题.
\end{exercise}
\begin{proof}[AI 独立解答]
(a) 复用定理\ref{thm:projection}的闭凸集证明. 令 $d=\inf_C\norm{u-w}<\infty$, 选 $v_n\in C$ 使距离趋于 $d$. 中点仍在 $C$, 平行四边形恒等式给出
\[
\norm{v_n-v_m}^2\le2\norm{u-v_n}^2+2\norm{u-v_m}^2-4d^2\longrightarrow0.
\]
完备性及闭性使极限 $v\in C$, 连续性使距离为 $d$. 两个不同最近点的中点将由同式给出更小距离, 故唯一.

(b) 固定 $u\in H,v\in C$. 若 $v$ 最优, 对 $w\in C$, $0<t\le1$, 凸性与原提示给出 $0\le-2t\operatorname{Re}\ip{u-v}{w-v}+t^2\norm{w-v}^2$. 除以 $t$ 并令 $t\downarrow0$, 得不等式. 反向若该不等式成立, 展开 $\norm{u-w}^2=\norm{u-v}^2+\norm{w-v}^2-2\operatorname{Re}\ip{u-v}{w-v}\ge\norm{u-v}^2$, 所以 $v$ 是唯一最近点.
\end{proof}
\begin{exercise}\label{ps09:5}
原题5. 设 $(h_k)$ 为任意非负实数列. (a) 证明 $C=\{b\in\ell^2:|b_k|\le h_k\text{对所有 }k\}$ 是闭凸集. (b) 对 $a\in\ell^2$, 求 $P_Ca$ 的各坐标.
\end{exercise}
\begin{proof}[AI 独立解答]
(a) 各坐标泛函连续, 闭圆盘原像的交闭. 逐坐标三角不等式保证凸性, $0\in C$ 保证非空. 不要求 $(h_k)\in\ell^2$.

(b) 令 $b_k=a_k$ 当 $|a_k|\le h_k$, 否则令 $b_k=h_ka_k/|a_k|$. 因 $|b_k|\le|a_k|$, $b\in\ell^2$, 且 $b\in C$. 对圆盘内任意 $z$, 若 $|a_k|>h_k$, 反三角不等式给出 $|a_k-z|\ge|a_k|-h_k=|a_k-b_k|$; 若 $|a_k|\le h_k$, 右端为零. 将平方不等式求和得 $\norm{a-c}_2\ge\norm{a-b}_2$ 对所有 $c\in C$ 成立. 由题\ref{ps09:4}唯一性, $P_Ca=b$; $a_k=0$ 时第一种情形已定义为零.
\end{proof}

\originalsection{作业10}\label{sec:ps10}
原件4页, 数学题面及前置定理3页. 原题提供 Arzelà--Ascoli 定理: 闭区间上连续函数列若一致有界且等度连续, 则有在上确界范数中收敛的子列. 等度连续指每个 $\eps>0$ 有共同 $\delta>0$, 使 $|x-y|<\delta$ 时所有函数的差均小于 $\eps$. 以下按原题允许使用这一实分析定理.
\begin{exercise}\label{ps10:1}
原题1. 设 $f_n\in C^1([0,1])$, $B=\sup_n(\norm{f_n}_2+\norm{f_n'}_2)<\infty$. (a) 证明 $\norm{f_n}_\infty\le B$. (b) 证明 $|f_n(x)-f_n(y)|\le B|x-y|^{1/2}$. (c) 证明有子列在 $L^2([0,1])$ 中收敛. 原提示(a): 先证 $|f_n(x)|\le|f_n(y)|+\norm{f_n'}_2$, 再对 $y$ 积分.
\end{exercise}
\begin{proof}[AI 独立解答]
(a) 微积分基本定理和 Cauchy--Schwarz 给出原提示的不等式; 对 $y\in[0,1]$ 积分, $\norm{f_n}_1\le\norm{f_n}_2$, 故 $|f_n(x)|\le\norm{f_n}_2+\norm{f_n'}_2\le B$. (b) 在 $x,y$ 间直接用 Cauchy--Schwarz: $|\int_y^x f_n'|\le\norm{f_n'}_2|x-y|^{1/2}$. (c) (a),(b)给出一致有界、等度连续, Arzelà--Ascoli 抽取一致收敛子列; 区间长度为1, $L^2$ 误差不超过一致误差.
\end{proof}
\begin{exercise}\label{ps10:2}
原题2. 若 $K\in C([a,b]^2)$, 定义 $Tf(x)=\int_a^bK(x,y)f(y)\dd y$. 证明 $T\in\B(L^2([a,b]))$ 且每个有界序列的像都有 $L^2$ 收敛子列, 即 $T$ 紧.
\end{exercise}
\begin{proof}[AI 独立解答]
设 $\ell=b-a$. 逐点 Cauchy--Schwarz 给出 $|Tf(x)|\le\sqrt\ell\norm K_\infty\norm f_2$, 故 $\norm{Tf}_2\le\ell\norm K_\infty\norm f_2$, 线性及对等价类良定义由积分得到. 若 $\norm{f_n}_2\le M$, 则 $(Tf_n)$ 一致有界. 又
$|Tf_n(x)-Tf_n(x')|\le\sqrt\ell M\sup_y|K(x,y)-K(x',y)|$. 核在紧正方形上一致连续, 右端对所有 $n$ 一致趋于零. Arzelà--Ascoli 给出一致收敛子列, 从而在 $L^2$ 收敛. 此为定理\ref{thm:continuouskernelcompact}在一般有限区间上的同一结论, 这里按原题给出的定理另证.
\end{proof}
\begin{exercise}\label{ps10:3}
原题3. 对 $s>0$, $B\ge0$, 证明 $K=\{f\in H^s(\mathbb T):\norm f_{H^s}\le B\}$ 是 $L^2([-\pi,\pi])$ 的紧子集. 因而嵌入 $H^s(\mathbb T)\to L^2(\mathbb T)$ 为紧算子.
\end{exercise}
\begin{proof}[AI 独立解答]
Parseval 给出 $\norm f_2^2\le2\pi B^2$, 且若 $P_N$ 保留 $|n|\le N$ 的 Fourier 项,
\[
\sup_{f\in K}\norm{f-P_Nf}_2^2\le2\pi B^2(1+(N+1)^2)^{-s}\longrightarrow0.
\]
再证闭性: 若 $f_j\in K$, $f_j\to f$ 于 $L^2$, 每个 Fourier 系数因其为有界泛函而收敛. 对每个有限 $N$, 取极限得 $\sum_{|n|\le N}(1+n^2)^s|\widehat f(n)|^2\le B^2$. 令 $N\to\infty$, 得 $f\in H^s$ 且在 $K$ 内. 所以 $K$ 闭、有界且具有一致有限维逼近, 定理\ref{thm:compact-approx}给出紧性. $s>0$ 正是尾界趋于零的条件.
\end{proof}
\begin{exercise}\label{ps10:4}
原题4. 设复矩阵 $(A_{ij})_{i,j\ge1}$ 满足 $\sum_i\sum_j|A_{ij}|^2<\infty$, 定义 $(Ta)_i=\sum_{j\ge1}A_{ij}a_j$ 于 $a\in\ell^2$. (a) 证明 $T\in\B(\ell^2)$ 且 $\norm T\le(\sum_{i,j}|A_{ij}|^2)^{1/2}$. (b) 定义 $(T_na)_i=\sum_{j=1}^nA_{ij}a_j$ 于 $i\le n$, 其余坐标零; 证明 $T_n$ 有限秩且 $\norm{T-T_n}\to0$, 从而 $T$ 紧.
\end{exercise}
\begin{proof}[AI 独立解答]
(a) 每一行平方可和, Cauchy--Schwarz 保证各坐标级数绝对收敛. 同一不等式给出 $\sum_i|(Ta)_i|^2\le\sum_{i,j}|A_{ij}|^2\norm a_2^2$, 故结果在 $\ell^2$, 线性有界及范数界成立.

(b) $T_n$ 的像包含于前 $n$ 个标准基的张成, 故有限秩. 对差矩阵用(a), 得 $\norm{T-T_n}^2\le\sum_{\max(i,j)>n}|A_{ij}|^2\to0$. 此尾和既截行也截列; 由定理\ref{thm:compactnorm}得紧. 这补全第\ref{ch:compact}章平方可和矩阵的计算.
\end{proof}
\begin{exercise}\label{ps10:5}
原题5. 在 $\ell^2$ 上令 $La=(a_2,a_3,\ldots)$, $Ra=(0,a_1,a_2,\ldots)$. (a) $L$ 或 $R$ 紧吗? 说明理由. (b) 证明 $R$ 无特征值. (c) 证明 $\sigma(R)=\{\lambda\in\C:|\lambda|\le1\}$. (d) 证明 $L$ 的特征值集为开单位圆盘, 并求各特征空间. (e) 证明 $\sigma(L)$ 为闭单位圆盘. 原提示(d): 解递推 $a_{k+1}=\lambda a_k$.
\end{exercise}
\begin{proof}[AI 独立解答]
(a) 二者都不紧. $Re_n=e_{n+1}$ 与 $Le_{n+1}=e_n$ 都是两两距离 $\sqrt2$ 的序列, 无收敛子列.

(b) 若 $Ra=\lambda a$, $\lambda=0$ 时 $R$ 单射使 $a=0$; $\lambda\ne0$ 时第一坐标给 $a_1=0$, 后续坐标递推全零, 无非零特征向量.

(c) $\norm R=1$, $|\lambda|>1$ 时 Neumann 级数使 $R-\lambda I$ 可逆. 若 $|\lambda|<1$, 取 $y=(1,\overline\lambda,\overline\lambda^2,\ldots)\in\ell^2\setminus\{0\}$. $R^*=L$, 所以 $(R-\lambda I)^*y=(L-\overline\lambda I)y=0$, 得 $\ran(R-\lambda I)\subset y^\perp\ne\ell^2$, 不可逆. 谱闭（命题\ref{prop:spectrumcompact}）, 所以单位圆边界也在谱中.

(d) 递推给出 $a=a_1(1,\lambda,\lambda^2,\ldots)$. 非零此列平方可和当且仅当 $|\lambda|<1$, 所以每个这样的特征空间正是该向量的一维张成, 其余 $\lambda$ 无特征向量.

(e) $\norm L=1$, 圆盘外由 Neumann 可逆; 圆盘内由(d)已有特征值, 边界由谱闭性收入. 不把“没有特征值”等同于“没有谱”.
\end{proof}
\begin{exercise}\label{ps10:6}
原题6（原注“不提交”）. 对 $f\in L^2([0,1])$, 定义 $Tf(x)=\int_0^xf(t)\dd t$. (a) 证明 $T$ 有界且紧; 原提示为 $|Tf(x)|\le\norm f_2$, $|Tf(x)-Tf(y)|\le\norm f_2|x-y|^{1/2}$. (b) 若 $Tf=0$, 证明每个 $0\le a<b\le1$ 都有 $\int_a^bf=0$, 推出 $\int_0^1f\chi=0$ 对每个简单函数 $\chi$ 成立, 从而 $f=0$. (c) 若 $\lambda\ne0$ 且 $(T-\lambda I)f=0$, 证明 $f$ 有 $C^1$ 代表元且满足 $f-\lambda f'=0$, $f(0)=0$, 从而 $f=0$. 因此 $T$ 没有特征值.
\end{exercise}
\begin{proof}[AI 独立解答]
(a) Cauchy--Schwarz 给出原提示的两个界, 特别地 $Tf$ 连续且 $\norm{Tf}_2\le\norm f_2$. 对有界输入列, 这些界给出一致有界与等度连续, Arzelà--Ascoli 使像有一致收敛子列, 因而有 $L^2$ 收敛子列, 算子紧.

(b) $Tf$ 连续, 若作为 $L^2$ 等价类为零, 则连续代表元处处零, 所以 $\int_a^bf=Tf(b)-Tf(a)=0$. 因而对每个区间阶梯函数 $\psi$ 有 $\int f\psi=0$. 对任一可测简单函数 $\chi$, 用题\ref{ps07:1}前置密度定理在 $L^2$ 中取阶梯 $\psi_n\to\chi$; Cauchy--Schwarz 保证 $\int f\psi_n\to\int f\chi=0$. 最后以阶梯逼近 $\overline f$ 得 $\int|f|^2=0$, 所以 $f=0$. 这里特意区分任意可测简单函数与区间阶梯函数, 不把二者误称相同.

(c) 等式给出 $f$ 的连续代表元 $g=Tf/\lambda$, 且 $g(0)=0$. 因 $f=g$ 几乎处处, $Tf=Tg$, 对连续 $g$ 用微积分基本定理得 $\lambda g'=g$, 所以 $g\in C^1$. 对 $\ee^{-x/\lambda}g(x)$ 求导得零, 初值又为零, 故 $g=0$, $f=0$ 于等价类.
\end{proof}

\originalchapter{期中测验}\label{ch:midterm}
原件3页, 题面2页, 5题共40分. 原开放书籍及禁止协作规定见课程信息; 以下是归档学习用的 AI 独立解答.
\begin{exercise}\label{midterm:1}
原题1（8分）. 对 $a<b$, 令 $\Lambda([a,b])$ 为所有 Lipschitz 连续函数, 即存在 $L\ge0$ 使 $|f(x)-f(y)|\le L|x-y|$. 证明它在范数 $\norm f=\norm f_\infty+\sup_{x\ne y}|f(x)-f(y)|/|x-y|$ 下为 Banach 空间. 原题允许使用上确界范数下 $C([a,b])$ 完备.
\end{exercise}
\begin{proof}[AI 独立解答]
记第二项为 $[f]_{\rm Lip}$. 和及数乘的 Lipschitz 常数受相应常数的和及绝对数乘控制, 所以这是向量空间. 上确界项保证正定性, 两项分别齐次且满足三角不等式, 因而确为范数.

设 $f_n$ 为本范数 Cauchy 列. 它在上确界范数下收敛于连续 $f$. 存在 $M$ 使 $[f_n]_{\rm Lip}\le M$, 对每对 $x,y$ 取极限得 $|f(x)-f(y)|\le M|x-y|$, 所以 $f\in\Lambda$. 又对 $n,m\ge N$, 有 $[f_n-f_m]_{\rm Lip}<\eps$. 固定 $n,x,y$, 令 $m\to\infty$, 再对 $x\ne y$ 取上确界得 $[f_n-f]_{\rm Lip}\le\eps$. 与一致收敛结合, 得本范数收敛.
\end{proof}
\begin{exercise}\label{midterm:2}
原题2（8分）. 对 $a\in\ell^\infty$, 证明 $\norm{a+c_0}_{\ell^\infty/c_0}=\limsup_{n\to\infty}|a_n|=\inf_{n\ge1}\sup_{k\ge n}|a_k|$. 原题允许使用: 加上趋零数列不改变绝对值的上极限.
\end{exercise}
\begin{proof}[AI 独立解答]
取 $b^{(N)}=(-a_1,\ldots,-a_N,0,\ldots)\in c_0$, 则商范数至多 $\sup_{k>N}|a_k|$, 令 $N\to\infty$ 得上界. 对任意 $b\in c_0$, $\limsup|a_n|=\limsup|a_n+b_n|\le\norm{a+b}_\infty$; 对 $b$ 取下确界得下界. 许可使用的上极限恒等式也由 $\big||a_n+b_n|-|a_n|\big|\le|b_n|\to0$ 直接得到.
\end{proof}
\begin{exercise}\label{midterm:3}
原题3. (a)（4分）若 $V,W$ 为 Banach 空间, $T_n\in\B(V,W)$, 线性 $T:V\to W$ 满足 $T_nx\to Tx$ 对每个 $x\in V$ 成立, 证明 $T$ 有界. 不假设算子范数收敛. (b)（4分）若同一向量空间的两范数满足 $\norm v_1\le\norm v_2$, 且两范数下都完备, 证明存在 $C>0$ 使 $\norm v_2\le C\norm v_1$.
\end{exercise}
\begin{proof}[AI 独立解答]
(a) 每条收敛的 $(T_nx)$ 有界. 统一有界原理\ref{thm:ub}给出 $M=\sup_n\norm{T_n}<\infty$, 故 $\norm{Tx}=\lim_n\norm{T_nx}\le M\norm x$. 此复用推论\ref{cor:stronglimit}.

(b) 恒等映射 $I:(V,\norm\cdot_2)\to(V,\norm\cdot_1)$ 为有界双射, 两端完备. 有界逆定理给出其逆有界, 即所需不等式; 零空间任选 $C>0$. 这也复用第\ref{ch:baire}章两完备范数的习题解答.
\end{proof}
\begin{exercise}\label{midterm:4}
原题4. 设 $E\subset\R$ 可测, $f,g:E\to\R$ 可测, 定义 $\norm f_\infty=\inf\{C\ge0:m(\{|f|>C\})=0\}$, 空集下确界取 $+\infty$. (a)（4分）证明 $|f|\le\norm f_\infty$ 几乎处处. (b)（4分）对 $c\in\R$, 证明 $\norm{cf}_\infty=|c|\norm f_\infty$, $\norm{f+g}_\infty\le\norm f_\infty+\norm g_\infty$.
\end{exercise}
\begin{proof}[AI 独立解答]
(a) 若本质上确界无穷, 不等式处处成立. 若 $M<\infty$, 对每个 $n$ 选可行本质界小于 $M+1/n$, 除去其零测例外的可数并, 得 $|f|\le M+1/n$ 对所有 $n$ 成立, 令 $n\to\infty$.

(b) $c\ne0$ 时可行界经乘 $|c|$ 一一对应, 包括没有有限可行界的情形. $c=0$ 时 $cf=0$ 的范数为零, 等式使用非负扩展实数约定 $0\cdot\infty=0$. 对三角不等式, 任一右端项为无穷时自动成立; 否则在两个本质界的共同零测例外外用逐点三角不等式, 再取本质界下确界.
\end{proof}
\begin{exercise}\label{midterm:5}
原题5. 设非负简单函数 $\phi=\sum_{i=1}^n a_i\ind_{A_i}$, $A_i\subset E$ 为可测分割. (a)（4分）若 $F\subset E$ 可测, 证明 $\int_F\phi=\int_E\phi\ind_F$. (b)（4分）若 $F,G\subset E$ 可测且不交, 证明 $\int_{F\cup G}\phi=\int_F\phi+\int_G\phi$.
\end{exercise}
\begin{proof}[AI 独立解答]
(a) $\ind_{A_i\cap F}=\ind_{A_i}\ind_F$, 所以两积分都等于 $\sum_i a_i m(A_i\cap F)$. (b) $A_i\cap(F\cup G)$ 为 $A_i\cap F$ 与 $A_i\cap G$ 的不交并; 对测度用可加性, 再作有限个非负求和, 即得等式. 所有加法允许无穷, 零系数乘无穷按定义取零.
\end{proof}

\originalchapter{期末作业}\label{ch:final}
原件3页, 题面2页, 5题共40分. 本章按官方资源页及大纲命名为期末作业; 原PDF内部沿用 exam 用语不改变这一归档分类.
\begin{exercise}\label{final:1}
原题1（8分）. 令 $f(x)=-1$ 于 $[-\pi,0]$, $f(x)=1$ 于 $(0,\pi]$. 求 $\widehat f(n)=(2\pi)^{-1}\int_{-\pi}^{\pi}f(x)\ee^{-\ii nx}\dd x$, $n\in\Z$, 并用结果计算 $\sum_{k=0}^\infty(2k+1)^{-2}$. 原提示: 用范数与 Fourier 系数的关系.
\end{exercise}
\begin{proof}[AI 独立解答]
零点取值不影响积分, $\widehat f(0)=0$. 对 $n\ne0$, 分别积分两段, 或利用奇性, 得 $\widehat f(n)=-\ii[1-(-1)^n]/(\pi n)$. 因而非零偶频系数为零, 奇频为 $-2\ii/(\pi n)$. Parseval 给出 $1=(2\pi)^{-1}\norm f_2^2=\sum_{n\in\Z}|\widehat f(n)|^2=(8/\pi^2)\sum_{k\ge0}(2k+1)^{-2}$, 所以所求和为 $\pi^2/8$. 此使用第\ref{ch:fourier}章已证明的完备性, 不是未经证明的点态级数求和.
\end{proof}
\begin{exercise}\label{final:2}
原题2（8分）. 用适当收敛定理计算
\[
\lim_{n\to\infty}\int_0^1\frac{n\cos x}{1+n^2x^2}\dd x,
\qquad
\lim_{n\to\infty}\int_0^\infty\frac{n\sin(x/n)}{x(1+x^2)}\dd x.
\]
原提示: 一个或两个积分中可先换元.
\end{exercise}
\begin{proof}[AI 独立解答]
第一个积分以 $t=nx$ 换元, 等于 $\int_0^\infty\ind_{[0,n]}(t)\cos(t/n)/(1+t^2)\dd t$. 被积函数对固定 $t$ 趋于 $(1+t^2)^{-1}$, 绝对值受此全轴可积函数控制. 控制收敛给出极限 $\int_0^\infty(1+t^2)^{-1}\dd t=\pi/2$. 原变量下虽几乎处处趋零, 却没有相应统一控制, 不能直接把极限带入.

第二个积分对每个 $x>0$ 有 $n\sin(x/n)/x\to1$, 且 $|n\sin(x/n)|\le x$, 所以受 $(1+x^2)^{-1}$ 控制. 同一定理给出极限 $\pi/2$. 单点 $x=0$ 可按连续延拓赋值, 不影响积分; 两个反正切积分的 Lebesgue 值由连续函数的截断积分及单调收敛得到.
\end{proof}
\begin{exercise}\label{final:3}
原题3（8分）. 若 $W$ 为 Hilbert 空间 $H$ 的线性子空间, 证明 $W^\perp=\{u:\ip uw=0\text{对所有 }w\in W\}$ 是闭子空间, 且 $(W^\perp)^\perp=\overline W$. 原PDF右端带闭包横线.
\end{exercise}
\begin{proof}[AI 独立解答]
此与题\ref{ps08:2}相同, 复用其完整证明: $W^\perp$ 是连续内积泛函闭核的交; 对 $\overline W$ 作正交分解 $x=v+z$, 若 $x\in(W^\perp)^\perp$, 则 $0=\ip xz=\norm z^2$, 故 $x\in\overline W$. 反向包含由内积连续性得到. 不去掉题面的闭包.
\end{proof}
\begin{exercise}\label{final:4}
原题4. 对有界复数列 $(\mu_k)$, 定义 $Ma=(\mu_ka_k)_k$ 于 $\ell^2$, 已知 $M$ 有界且 $\norm M\le\sup_k|\mu_k|$. (a)（4分）证明 $M$ 自伴当且仅当每个 $\mu_k$ 为实数. (b)（4分）若 $\mu_k\to0$, 证明 $M$ 紧. 原题允许不加证明引用 Piazza 中讨论的性质: 对 $A\in\B(\ell^2)$,
\[
\overline{\{Ab:\norm b_2\le1\}}=\{Ab:\norm b_2\le1\}.
\]
以下证明不需要该额外提示.
\end{exercise}
\begin{proof}[AI 独立解答]
(a) 绝对收敛的内积计算给出 $M^*b=(\overline{\mu_k}b_k)_k$. 故实数列使 $M=M^*$. 反之测试 $e_k$ 得 $\mu_k=\overline{\mu_k}$, 每项实.

(b) 令 $M_N$ 只保留前 $N$ 个对角项, 它有限秩, 且 $\norm{M-M_N}=\sup_{k>N}|\mu_k|\to0$; 等式下界由测试 $e_k$ 后取上确界得到. 由定理\ref{thm:compactnorm}得紧. 此复用第\ref{ch:compact}章对角紧算子的习题, 不要求 $(\mu_k)$ 平方可和.
\end{proof}
\begin{exercise}\label{final:5}
原题5（8分）. 在 $L^2([0,1])$ 上令 $Mf(x)=xf(x)$, 已知 $M$ 有界且 $\norm M\le1$. 证明 $\sigma(M)=[0,1]$, 且 $M$ 无特征值. 原提示: 对 $\lambda\in[0,1]$ 考察常数函数1是否在 $\ran(M-\lambda I)$ 内.
\end{exercise}
\begin{proof}[AI 独立解答]
若 $\lambda\notin[0,1]$, 距离 $d=\dist(\lambda,[0,1])>0$, 乘法 $f\mapsto f/(x-\lambda)$ 为有界双向逆, 范数至多 $1/d$, 故 $\lambda$ 在预解集. 若 $\lambda\in[0,1]$, 方程 $(x-\lambda)f(x)=1$ 将要求 $f(x)=1/(x-\lambda)$ 几乎处处. 此函数在 $\lambda$ 的单侧或双侧邻域平方不可积, 包括端点0、1, 故常数1不在值域, $M-\lambda I$ 不满射. 因而谱正是 $[0,1]$.

若 $(M-\lambda I)f=0$, 在 $x\ne\lambda$ 处 $f=0$ 几乎处处; 可能的单点支集是零测集, 所以 $f=0$ 于 $L^2$ 等价类. 对非实或区间外 $\lambda$ 更是处处没有零因子, 同样无非零特征向量. 这补全第\ref{ch:spectrum}章乘法谱例子的区间端点核对.
\end{proof}

\originalchapter{来源与改编}\label{ch:source}
\originalsection{逐讲对应}
本书主源是 Casey Rodriguez 授课、Andrew Lin 记录的 MIT 18.102 Spring 2021. 合订本125个 PDF 物理页, 23份独立讲义合计147个物理页; 差别来自分讲排版及前后信息, 并不是另外147页独立课程内容. 官方目录重复链接第1讲, 已按实际下载文件去重. 下表中页码为125页合订本的物理页序; 讲次交界页可能同时包含前讲末尾和后讲日期, 因此边界按下一讲日期所在页划分, 只用作定位.

\begin{longtable}{p{10mm}p{20mm}p{125mm}}
\toprule
讲次 & 合订页序 & 内容及本书对应\\
\midrule
\endhead
1 & 1--6 & 向量空间、范数、连续有界函数和 Banach 空间: 第\ref{ch:norm}章.\\
2 & 7--13 & 绝对级数、有界算子、算子空间和对偶: 第\ref{ch:norm}、\ref{ch:operators}、\ref{ch:dual}章.\\
3 & 14--17 & 子空间、商空间、Baire 与统一有界: 第\ref{ch:operators}、\ref{ch:baire}章.\\
4 & 18--21 & 开映射、闭图及 Zorn 引理: 第\ref{ch:baire}、\ref{ch:dual}章.\\
5 & 22--25 & Hamel 基与 Hahn--Banach: 第\ref{ch:dual}章, 补全复数保范数证明.\\
6 & 26--30 & 双对偶及外测度: 双对偶在第\ref{ch:dual}章完整展开; 外测度在第\ref{ch:measure}章给定义接口.\\
7 & 31--35 & $\sigma$ 代数、可测集和 Borel 集: 第\ref{ch:measure}章定义接口; 构造证明交由实变科.\\
8 & 36--40 & Lebesgue 测度的构造与性质: 第\ref{ch:measure}章明示调用, 未逐项重编构造.\\
9 & 41--46 & 可测函数及其运算: 第\ref{ch:measure}章必要定义与闭包性质.\\
10 & 47--52 & 简单函数逼近及非负积分: 第\ref{ch:measure}章.\\
11 & 53--58 & 单调收敛、Fatou: 第\ref{ch:measure}章完整证明.\\
12 & 59--64 & 可积函数、控制收敛和 Riemann 积分关系: 第\ref{ch:measure}章完整证明控制收敛及连续函数的两种积分一致性.\\
13 & 65--70 & $L^p$、Hölder、Minkowski、完备性、内积: 第\ref{ch:measure}、\ref{ch:hilbert}章完整证明.\\
14 & 71--76 & Hilbert 空间、正交系、Gram--Schmidt、Bessel: 第\ref{ch:hilbert}章.\\
15 & 77--82 & 正交基、Parseval、可分分类、Fourier 与 Dirichlet 核: 第\ref{ch:hilbert}、\ref{ch:fourier}章.\\
16 & 83--88 & Fejér 核与 Fourier 完整性: 第\ref{ch:fourier}章, 图形按定义重绘.\\
17 & 89--93 & 闭凸集极小化、正交分解、Riesz 表示: 第\ref{ch:adjoint}章.\\
18 & 94--99 & 伴随、紧集与一致小尾项: 第\ref{ch:adjoint}、\ref{ch:compact}章.\\
19 & 100--104 & 紧集判别、有限秩算子: 第\ref{ch:compact}章.\\
20 & 105--109 & 紧算子、谱、Neumann 级数、自伴范数: 第\ref{ch:compact}、\ref{ch:spectrum}章.\\
21 & 110--113 & 自伴谱端点、紧自伴特征空间: 第\ref{ch:spectrum}章, 补全端点向量范数估计.\\
22 & 114--118 & Fredholm 条件、最大原则、紧自伴谱定理: 第\ref{ch:spectrum}章.\\
23 & 119--125 & Green 核、正弦谱、平方根、Dirichlet 问题: 第\ref{ch:dirichlet}章, 校正核符号并补全正则性.\\
\bottomrule
\end{longtable}

\originalsection{补充、勘误与边界}
编者补充包括赋范空间完备化及稠密延拓, 数列空间对偶的计算证明, 平行四边形恒等式的逆向证明, 一般弱与弱星拓扑、可分对偶球弱星紧性, 及配套例题、习题的完整解答. 第\ref{ch:weak}章的 Hilbert 弱收敛与第\ref{ch:compact}章的紧算子关系另参 Richard Melrose 2020讲义第3章第15节, 印刷页82--84. 这些补充均围绕本册的依赖关系, 不把辅助讲义的所有章节纳入主源范围.

原第4讲的一处说明把标准坐标向量称作 $\ell^1$ 的 Hamel 基; 它们的有限张成实际是 $c_{00}$, 稠密却不是全空间. 原第5讲对复数 Hahn--Banach 的实虚部相容性略写, 本册改用实化与单位相位估计. 原第13讲某个衰减条件的不等号写成下界, 其证明使用的是上界; 本册的积分例子以直接积分给出准确条件. 原第16讲 Fejér 核质量陈述的积分区间误写为整条实轴, 而周期核在实轴上的积分不是1; 本册使用一个周期的积分. 该讲的半角公式也按正确的 $1-\cos t=2\sin^2(t/2)$ 使用.

原第21讲从二次型趋于零到 $\norm{(A-\lambda I)x_n}\to0$ 的推理需要额外估计, 已在自伴谱端点证明中补全. 原第22讲的谱基先在非闭值域上讨论, 本册明确把完整 Hilbert 展开放在值域闭包 $(\ker A)^\perp$ 中. 原第23讲的分段 Green 核符号与 $-u''=f$ 不一致, 已以正核 $\min(x,y)(1-\max(x,y))$ 修正; 非零特征值为 $(k\pi)^{-2}$, 平方根的上确界估计依赖 $\norm{f}_2$, 且乘法算子只需有界, 不应被默认称为紧算子. 本书通过已证明的正弦基、统一收敛及 Green 表示补齐这些论证.

本册完整覆盖主源的 Banach 与 Hilbert 空间主线、Fourier 完备性、紧算子、自伴谱及连续非负势的区间 Dirichlet 问题, 并收入官方全部10份作业、期中测验和期末作业的题面及完整解答. 原第6--12讲全部测度构造和非可测集细节仍未逐页重编, 但考核实际要求的外测度平移与缩放、开集区间分解、有限区间并逼近、Littlewood 三原则、$L^\infty$ 完备性与 $L^p$ 可分性均在对应题解内证明.

未覆盖视频文字稿, Carleson 定理的证明, 一般不可分 Banach--Alaoglu 定理, 一般 Banach 空间自反性的等价刻画, 非自伴 Fredholm 全理论, 一般非紧自伴谱测度、无界算子的域理论、一般 Sobolev 与分布、迹类和 Schatten 理论. 原作业指定的周期 Fourier 加权 $H^s$ 空间、$s>1/2$ 连续嵌入、$s>0$ 紧嵌入及磨光逼近已收入, 不能再把这几个具体结果列作未覆盖, 也不能由此宣称覆盖一般 Sobolev 理论. 更广专题有些出现在176页 Melrose 辅助讲义中, 不由参考书目录推断本册完成范围.

\originalsection{作业与考核对应}
下表中的页数为各原PDF的物理页数, 包含最后的 MIT OCW 署名页. “单元”指显式字母子问, 或没有字母分问的原编号单题. 期末第2题的两个积分按一个原编号单元登记, 解答中逐个计算. 完整原题条件、子问及中文定位另见 assessment-inventory.json; 原件没有官方答案, 本书的新增解答统一标为 AI 独立解答.
\begin{longtable}{p{32mm}r r r p{63mm}}
\toprule
材料 & 原页数 & 题数 & 单元 & 本书定位\\
\midrule\endhead
作业1 & 3 & 5 & 9 & 第\ref{sec:ps01}节, Hölder、完备、数列对偶.\\
作业2 & 3 & 5 & 10 & 第\ref{sec:ps02}节, 可逆性、商空间、闭值域和反例.\\
作业3 & 3 & 4 & 9 & 第\ref{sec:ps03}节, 对偶转置、外测度、开集.\\
作业4 & 2 & 3 & 4 & 第\ref{sec:ps04}节, 可测性与集合逼近.\\
作业5 & 3 & 4 & 9 & 第\ref{sec:ps05}节, 扩展实值运算与 Egorov.\\
作业6 & 3 & 4 & 7 & 第\ref{sec:ps06}节, 连续逼近、Riemann--Lebesgue.\\
作业7 & 2 & 4 & 9 & 第\ref{sec:ps07}节, 可分性、$L^\infty$、极化.\\
作业8 & 3 & 4 & 9 & 第\ref{sec:ps08}节, 正交、$H^s$ 与连续嵌入.\\
作业9 & 4 & 5 & 10 & 第\ref{sec:ps09}节, 磨光、凸集投影.\\
作业10 & 4 & 6 & 15 & 第\ref{sec:ps10}节, 紧性、移位谱、Volterra.\\*
期中测验 & 3 & 5 & 8 & 第\ref{ch:midterm}章.\\*
期末作业 & 3 & 5 & 6 & 第\ref{ch:final}章.\\
\bottomrule
\end{longtable}
12份合计36个原PDF物理页、54道编号题、105个解答单元, 即87个字母子问与18道未分问单题.

原12章的30份例题与习题解答全部保留, 仅增交叉引用标签. 直接对应的复用包括: 题\ref{ps03:1}复用例题\ref{legacy:ch03:example1}的延拓泛函; 题\ref{ps07:1}(a)复用习题\ref{legacy:ch05:exercise1}的有限测度包含; 题\ref{midterm:3}(b)复用习题\ref{legacy:ch04:exercise1}的完备范数比较; 题\ref{ps09:3}复用习题\ref{legacy:ch07:exercise1}的自伴幂等投影; 题\ref{final:4}(b)复用习题\ref{legacy:ch10:exercise1}的对角紧性. 题\ref{ps08:1}(b)把习题\ref{legacy:ch06:exercise1}的有限正交近似推广到闭包, 题\ref{final:5}细化例题\ref{legacy:ch11:example1}的乘法谱; 都保留所需扩展论证, 不以一个章节号代替新增目标的证明. 其余题目对既有定理的精确引用均在解答内列明.

源题勘误还包括作业2第4题同构的箭头, 作业5递减测度连续性的名称, 作业8第1题漏印范数, 作业9前置子列定理的字母与阈值混写. 作业9第4题缺少非空假设使原命题按字面为假, 本书明确给出空集反例再补条件. 作业8第2题与期末第3题的闭包横线按原PDF保留. 所有这些修订均在题面附近说明.

原件目录保留实际下载的 PDF、TeX 与原源码 ZIP, 每个文件的字节、SHA256、页数和许可依据均可在源清单定位. 本书中文源码 ZIP 是独立可编译的干净源码, 不包含字体、编译缓存、页面图像或原 ZIP 内的 macOS 元数据. 原源归档与中文成册的范围在两个目录的说明中分别列明.

\begin{thebibliography}{9}
\bibitem{lin} Casey Rodriguez (授课), Andrew Lin (笔记). \textit{18.102: Introduction to Functional Analysis}. Spring 2021. MIT OpenCourseWare. \href{https://ocw.mit.edu/courses/18-102-introduction-to-functional-analysis-spring-2021/resources/complete-lecture-notes/}{125页完整讲义}及23讲源文件.
\bibitem{melrose} Richard Melrose. \textit{Functional Analysis: Lecture Notes for 18.102, Spring 2020}. Version 0.9A, run May 16, 2020. MIT. \href{https://ocw.mit.edu/courses/18-102-introduction-to-functional-analysis-spring-2021/resources/mit18_102s20_lec_fa/}{OCW提供的176页辅助讲义}. 本册仅用作关键条件和弱收敛的交叉核对.
\bibitem{assessments} Casey Rodriguez 课程. \textit{18.102/18.1021 Assignments and Exams}. Spring 2021. MIT OpenCourseWare. \href{https://ocw.mit.edu/courses/18-102-introduction-to-functional-analysis-spring-2021/pages/assignments-and-exams/}{10份作业、期中测验与期末作业}, 12份原PDF共36页; 原件没有提供官方答案.
\bibitem{license} MIT OpenCourseWare. \href{https://ocw.mit.edu/pages/privacy-and-terms-of-use/}{使用条款}; \href{https://creativecommons.org/licenses/by-nc-sa/4.0/}{CC BY-NC-SA 4.0 International}.
\end{thebibliography}
\end{document}
